% concatenated from versione_finale.tex, proof_flin_G.tex
\documentclass{article}


% if you need to pass options to natbib, use, e.g.:
%     \PassOptionsToPackage{numbers, compress}{natbib}
% before loading neurips_2025


% ready for submission
%\usepackage{neurips_2025}


% to compile a preprint version, e.g., for submission to arXiv, add add the
% [preprint] option:
\usepackage[final]{neurips_2025}


% to compile a camera-ready version, add the [final] option, e.g.:
%     \usepackage[final]{neurips_2025}


% to avoid loading the natbib package, add option nonatbib:
%    \usepackage[nonatbib]{neurips_2025}


\usepackage[utf8]{inputenc} % allow utf-8 input
\usepackage[T1]{fontenc}    % use 8-bit T1 fonts
\usepackage[hidelinks]{hyperref}       % hyperlinks
\usepackage{url}            % simple URL typesetting
\usepackage{booktabs}       % professional-quality tables
\usepackage{amsfonts}       % blackboard math symbols
\usepackage{nicefrac}       % compact symbols for 1/2, etc.
\usepackage{microtype}      % microtypography
\usepackage{xcolor}         % colors

% EM: PACKAGES:
\usepackage{amssymb, amsmath, amscd, amsthm, enumitem, mathtools, xcolor, comment,bbm, tikz-cd}
%\usepackage{nicematrix}
\usepackage{graphicx}
%\graphicspath{ {./neurips/} }

% EM: THM STYLES:
\newtheorem{thm}{Theorem}[section]
\newtheorem{cor}[thm]{Corollary}
\newtheorem{lem}[thm]{Lemma}
\newtheorem{prop}[thm]{Proposition}
\theoremstyle{definition}
\newtheorem{defn}[thm]{Definition}
\newtheorem{ass}{Assumption}
\theoremstyle{remark}
\newtheorem{rem}[thm]{Remark}
\newtheorem{claim}[thm]{Claim}
\newtheorem{ex}[thm]{Example}
\numberwithin{equation}{section}

% EM: SOME MACROS:
%Alcune macros
\DeclareMathOperator*{\argmax}{arg\,max}
\DeclareMathOperator*{\argmin}{arg\,min}
\DeclareMathOperator{\erf}{erf}
\newcommand{\D}{\mathcal{D}}
\newcommand{\X}{\mathcal{X}}
\newcommand{\Y}{\mathcal{Y}}
\newcommand{\R}{\mathbb{R}}
\newcommand{\N}{\mathbb{N}}
\newcommand{\Z}{\mathbb{Z}}
\newcommand{\K}{\mathbb{K}}
\newcommand{\E}{\mathbb{E}}
\newcommand{\p}{\varphi}
\newcommand{\id}{\mathrm{id}}
\newcommand{\s}{\psi}
\newcommand{\eps}{{\varepsilon}}
\newcommand{\hilb}{\mathcal{H}}
\newcommand{\test}{\textup{TE}}
\newcommand{\train}{\textup{TR}}
\newcommand{\kttest}{k_{xx}}
\newcommand{\ktest}{k_{x\X}}
\newcommand{\ktrain}{k_{\X\X}}
\newcommand{\lin}{{\textup{lin}}}
\newcommand{\lammax}{\lambda_{\textup{max}}}
\newcommand{\lammin}{\lambda_{\textup{min}}}
\newcommand{\sigmin}{\sigma_{\textup{min}}}
\newcommand{\sigmax}{\sigma_{\textup{max}}}

\newcommand{\norm}[2]{
    \|#1\|_{#2}
}
\newcommand{\partv}[1]{
    \frac{\partial}{\partial #1}
}
\newcommand{\wass}[3]{
    \mathcal W_{#1} \left( #2, #3 \right)
}

\title{Quantitative convergence of trained single layer neural networks to Gaussian processes}


% The \author macro works with any number of authors. There are two commands
% used to separate the names and addresses of multiple authors: \And and \AND.
%
% Using \And between authors leaves it to LaTeX to determine where to break the
% lines. Using \AND forces a line break at that point. So, if LaTeX puts 3 of 4
% authors names on the first line, and the last on the second line, try using
% \AND instead of \And before the third author name.


\author{%
Eloy Mosig García\\
  Department of Mathematics\\
  University of Pisa\\
  Largo Bruno Pontecorvo, 5, 56127 Pisa PI, Italia\\
  \texttt{eloy.mosig@phd.unipi.it}\\
  \And
Andrea Agazzi\\
Department of Mathematics and Statistics\\
University of Bern\\
Alpeneggstrasse 22, 3012 Bern\\
\texttt{andrea.agazzi@unibe.ch}\\
  %David S.~Hippocampus\thanks{Use footnote for providing further information
  %  about author (webpage, alternative address)---\emph{not} for acknowledging
  %  funding agencies.} \\
  %Department of Computer Science\\
  %Cranberry-Lemon University\\
  %Pittsburgh, PA 15213 \\
  %\texttt{hippo@cs.cranberry-lemon.edu} \\
  % examples of more authors
  \And
  Dario Trevisan\\
  Department of Mathematics\\
  University of Pisa\\
  Largo Bruno Pontecorvo, 5, 56127 Pisa PI, Italia\\
  \texttt{dario.trevisan@unipi.it}\\
  % \AND
  % Coauthor \\
  % Affiliation \\
  % Address \\
  % \texttt{email} \\
  % \And
  % Coauthor \\
  % Affiliation \\
  % Address \\
  % \texttt{email} \\
  % \And
  % Coauthor \\
  % Affiliation \\
  % Address \\
  % \texttt{email} \\
}

\begin{document}

\maketitle

\begin{abstract}

    In this paper, we study the quantitative convergence of shallow neural networks trained via gradient descent to their associated Gaussian processes in the infinite-width limit. 
    While previous work has established qualitative convergence under broad settings, precise, finite-width estimates remain limited, particularly during training. 
    We provide explicit upper bounds on the quadratic Wasserstein distance between the network output and its Gaussian approximation at any training time $t \ge 0$, demonstrating polynomial decay with network width. 
    Our results quantify how architectural parameters, such as width and input dimension, influence convergence, and how training dynamics affect the approximation error.

\begin{comment}
    The abstract paragraph should be indented \nicefrac{1}{2}~inch (3~picas) on
  both the left- and right-hand margins. Use 10~point type, with a vertical
  spacing (leading) of 11~points.  The word \textbf{Abstract} must be centered,
  bold, and in point size 12. Two line spaces precede the abstract. The abstract
  must be limited to one paragraph.
\end{comment}
\end{abstract}

\section{Introduction}
\label{sec:introduction}
\begin{comment}
    
In recent years, deep neural networks have attracted significant attention due to their exceptional performance across various tasks \cite{intro1, intro2}. 
This has sparked a growing interest in understanding the structure and behavior of these overparametrized models, and an auspicious research direction, dating back to \cite{neal96} in the shallow case, is the infinite-width limit, i.e., when the size of the hidden layers diverges. 

From a theoretical standpoint, it has been demonstrated that the wide limit of a neural network with Gaussian initialization exhibits characteristics of a Gaussian process, not only during the initialization phase \cite{Matthews2018GaussianPB} but also at any point during training via gradient descent and employing the mean square error loss function \cite{Lee2020}.  
Indeed, a significant breakthrough in understanding the dynamics of gradient descent in the wide limit was achieved through the investigation by Jacot et al. \cite{jacot} and Lee et al. \cite{Lee2020}, where the Neural Tangent Kernel (NTK) was recognized to play a key role in the training behavior of the network, under wide limit conditions. 
These theoretical insights form the foundation of our investigation. However, our study is also driven by practical implications: previous research has shown that wider networks tend to possess superior generalization capabilities \cite{overparametrization1, overparametrization2}, making them desirable in some real-world applications \cite{intro1, intro2}. 

Quantifying such convergence results in deep neural networks in the wide limit appears then of paramount importance for several reasons. 
Quantitative bounds that take into account the network structure allow us to assess the reliability and robustness of trained models, gaining from the well-established theory of Gaussian processes \cite{RasmussenBook}. 
Furthermore, by quantifying convergence, one can better understand the limitations and trade-offs associated with different network architectures, initialization schemes, and training algorithms.

% Deep neural networks have become a central topic in recent years, owing to their unparalleled performance across a diverse array of tasks \cite{intro1,intro2}.
% Our approach is to study the stochastic behavior of the network when the width of the hidden layers diverges.
% This approach is motivated by both theoretical and practical implications.
% On one hand, the so-called wide limit of the neural network with Gaussian initialization has been shown to behave as a Gaussian process both at initialization \cite{trevisan2023wide} and at arbitrary positive times during training when using square loss \cite{neal96, Lee2020}.
% For the choice of the square loss, the training dynamics admit a closed-form solution \cite{Lee2020}.
% The results we introduce in this paper are also desirable from the applicational point of view.
% Indeed, it has been shown that wider networks have better generalizing properties \cite{overparametrization1,overparametrization2}.

% A significant breakthrough was made with the study of gradient descent dynamics in the wide limit of the network \cite{Lee2020, jacot}.
% In particular the Neural Tangent Kernel (NTK) operator introduced by \cite{jacot} that determines the behavior of the network in the wide limit will be the main object of our study.

\end{comment}

Deep neural networks have achieved remarkable success across a wide range of tasks in computer vision, 
natural language processing, and scientific computing, often surpassing traditional models by large margins \cite{lecun2015deep}, \cite{goodfellow2016deep}. 
This empirical progress has sparked substantial interest in understanding the theoretical principles underlying their behavior, 
particularly in the overparameterized regime, where the number of parameters is larger than the one of training samples.

A major avenue of theoretical investigation in this sense focuses on studying the properties of neural networks in the infinite-width limit. 
For instance, when the network's parameters at initialization are sampled from a Gaussian distribution, it was shown \cite{neal96}, \cite{Matthews2018GaussianPB} 
that the network's output converges to a Gaussian process as width tends to infinity, 
providing a tractable framework for theoretical analysis.

This perspective was significantly extended by the introduction of the Neural Tangent Kernel (NTK) framework \cite{jacot}, 
allowing to characterize the training dynamics of infinitely wide neural networks under gradient descent in function space. 
In this limit, the network evolves approximately linearly around its initialization, 
and training can be understood as kernel regression with a fixed 
%\textcolor{lightgray}{data-dependent kernel.}
kernel which depends on the the architecture only, and is evaluated on input data.
This linearization dramatically simplifies the analysis of generalization and convergence, 
and has led to a large body of theoretical work on the expressivity and limitations of infinitely wide networks.

However, the practical relevance of NTK-based analyses hinges on the accuracy of their approximation at finite width. 
While existing literature has established qualitative convergence of wide neural networks in the NTK regime to Gaussian processes at positive training time \cite{Lee2020}, 
rigorous quantitative results — providing explicit finite-width error bounds — remain scarce. 
This gap limits the applicability of NTK theory to realistic settings where network width is large but finite.

Indeed, quantitative convergence guarantees are crucial to bridge theory and practice. 
They allow one to bound the discrepancy between the predictions of a finite-width network and its infinite-width NTK counterpart, 
thus enabling quantitative uncertainty quantification estimates and the safe deployment of theoretical insights to real-world architectures. 
Moreover, such estimates reveal how network width, depth, initialization, and training hyperparameters impact the validity of linear approximations during training. 
These insights are essential for developing principled training strategies and for diagnosing when the NTK regime offers a reliable approximation, or when nonlinear effects beyond NTK must be taken into account.

\subsection{Our contributions}

This paper provides rigorous quantitative estimates for the convergence of trained shallow neural networks towards their Gaussian process counterparts, measured in terms of quadratic Wasserstein distances.
Specifically, we extend previously established convergence bounds at initialization obtained by \cite{BasteriTrevisan2023}, \cite{Favaro2025} and \cite{trevisan2023wide} to arbitrary positive training times. 
Our results deliver explicit convergence rates that decay polynomially with network width, clearly delineating how the approximation error evolves during training.

Concretely, we demonstrate that the distance of the distribution of a shallow neural network's output trained via gradient descent 
to its Gaussian process approximation at any training time $t > 0$ satisfies explicit quantitative bounds dependent on network width. 
Our main theorem (\ref{thm:main}) shows that for any test point $x$ , under mild assumptions on the hidden layer width $n_1$ and on the regularity of the activation function we have:

\[
    \mathcal W_2^2(f_t(x),G_t(x)) = \mathcal O\left(\frac{\log n_1}{n_1}\right).
\]

We also address long-term training dynamics explicitly, characterizing convergence rates as training time diverges. 
Indeed, the above result continues to hold on timescales $t$ growing polynomially in the network width $n_1$ , 
as discussed in Remark \ref{rem:time}.

These results significantly refine prior qualitative statements, 
providing actionable quantitative guidelines on how network parameters and training duration determine the extent to which finite networks emulate their infinite-width limits.

\subsection{Related work}

The convergence of randomly initialized neural networks to Gaussian processes in the infinite-width limit has been a foundational result in the theory of neural networks. 
This phenomenon was first suggested by \cite{neal96} and later formalized for deep architectures by \cite{Matthews2018GaussianPB}.
The key insight that this correspondence extends beyond initialization was introduced by \cite{jacot}, 
who demonstrated that training dynamics in the infinite-width limit are governed by the so-called Neural Tangent Kernel (NTK), 
a deterministic kernel that linearizes the training trajectory. 
This sparked significant interest in the use of kernel methods to analyze deep learning models.

Following these developments, several works studied the convergence of finite-width neural networks to their limiting Gaussian processes. 
In particular, \cite{Lee2020} established that gradient descent dynamics in the NTK regime converge to those of a linearized model. 
More recently, the works of \cite{BasteriTrevisan2023} and \cite{trevisan2023wide} provided quantitative convergence rates at initialization, 
measured in Wasserstein distance, which laid the groundwork for a more refined understanding of the finite-width behavior of neural networks. 
Moreover, in the also recent work by \cite{Favaro2025}, additional quantitative results were obtained for total variation and convex distances.
Complementary to these works, \cite{bordino2025} used second-order Poincaré inequalities to derive QCLTs for Gaussian neural networks, 
obtaining a general but suboptimal convergence rate compared to the optimal $n^{-1}$.

However, these results were largely confined to the initialization regime. 
To this day, extensions to the full training trajectory remained limited, with few works addressing how approximation errors evolve over time or depend on architectural features such as width and depth. 
The present work builds on this gap by extending the quantitative convergence discussed above to trained networks, providing explicit bounds on the Wasserstein distance between the network output and the associated Gaussian process for any positive training time.

From a spectral perspective, the NTK's conditioning plays a central role in understanding convergence rates and generalization. 
Lower bounds on the smallest eigenvalue of the empirical NTK have been derived under various conditions. 
For instance, \cite{montufar} and \cite{bombari2022memorization} provide sharp bounds in the context of ReLU and smooth activation functions, respectively. 

Additionaly, \cite{carvalho2024} showed that under very mild assumptions on the non-linearity and non-proportionality of the training data, the analytic NTK is not degenerate.
These results are essential for establishing the stability of the gradient flow and, hence, for deriving quantitative convergence guarantees.

%\color{lightgray}

%A different line of work has focused on linear approximations of neural networks during training. 
%\cite{Bartlett_Montanari_Rakhlin_2021} and \cite{bachchizat} rigorously established that, under suitable conditions, networks trained with gradient descent remain close to their linearizations throughout training, provided the NTK remains sufficiently well-conditioned. 
%Our results exploit and extend these insights, combining them with probabilistic estimates to bound deviations from the Gaussian limit in a finite-width regime.

Our results are closely related to the work of \cite{Matthews2018GaussianPB}, who proved weak convergence of fully-connected BNNs at initialization to a Gaussian process under the metric 
$\rho(f,f') = \sum_{i \in \mathbb{N}} 2^{-i} \min\{1, |f(x_i)-f'(x_i)|\}$, defined on a countable input set. 
In our setting, the input set is finite; considering the restriction $\rho_F$, it follows that convergence in $\mathcal{W}_2$ implies convergence in $\rho_F$. 
%we show 
%$\rho_F(f,f') < \| f - f'\|_2$, hence $\mathbb{E}[\rho_F(f,f')] < \mathcal{W}_2(\mu,\nu)$. 
%Thus, convergence in $\mathcal{W}_2$ implies convergence in $\rho_F$. 
\cite{Matthews2018GaussianPB} also analyzed convergence under the maximum mean discrepancy (MMD). 
While MMD is not generally controlled by Wasserstein distances, connections have been established via regularized OT divergences (\cite{Feydy19,nietert2021smooth}).
Moreover, \cite{vayer} identified conditions on the RKHS kernel $k_H$ under which $MMD \lesssim \mathcal{W}_2$. 
Consequently, our bounds also imply MMD convergence under these conditions. 
The metric $\rho_F$ which offers a notion of pointwise convergence and is oblivious of the tails of the distributions, which helps stablish the results in \cite{Matthews2018GaussianPB}.  
On the other hand, $\mathcal{W}_2$ captures the geometric structure and scaling of the output space. 
Finally, while \cite{Matthews2018GaussianPB} address the more general setting deep networks, our analysis focuses on the shallow case, yielding new quantitative rates which improve previous ones in our setting.


%A foundational line of work has rigorously established that, under suitable overparameterization, neural networks trained with gradient descent converge to a global minimum. 
%Seminal contributions by \cite{du2019gradientflow} and \cite{Arora2019Fine}, among others, demonstrated that for wide networks, the empirical Neural Tangent Kernel (NTK) remains well-conditioned throughout training. 
%This ensures the stability of the gradient flow and guarantees convergence. 
%Building on this, works like \cite{Bartlett_Montanari_Rakhlin_2021} and \cite{bachchizat} showed that the network's trajectory stays close to its linearization around initialization. 
%These results are critical as they define the "NTK regime" where training is tractable and behaves linearly.

%\textcolor{teal}{Più condensato -- Guardare su google scholar lavori più citati che citano Du e Arora. Dire che noi invece facciamo stime quantitative senza entrare il dettaglio < 10 righe}
%A foundational line of work has rigorously established that, under suitable overparameterization, gradient-based training of neural networks converges to a global minimum.
%Seminal contributions by \cite{du2019gradientflow} and \cite{Arora2019Fine} demonstrated that, for sufficiently wide networks, the empirical NTK remains well-conditioned throughout training, ensuring convergence via kernel regression. 
%Building on this, works like \cite{allen2019_convergence} provided polynomial-width overparameterization guarantees for deep architectures, significantly relaxing previous width requirements. 
%\cite{sankararaman2020_confusion} introduced the concept of gradient confusion to explain how increasing network width reduces such conflicts among stochastic gradients, thereby improving convergence speed. 
%More recent work by \cite{noy2021convergence} offers a specialized convergence theory for deep overparameterized ReLU networks, showing global minima can be achieved in widths quadratic in sample size and linear in depth, with logarithmic time complexity. 
%Additionally, \cite{zhang2019_natgrad} analyzed natural gradient descent in this regime, proving fast convergence under stability of the Jacobian.
% and extending results to adaptive methods like K‑FAC.
%These works consistently reinforce that, in the NTK regime, the network trajectory remains close to its linearization around initialization. 
%Our contributions align with and extend this literature by providing new quantitative bounds on the Wasserstein‑2 distances between trained neural network outputs and their Gaussian process approximations.


A foundational stream of research has shown that, under sufficient overparameterization, gradient-based training of neural networks converges to a global minimum. 
Seminal results by \cite{du2019gradientflow} and \cite{Arora2019Fine} established that for wide two‑layer networks with ReLU activation, the empirical NTK remains well-conditioned, enabling convergence via kernel regression. 
Subsequent advances generalized these results to deep architectures in different directions, such as \cite{allen2019_convergence,zou2019_improved,sankararaman2020_confusion,wu2019,wei_lee_liu,Zou2020},
which provide guarantees that hold with high probability over parameter initalization.
%proved that SGD achieves global convergence in polynomial time under width scaling polynomial in depth and sample size~\cite{allen2019_convergence}. 
%Zou and Gu (2019),
% further relaxed width requirements and sharpened convergence rates~\cite{zou2019_improved}. 
%Sankararaman et al.~(2020),
% introduced the notion of “gradient confusion” and showed that wider networks reduce such conflicts, accelerating SGD~\cite{sankararaman2020_confusion}. 
%Wu et al.~(2019),
% established convergence guarantees for adaptive methods such as AdaGrad in overparameterized regimes~\cite{wu2019_adagrad}. 
These works reinforce that in the NTK regime, the network trajectory stays close to its linearization around initialization. 
%Beyond the strict NTK regime, \cite{li_liang18} and  \cite{li_soltanolkotabi_oymak} build on \cite{Arora2019Fine} to extend the global convergence guarantees outside the NTK regime, 
%showing that 
%global convergence is not only a consequence of width (NTK) but can also exploit data geometry and optimization dynamics as a form of regularization.
%optimization dynamics act as a form of regularization complementary to the NTK theory
Our contributions align with this body of work 
%in the NTK regime 
and further extend this literature by providing novel finite-sample quantitative bounds on the Wasserstein‑2 distance between neural network outputs and their Gaussian process approximations.
%Unlike previous efforts in this direction, our results take into account the tails of the distributions.


%Finally, while many recent theoretical works adopt the NTK parametrization --- rescaling weights to facilitate convergence proofs --- we follow the standard parametrization, in line with \cite{BasteriTrevisan2023, trevisan2023wide}, to ensure compatibility with standard training practices and facilitate interpretation of the bounds in practical settings.

\begin{comment}
    We provide explicit computations for the theoretical results in  Lee et al. \cite{Lee2020} for the quadratic loss and expand them for the choice of L2 regularization.

On the other hand, we build on the very recent work by Basteri and Trevisan \cite{BasteriTrevisan2023}, \cite{trevisan2023wide} that quantitatively characterizes the convergence of a deep, fully connected neural network to its associated Gaussian process. We show a result analogous to their main theorem for the joint probability of the neural network and the NTK in terms of the $p$-Wasserstein distance.
We make use of the technical properties regarding this metric that are presented in the aforementioned articles.
Moreover, we generalize the main theorem in \cite{BasteriTrevisan2023}, which regards the neural network at initialization, to any arbitrary training time $t$ and arbitrary choice of the trained layers.
Alternatively, recent work by Favaro et al. \cite{Favaro2025} provides quantitative bounds to the Gaussian approximation of the neural network for finite widths.

Our Theorems \ref{thm:convergence}, \ref{thm:convergence2} and \ref{thm:convergence_general} on convergence when $t$ diverges depend on the eigenvalues of the limiting kernel.
We make use of a result by Jacot et al. \cite{jacot} to ensure that the smallest eigenvalue of the limiting kernel is bounded away from zero under the hypothesis of compactness on the training data.
Alternatively, Karhadkar et al. \cite{montufar} provide explicit and sharp lower bounds for the smallest singular value of the feature matrix for deep ReLU networks. Bombari et al. \cite{bombari2022memorization} do it for general smooth activations under mild hypothesis on the widths' growth and on the distribution of the training data.

Although many recent works such as \cite{jacot, arora, nguyen, Lee2020, du2019gradientdescent} choose the so-called NTK parametrization (see Section \ref{sec:mainresult}) we present our results in the standard parametrization, in concordance with the authors of \cite{trevisan2023wide, BasteriTrevisan2023}.
The results we borrow from \cite{Lee2020} have been proven to hold both in the standard and NTK parametrization.
\end{comment}

\subsection{Structure of the paper}

Section \ref{sec:math-setting} introduces our notation and mathematical preliminaries. 
In Section \ref{sec:mainresult}, we present our primary theoretical contributions, including our main quantitative convergence theorem. 
The key technical proofs and intermediate results are outlined succinctly,
referring the reader to the relevant lemmas in the Supplementary Material. 
Numerical experiments validating our theoretical predictions appear in Section \ref{sec:experiments}. 
Section \ref{sec:discussion} discusses implications and future research directions.

\section{Notation}
\label{sec:math-setting}

In the following, given a matrix $A \in \R^{p \times q}$ we will denote by $\norm{A}{}$ its Frobenius norm and by $\norm{A}{op}$ its operator norm.
$A_{i\_} \in \R^q$ will denote the $i$-th row of $A$ and $A_{\_ j} \in \R^p$ will denote the $j$-th column of $A$, for $1 \le i \le p$ and $1 \le j \le q$.
The symbol $\cdot$ denotes the usual matrix product.
$\sigma_\textup{min}(A), \sigma_\textup{max}(A)$ are the smallest and largest singular value of $A$; and if $p=q$, $\lammin(A)$ and $\lammax(A)$ denote the smallest and largest eigenvalue of $A$, respectively.
%For any function $\phi_t$ depending on time, $\frac{\partial}{\partial t} \phi_t$ will denote its time derivative; and 
For any vector-valued function $f$,
$f(z)_u$ denotes the $u$-th coordinate of $f(z)$.
For any Polish metric space $X$, $\mathcal P(X)$ will denote the space of Borel probability measures on $X$.

\subsection{Shallow neural networks and associated Gaussian process}
\label{subseq:notation}

%Let $n_0$ be the input dimension.
We consider a fully connected, shallow (i.e. single hidden layer) neural network of \emph{width} $n_1$ and input dimension $n_0$.
We assume the output dimension $n_2$ to be equal to $1$ for simplicity.
The output of the neural network as a function of its parameters is given by:
\[
f(x;\theta) = \frac{1}{\sqrt{n_{1}}}\Phi\left(\frac{1}{\sqrt{n_0}}x \theta^{(0)}\right) \theta^{(1)} \in \R,
\]
where $\theta^{(0)} \in \R^{n_0 \times n_1}$ and $\theta^{(1)}\in \R^{n_1}$ denote the inner and outer (respectively) \emph{weights} or \emph{parameters}, 
$\Phi(z)$ is the \emph{activation function}, which acts entrywise on its input,
and $x\in \R^{n_0}$ from now on denotes a test input.
Note that our model implicitly covers neural networks with biases $b^{(0)}\in\R^{n_1}, b^{(1)}\in \R$, by substituting the input $x$ with $(x,1)$, using the parameters $\tilde \theta^{(0)} = (\theta^{(0)},b^{(0)}) \in \R^{(n_0+1)\times n_1}, \tilde \theta^{(1)} = (\theta^{(1)},b^{(1)}) \in \R^{n_1+1}$ and using the activation function $\tilde \Phi(z) = (\Phi(z),1)$.
We will denote by $N = n_0n_1+n_1$ the total dimension of the parameters.
For any ordered set of inputs $X=(x_1, \ldots, x_d)\in \R^{n_0\times d}$ we will use the notation $f(X;\theta) = (f(x_1;\theta), \ldots, f(x_d;\theta)) \in \R^{n_0 \times d}$.
%$\theta^{(\ell)}_t = (W_t^{(\ell)}, b_t^{(\ell)})^\top \in \R^{(n_\ell+1) \times n_{\ell+1}}$, with $W^{(\ell)}_t \in \R^{n_\ell \times n_{\ell+1}}$ and $b^{(\ell)}_t \in \R^{n_{\ell+1}}$ denoting the \emph{weights} and \emph{biases} of the network at layer $\ell \in \{0,1\}$ and time $t$.
%We will refer to the matrix $\theta^{(\ell)}_t$ as the \emph{parameters} at layer $\ell$ and time $t$.
%The entries 
%$w_{ij}^{(\ell)}$ and $b_{i}^{(\ell)}$ of the weights and biases at initialization $t=0$ are drawn i.i.d. standard Gaussian.
%$\theta^{(\ell)}_{ij}$ are drawn independent and identically distributed (i.i.d.) standard Gaussian at initialization.
%The function $\Phi(z) = (\sigma(z),1)$, where $\sigma \colon \R \to \R$ acts entrywise on its input is called the \emph{activation function} or non-linearity.
%We impose $\Phi$ (equivalently, $\sigma$) to be Lipschitz continuous and not linear.
%We impose $\Phi$ (equivalently, $\sigma$) to be Lipschitz continuous and not linear.
In what follows, parameters $\theta^{(0)}_{ij}, \theta^{(1)}_j$, for $1 \le i \le n_0$ and $1 \le j \le n_1$, are drawn independent and identically distributed (i.i.d.) from standard Gaussian random variables at initialization.

%The parameters are initialized as follows:
%\begin{ass}
%    \label{ass:gaussian_initialization}
%    The parameters $\theta^{(0)}_{ij}, \theta^{(1)}_j$, for $1 \le i \le n_0$ and $1 \le j \le n_1$, are drawn independent and identically distributed (i.i.d.) from standard Gaussian random variables at initialization.
%\end{ass}

We will denote by $h_i$ the \emph{preactivation} of $i$-th hidden neuron, for $1\le i\le n_1$:
\[
    %f^{(1)}_i(x;\theta_t) = \frac{1}{\sqrt{n_0}}(x,1) (\theta_t^{(0)})_{\_ i} \in \R^{n_1}.
    h_i(x;\theta) = \frac{1}{\sqrt{n_0}}x (\theta^{(0)})_{\_ i} \in \R.
\]

%\color{lightgray}
%We will denote by $h_i$ the \emph{preactivation} of $i$-th hidden neuron, for $1\le i\le n_1$:
%\[
%    %f^{(1)}_i(x;\theta_t) = \frac{1}{\sqrt{n_0}}(x,1) (\theta_t^{(0)})_{\_ i} \in \R^{n_1}.
%    h_i(x;\theta) = \frac{1}{\sqrt{n_0}}x (\theta^{(0)})_{\_ i} \in \R^{n_1}.
%\]

%Sometimes, it will be more convenient to deal with scalar quantities, and we will write
%\[
%    h_{ij}(x;\theta) = \frac{1}{\sqrt{n_0}}x_i (\theta^{(0)})_{ij} \in \R,
%\]
%for $1 \le i \le n_0$ and $1 \le j \le n_1$.
%, with the convention $x_{n_0+1}=1$.

Now we introduce the \emph{Gaussian approximation} $G$ of the neural network $f$ as the centered Gaussian process associated to the covariance operator $\mathcal K$ given by:
\begin{align*}
    \tilde {\mathcal K}(x,x') &= \frac{1}{n_0}x^\top x', \\
    \mathcal K(x,x') &= \E_{(u,v) \sim \mathcal N(0,\mathcal T(x,x'))}[\Phi(u)\Phi(v)],
    %\mathcal K(x,x') &= \E_{g \sim \mathcal N(0,\tilde {\mathcal K})}[\Phi(g(x))\Phi(g(x'))].
\end{align*}
with \[
\mathcal T(x,x') = \begin{pmatrix}
    \tilde {\mathcal K}(x,x) & \tilde {\mathcal K}(x,x') \\
    \tilde {\mathcal K}(x',x) & \tilde {\mathcal K}(x',x')
\end{pmatrix}.\]

Explicit convergence rates for the Gaussian approximation at initialization can be found in \cite{BasteriTrevisan2023,Favaro2025}.

\subsection{Training, empirical NTK and limiting kernel}

Let $\mathcal D = \{(x_i,y_i)\}_{i=1}^n \subset \R^{n_0} \times \R$ be a given dataset.
Denote by $\X = (x_1, \ldots, x_n) \in \R^{n_0 \times n}$ the vector of training inputs, and by $y = (y_1, \ldots, y_{n}) \in \R^{n}$ the vector of outputs.
From now on, we will consider the parameters $\theta = \theta_t$ to evolve on training time.
% and we will use the notation $f_t(x) = f(x;\theta_t)$ for any $t \ge 0$. 
%We will use the symbols $\ktrain, \ktest$ and $\kttest$ to denote the objects $(k^{(L)}(x_i,x_j))_{i,j=1}^n \in \R^{n \times n}$, $(k^{(L)}(x,x_j))_{j = 1, \ldots, n} \in \R^{1 \times n}$ and $(k^{(L)}(x,x))_{i,j=1}^m \in \R$, respectively.
Consider the empirical risk for the mean squared error loss: 
\[\mathcal R_\mathcal{D}[f,\theta] = \frac{1}{2}(f(\X;\theta)-y)^\top(f(\X;\theta)-y).\]
Continuous time gradient descent with respect to this objective function yields the following dynamics for the parameters and the network:
\begin{align}
    \begin{split}
        \label{pre_eq1}
    \frac{\partial}{\partial t} \theta_t &= - \nabla_{\theta} \mathcal R_\mathcal{D}[f,\theta_t] = - \nabla_{\theta}f(\X;\theta_t) (f(\X;\theta_t) - y),\\
    \frac{\partial}{\partial t} f(\X;\theta) &= -\nabla_{\theta}f(\X;\theta_t)^\top\nabla_{\theta}f(\X;\theta_t) (f(\X;\theta_t) - y),\\
    \frac{\partial}{\partial t} f(x;\theta) &= -\nabla_{\theta}f(x;\theta_t)^\top\nabla_{\theta}f(\X;\theta_t) (f(\X;\theta_t) - y).
    \end{split}
\end{align}

In \cite{Lee2020} the authors showed that the dynamics in \eqref{pre_eq1} converge to those of a linearized network, which we now introduce.

\begin{comment}
    First, we introduce linear networks.
\begin{defn}
    We say that a neural network $f(x;\theta_t)$ is \emph{linear on its parameters} if it is linear as a function of $\theta_t$, 
    that is, there exist functions $a \colon \R^{n_0} \to \R$ and $b \colon \R^{n_0} \to \R^{1 \times N}$ such that:
    \[
        f(x;\theta_t) = a(x,\theta_0)+b(x,\theta_0)(\theta_t-\theta_0)
    \]
    for each $(x,\theta_t) \in \R^{n_0 \times N}$.
\end{defn}

Given a neural network $f$, one can construct a linear approximation $f^\lin$ of $f$.
\end{comment}

\begin{defn}
    Given a neural network $f$ we define its associated \emph{linearized network}:
\[
    f^\lin(x;\theta_t) = f(x;\theta_0)+ \nabla_{\theta} f(x;\theta_0) \vert_{\theta = \theta_0} \omega_t
\]
with the change of parameters $\omega_t = \theta_t-\theta_0$.
In the following, we will also consider the \emph{linearized gradient flow}, given by
\[
    \frac{\partial}{\partial t} {\overline \theta_t} = -\nabla_\theta f(\X;\theta_0) \cdot (f^\lin(\X;\overline \theta_t)-y).
\]
%with $f^\lin_t(x) = f^\lin(x;\overline \theta_t)$ for each $x \in \R^{n_0}$.
\end{defn}

The linearized network is known to approximate arbitrarily well the real training dynamics in the wide limit under some stability conditions and positive-definiteness of the NTK 
(see Theorem 5.1 in \cite{Bartlett_Montanari_Rakhlin_2021} and Theorem 2.2 in \cite{bachchizat}).
In the Supplementary Material we prove an analogue result (Proposition \ref{thm:montanari}) adapted to our setting,
In particular, our result applies to shallow networks with inner and outer weights and for the MSE loss without scaling over the number of training points, as opposed to the cited results.

An alternative formulation of this asymptotic linearzation phenomenon is provided in \cite{jacot}, 
where the neural tangent kernel (NTK) was introduced.
The authors showed that under 
%Assumption \ref{ass:gaussian_initialization}, 
Gaussian initialization,
the dynamics \eqref{pre_eq1} are governed by this operator which we now define in our setting.
Define the \emph{empirical kernel}, or \emph{NTK} $k \colon \R^{n_0} \times \R^{n_0} \times \R^N \to \R$ as:
\[
    k(x,x';\theta) = \nabla_{\theta}f(x;\theta) \nabla_{\theta}f(x';\theta)^\top,
\]
where $\theta \in \R^N$ denotes the matrix of \emph{parameters}.
The empirical kernel at the hidden layer can also be defined as a function of $\R^{n_0} \times \R^{n_0} \times \R^N\to \R^{n_1 \times n_1}$:
\[
    \tilde k_{uv}(x,x';\theta) = \nabla_{\theta^{(0)}}h(x;\theta)_u \nabla_{\theta^{(0)}}h(x';\theta)_v^\top, \quad 1 \le u,v \le n_1,
\]

During training, the parameter $\theta_t$ will be omitted when it is clear from the context and we will simply write $k_t(x,x') = k_t(x,x';\theta_t)$ and $\tilde k_t(x,x') = \tilde k_t(x,x';\theta_t)$.
%Explicit expressions for these kernels are available in Section \textcolor{red}{chissa}.
%The parameters $\theta_t$ and the superscript$\phantom{k}^{(2)}$ will be omitted in the definitions of $k$ and $f$ when they are clear from the context.
%In this case we will also write $f_t(x) = f(x;\theta_t)$ and $k_t(x,x')=k(x,x';\theta_t)$ for any $x \in \R^{n_0}$ and $t>0$.

%The other key object in this article will be the limit of the empirical kernel when the width $n_1$ tends to infinity.
%The empirical kernel converges in law to a limiting kernel.
%The existence of this limiting kernel and the convergence in law of the empirical kernel to its limit were first proven by the authors of \cite{jacot}.
The convergence of the NTK to this limiting kernel was first proven by the authors of \cite{jacot}.
Define the \emph{analytical}, or \emph{limiting kernel} $k_\infty$ as follows:

\[
    k_\infty(x,x') = \mathcal K(x,x') + \tilde{\mathcal K}(x,x') \E_{(u,v) \sim \mathcal N(0,\mathcal T(x,x'))}[\Phi'(u)\Phi'(v)],
\]
with $\mathcal T(x,x') = \begin{pmatrix}
    \tilde {\mathcal K}(x,x) & \tilde {\mathcal K}(x,x') \\
    \tilde {\mathcal K}(x',x) & \tilde {\mathcal K}(x',x')
\end{pmatrix}$.
From now on, let $k_\infty = k_\infty(\X,\X)$ denote the limiting kernel valued on the training set.

%Again, the superscript$\phantom{k}^{(2)}$ will be omitted when it is clear from the context.

Note that, in the linearized regime, Equation \eqref{pre_eq1} can be solved analytically.
%Note that the empirical NTK for any linear networks is independent of time.
%In particular, the NTK for $f^\lin$ coincides with $k_0$ at any time $t \ge 0$.
%This allows to solve the gradient flow equations \eqref{pre_eq1} analytically in the linearized setting.
%In effect, letting $\ktrain = \nabla_{\theta}f_0(\X) \nabla_{\theta}f_0(\X)^\top$ and $\ktest = \nabla_{\theta}f_0(x) \nabla_{\theta}f_0(\X)^\top$,
In effect, letting $\ktrain = \nabla_{\theta}f(\X;\theta_0) \nabla_{\theta}f_0(\X;\theta_0)^\top$ and $\ktest = \nabla_{\theta}f(x;\theta_0) \nabla_{\theta}f(\X;\theta_0)^\top$,
the gradient flow equations \eqref{pre_eq1} for $f^\lin$ can be rewritten as:
\begin{align}
    \begin{split}
        \label{eq1} \frac{\partial}{\partial t} {\overline \theta_t} &= - \nabla_{\theta}f(\X;\theta_0)^\top (f^\lin(\X;\overline \theta_t) - y),\\
        \frac{\partial}{\partial t} {f^\lin}(\X;\overline \theta_t) &= -\ktrain (f^\lin(\X;\overline \theta_t) - y),\\
         \frac{\partial}{\partial t} {f^\lin}(x;\overline \theta_t) &= -\ktest (f^\lin(\X;\overline \theta_t) - y).
    \end{split}
\end{align}

The inverse of matrix $\ktrain$, which is random in the initialization of the network and may not be positive definite for some $\theta$,
 appears in the solution of Equation \eqref{eq1}.
 Thus we introduce the following auxiliary operator:

\begin{defn}
\label{def:It}
    For any symmetric, invertible matrix $B \in \R^{n \times n}$ and for any $t > 0$, define the $n\times n$ real matrix
    \[
        I_t(B) = (\mathbbm{1}_n - e^{-Bt})B^{-1}.
    \]

    Note that $I_t(B)$ is invertible and symmetric since the matrix exponential of $B$ is positive definite.
    The operator $I_t$ can be extended to general symmetric matrices in the following way: define, for each $a \in \R$,
    \[
    I_t(a) = \int_0^te^{-as}ds = \begin{cases}
        \frac{1-e^{-at}}{a} \quad &\textup{if }a \ne 0, \\
        t \quad &\textup{if }a=0.
    \end{cases}
    \]
    
    Let $B\in \R^{n \times n}$ be symmetric, not necesarily non-degenerate.
    Consider the eigenvalue decomposition of $B$, $B = UDU^\top $ with $D=\textup{diag}(a_1, \ldots a_n)$ and $U$ orthogonal, then put $I_t(B) = U \textup{diag}(I_t(a_1), \ldots I_t(a_n)) U^\top $.
    \end{defn}

Lemma \ref{lem_It_properties} shows some properties and well-posedness of the operator $I_t$ defined in Definition \ref{def:It}.
The matrix $I_t(B)$ can be thought of as the analytic continuation of the matrix function $(\mathbbm{1}_n - e^{-Bt})B^{-1}$, for any $B \in \R^{n \times n}$ symmetric.
With this definition, the solution to \eqref{eq1} reads:
\begin{align}
    \label{eq1:soltrain} f^\lin(\X;\overline \theta_t) &=\exp(-\ktrain t)f(\X;\theta_0) + (\mathbbm 1_n - \exp(-\ktrain t))y, \\
    \label{eq1:soltest} f^\lin(x;\overline \theta_t)&= f(x;\theta_0) -\ktest I_t(\ktrain)(f(x;\theta_0)-y).
\end{align}
The computation leading to this expression is contained in Supplementary Material \ref{ap:eq1}.

In \cite{Lee2020}, the authors showed that when $f$ is linear in its parameters, its output distribution at a test point $x \in \R^{n_0}$ converges weakly,
 as $n_1$ diverges, to a Gaussian process $G_t$ with mean and covariance given by:
\begin{align}
    \begin{split}
        \label{characterization}
    \mu_t(x) &= k_\infty(x,\X) I_t(k_\infty)y,\\
    \Sigma_t(x,x')&=\mathcal K(x,x') - \mathcal K(x,\X)I_t(k_\infty)k_\infty(\X,x')- k_\infty(x,\X)I_t(k_\infty)\mathcal K(\X,x')\\
    &\quad + k_\infty(x,\X)I_t(k_\infty)\mathcal K(\X,\X)I_t(k_\infty)k_\infty(\X,x'),
    \end{split}
    \end{align}
for every positive training time $t$.
    %This formula agrees with the formula found by the authors of \cite{Lee2020}, which is a generalization for deep neural networks.
%This characterization was introduced in \cite{Lee2020}.
For the sake of completeness we included a proof for the above formula in Supplementary Material \ref{ap:sub:proofGt}.


The solution of \eqref{eq1} completely characterizes the dynamics of the Gaussian process $G_t$ for any time $t\ge 0$, as a consequence of the Central Limit Theorem.
%Before presenting this characterization we make the following simplifying assumption:
%\begin{ass}
%    \textcolor{teal}{mettere nel testo}
%    \label{ass:positivity_kinf}
%    The limiting kernel $k_\infty(\X,\X)$ is positive definite.
%\end{ass}
%\begin{rem}
%    Assumption \ref{ass:positivity_kinf} is mild and holds in a very general setting.
%    Indeed, the authors of \cite{carvalho2024} showed that when the training data is in general position in $\R^{n_0}$ and $\Phi$ is not a polynomial the smallest eigenvalue of $k_\infty(\X,\X)$ is strictly greater than zero.
%\end{rem}
In the rest of the paper we will assume that $k_\infty(\X,\X)$ is positive definite.
This is a mild assumption and holds in a very general setting.
Indeed, the authors of \cite{carvalho2024} showed that when the training data is in general position in $\R^{n_0}$ and $\Phi$ is not a polynomial the smallest eigenvalue of $k_\infty(\X,\X)$ is strictly greater than zero.


\section{Assumptions and main result}
\label{sec:mainresult}

%Before stating our main result we give a generalization of the main result in \cite{trevisan2023wide} for linear networks positive training time,
%which serves as a middle step in the proof of the next theorem.

In this section we state our assumptions and main result.
%First we introduce two additional assumptions:
\begin{ass}
    \label{ass:gaussian_initialization}
    The parameters $\theta^{(0)}_{ij}, \theta^{(1)}_j$, for $1 \le i \le n_0$ and $1 \le j \le n_1$, are drawn independent and identically distributed (i.i.d.) from standard Gaussian random variables at initialization.
\end{ass}
\begin{ass}
%    \textcolor{teal}{mettere nel testo}
    \label{ass:positivity_kinf}
    The limiting kernel $k_\infty(\X,\X)$ is positive definite.
\end{ass}
%\begin{rem}
%    Assumption \ref{ass:positivity_kinf} is mild and holds in a very general setting.
%    Indeed, the authors of \cite{carvalho2024} showed that when the training data is in general position in $\R^{n_0}$ and $\Phi$ is not a polynomial the smallest eigenvalue of $k_\infty(\X,\X)$ is strictly greater than zero.
%\end{rem}

\begin{ass}
    \label{ass:bounded}
    $\Phi$ and $\Phi'$ are Lipschitz continuous  and bounded.
\end{ass}
\begin{ass}
    \label{ass:r}
    For some fixed $r\ge 5$ the following inequality holds:
    \[
    \frac{4\norm{\X}{}(\sqrt 5\norm{\Phi}{\infty}+\norm{y}{})}{\sqrt{n_1n_0}}
        \left(\norm{\Phi'}{\infty} + \textup{Lip}\Phi + \frac{\norm{\X}{}\textup{Lip}\Phi'\sqrt{r \log n_1}}{\sqrt{n_0}}\right) < \lammin^\infty,
    \]
\end{ass}
\begin{rem}
    Note that these are rather mild assumptions.
    Assumption \ref{ass:positivity_kinf} is mild and holds in a very general setting.
    Indeed, the authors of \cite{carvalho2024} showed that when the training data is in general position in $\R^{n_0}$ and $\Phi$ is not a polynomial the smallest eigenvalue of $k_\infty(\X,\X)$ is strictly greater than zero.
Assumption \ref{ass:bounded} is standard in literature and is
 satisfied by most activation functions, such as the sigmoid function and other logistic activations, 
hyperbolic tangent, Gaussian activation or sinusoid, among others.
The ReLu family is a notable exception, although we expect our result to also hold in that case.
As for Assumption \ref{ass:r},
notice that the left hand side tends to zero as $\min\{n_1,n_0\}$ grows.
This implies that our hypothesis is satisfied for sufficiently large $n_0$ or $n_1$.
In particular, it holds for sufficiently overparametrized networks.
\end{rem}
\begin{rem}
    Assumption \ref{ass:r} may appear somewhat artificial, so we provide an informal intuition on its use.
    This assumption is needed to control the fluctuations of $\norm{k_0-k_\infty}{}$ with the (deterministic) smallest eigenvalue of the limiting kernel.
    This enables the use of Proposition \ref{thm:montanari}, which provides a quenched estimation of the $L^2$ distance between $f$ and $f^\lin$, which in turn plays a central role in the proof of Theorem \ref{thm:main}.
    In particular, the $L^2$ norm of this difference is controlled with a function of the Lipschitz constant of the Jacobian $\nabla_\theta f$ at initialization,
     and this Lipschitz constant is estimated with an expression in terms of  $\norm{k_0-k_\infty}{}$ bounded by the left-hand side in Assumption \ref{ass:r}.
    This is the content of Lemmas \ref{lem:jacobian} and \ref{lem:lammin}.
\end{rem}

To measure how well the Gaussian process $G_t$ approximates the network $f$ at time $t$, we will use a well-known family of metrics between probability distributions, the Wasserstein distances:
\begin{defn}
    Let $p \in [1, \infty]$ and let $\mu, \nu$ be two probability measures defined on a Polish space $(M,d_M)$ with finite $p$-moment.
    The $p$-Wasserstein distance between $\mu$ and $\nu$ is given by
    \[
    \mathcal W_p (\mu, \nu) = \inf_{\gamma \in \Gamma(\mu, \nu), X\sim \mu, Y \sim \nu} \left( \E_\gamma \left[ d_M(X,Y)^p \right] \right)^{\frac{1}{p}},
    \]
    where $\Gamma(\mu, \nu)$ denotes the set of joint probability measures $\gamma$ defined on $M \times M$ with marginal laws $\mu$ and $\nu$.
    With a slight abuse of notation, we will often write $\mathcal W_p (X, Y) =\mathcal W_p (\mu, \nu)$ for any $X\sim \mu$ and $Y\sim \nu$.
\end{defn}

Now we can state our main theorem:

%\color{red}
%We collect all the assumptions for our main theorem here for clarity:
%\begin{enumerate}
%    \item The parameters $\theta^{(0)}_{ij}, \theta^{(1)}_j$, for $1 \le i \le n_0$ and $1 \le j \le n_1$, are drawn independent and identically distributed (i.i.d.) from standard Gaussian random variables at initialization.
%    \item $\Phi$ and $\Phi'$ are Lipschitz continuous  and bounded.
%    \item For some fixed $r\ge 5$ the following inequality holds:
%    \[
%    \frac{4\norm{\X}{}(\sqrt 5\norm{\Phi}{\infty}+\norm{y}{})}{\sqrt{n_1n_0}}
%        \left(\norm{\Phi'}{\infty} + \textup{Lip}\Phi + \frac{\norm{\X}{}\textup{Lip}\Phi'\sqrt{r \log n_1}}{\sqrt{n_0}}\right) < \lammin^\infty,
%    \]
%\end{enumerate}

\color{black}
\begin{thm}
    \label{thm:main}
    Under Assumptions \ref{ass:gaussian_initialization},
    \ref{ass:positivity_kinf},
    \ref{ass:bounded} and \ref{ass:r},
    for each test point $x\in \R^{n_0}$ there exist positive constants $a_1$ and $a_2$ not depending on $n_0,n_1$ nor $t$ such that:
    \[
        \mathcal W_2^2(f(x;\theta_t),G_t(x)) \le r\left(\frac{a_1\log n_1}{(\lammin^\infty)^3n_1n_0}+\frac{a_2  n_0}{(\lammin^\infty)^rn_1^\frac{r}{4}}(1+t^8)\right).
    \]
    Here $r$ is the constant appearing in Assumption \ref{ass:r}.
\end{thm}

\begin{rem}
    \label{rem:time}
    Note that our result is not limited to fixed training time $t$, but holds for values of $t$ growing polynomially on $n_1$.
    Indeed, provided that $t$ grows at most polynomially in $n_1$, the constant $r$ can be chosen arbitrarily big making the term dependent on time negligible.
    In particular, as long as $t$ grows polynomially on $n_1$, a sufficiently big $r$ can be chosen so that the right-hand side in Theorem \ref{thm:main} tends to zero as $n_1$ diverges.

    The term $t^8$ in the right hand side of the inequality is due to Lemma \ref{lem:gronwallbeta} in Supplementary Material \ref{ap:auxiliary}.
    Lemma \ref{lem:gronwallbeta} provides upper bounds of the entries $\theta^{(0)}_t$ and $\theta^{(1)}_t$
    that account for perturbations that occur on tail events with respect to the initalization distribution (i.e. in the ``bad event" $S^C$).
    A finer control is possible if one is interested in a result that holds $S$ on only, which has high probability, such as the ones in \cite{Bartlett_Montanari_Rakhlin_2021,bachchizat}. 
    This finer control corresponds to Theorem \ref{thm:montanari}.
    %by integrating the gradient flow equations.
    %The exponent is due to products between $\theta^{(0)}_t$ and $\theta^{(1)}_t$ resulting from applying Hölder and Young inequalities.
\end{rem}

\subsection{Sketch of proof of Theorem \ref{thm:main}}
\label{subsec:sketch}
The proof of our main theorem is as follows: we bound by triangle inequality

\[
    \mathcal W_2(f(x;\theta_t),G_t(x)) \le \mathcal W_2(f(x;\theta_t),f^\lin(x;\overline \theta_t))+\mathcal W_2(f^\lin(x;\overline \theta_t),G_t(x)),
\]
%\textcolor{red}{where $\overline{f_t}(x) = f^\lin(x;\bar \theta_t)$,}
and proceed to control the two terms appearing in the right-hand side separately.

To bound the first summand, we partition $\R^N$ into a ``good" event $S \subset \R^N$, in which the assumptions of Proposition \ref{thm:montanari} hold, along some other concentration properties of the parameters,
and a ``bad" event $S^C$ in which they do not; so the estimation of the first summand reduces to the estimation of integrals over $S$ and $S^C$, respectively.
Moreover, as $n_1$ diverges, $\mathbb P(S)$ converges to 1.
Proposition \ref{thm:montanari} consists on an upper bound of $\norm{f(x;\theta_t)-f^\lin(x;\overline \theta_t)}{}^2$ on this ``good event''; which is a version of Theorem 5.1 in \cite{Bartlett_Montanari_Rakhlin_2021} adapted to our setting.
This allows to bound the first integral.
In $S^C$ the strategy is to show that the measure of $S^C$ decreases faster than how the upper bounds in Theorem \ref{thm:stime_polinomiali} grow.
Theorem \ref{thm:stime_polinomiali} provides estimations for $\norm{f(x;\theta_t)-f^\lin(x;\overline \theta_t)}{}^2$ that are rougher than the ones in Proposition \ref{thm:montanari} in the sense that do not vanish in the wide limit, but hold in $S^C$.
We estimate the second integral by partitioning $S^C$ into countably many discs parametrized by $\gamma \in \N$, and summing over $\gamma$ while exploiting concentration inequalities that hold on each disc.
The result of this estimation is summarized in the following Theorem:

\begin{prop}
    \label{quasi_main}
    On the hypothesis of Theorem \ref{thm:main}, there exist positive constants $a_1$ and $a_2$ not depending on $n_0,n_1$ nor $t$ such that:
    \[
        \mathcal W_2^2(f(x;\theta_t),f^\lin(x;\overline \theta_t)) \le r\left(\frac{a_1\log n_1}{(\lammin^\infty)^3n_1n_0}+\frac{a_2  n_0}{(\lammin^\infty)^rn_1^\frac{r}{4}}(1+t^8)\right).
    \]
\end{prop}

The dependence on time of the right-hand side of statement of the Theorem comes from Lemma \ref{lem:gronwallbeta}.
In particular, in the ``bad'' event $S^C$ a lower bound of the smallest eigenvalue of the random matrix $k_0$ is not available, 
and hence by Definiton \ref{def:It} the sharpest upper bound for $\norm{I_t(k_0)}{}$ is $t$.

The second summand, instead, is estimated with the following result:

\begin{prop}
    \label{prop:main_linear}
    Let $f^\lin$ be the linearization of $f$, and let $x\in \R^{n_0}$ be a test point.
    Then, under assumptions \ref{ass:gaussian_initialization} and \ref{ass:bounded}, there exist positive constants 
    %$C, \overline C, \overline D$ 
    $C$ and $C'$
    not depending on $n_1$ nor $t$ such that:
    \begin{align*}
        \mathcal W_2^2(f^\lin(\X;\overline \theta_t),G_t(\X)) &\le \frac{C}{n_1},\\
        \mathcal W_2^2(f^\lin(x;\overline \theta_t),G_t(x)) &\le \frac{C'}{n_1}(1+t^2).
        %\mathcal W_2^2(f^\lin(\X;\overline \theta_t),G_t(\X)) &\le \frac{1}{n_1}C (t+1)e^{-\lammin^\infty t},\\
        %\mathcal W_2^2(f^\lin(x;\overline \theta_t),G_t(x)) &\le \frac{1}{n_1}(\overline C + \overline D(t+1)e^{-\lammin^\infty t}).
    \end{align*}
    \end{prop}
    
    The first statement in Proposition \ref{prop:main_linear} is proven by bounding $\frac{\partial}{\partial t}\norm{f^\lin(\X;\overline \theta_t) - G_t(\X)}{}^2$ with Young's inequality and gradient flow equations.
    Next, we apply Grönwall and Hölder's inequalities to decompose the problem in some expected values of $L^2$ and $L^4$ norms of the differences between kernels and between $f^\lin$ and $G$, both at initialization.
    Lastly, the main result in \cite{BasteriTrevisan2023} and Proposition \ref{prop:kernelp} providing $L^p$ estimations for the difference between the empirical NTK and the limiting kernel complete the proof.   
    The proof for the second statement in Proposition \ref{prop:main_linear} uses the bound for training points in a triangular system of differential inqualities, inherited by the solution of Equation \eqref{eq1}.
    %After that, by integration and Hölder inequality the proof finishes in an analogous way to the proof of the first statement.
    
Theorem \ref{thm:main} and Propositions \ref{quasi_main} and \ref{prop:main_linear} are proven in full detail in the Supplementary Material \ref{ap:proof_1} and \ref{ap:proof_2}.

\section{Numerical Experiments}
\label{sec:experiments}

We conduct some numerical experiments in Figure \ref{fig:experiment} to support our theoretical results.
In both experiments, $t$ is taken as the product of the learning rate and the number of iterations or epochs.
For simplicity, we considered training and test inputs on the real line.
The training set and test set were drawn from a uniform distribution on a fixed interval, 
and the labels $y$ of the training points correspond to a sine function with additive noise.
The code is available at \url{https://github.com/emosig/quantitative_gaussian_trainedNN}.

\subsection{Experiment 1: Gaussian approximation of $f$}

The leftmost and center plots in 
%The plots in
Figure \ref{fig:experiment} represent 100 trained shallow neural networks with sigmoid activation of width $n_1=700$ on the leftmost plot and $n_1=1000$ in the central plot.
The networks have been trained  for $2\cdot10^4$ epochs with learning rates of $\frac{1}{700}$ and $\frac{7}{1000}$ (hence $t \approx 28.571$ on the leftmost plot and $t=140$ in the central one) to fit two training points, corresponding to different random seeds on each case.
Together with the networks, the plot depicts the mean (in black) of $G_t$ and a 95\% confidence interval (in grey) over 200 equally spaced test points on the interval $[-10,10]$.
%The activation chosen for this experiment was the sigmoid function, 
%the left plot corresponds to a network with $n_1=500$ and the right one to a network with $n_1=100$.
%Both networks have been trained for $10^4$ epochs with a learning rate of $0.1$.
%and in Figure \ref{fig:experiment1} we show two instances of the experiments for 300 iterations and $n_1=100$ (left) and 500 iterations and $n_1 = 200$ (right).
%The networks were programmed with \emph{TensorFlow} \cite{tensorflow2015-whitepaper} and the Gaussian process $G_t$ was constructed using the library \emph{neural tangents} \cite{neuraltangents2020} to construct the kernels $\mathcal K$ and $k_\infty$, 
The networks were programmed with \href{https://download.pytorch.org/whl/cu124/torch-2.6.0%2Bcu124-cp311-cp311-linux_x86_64.whl}{PyTorch 2.6.0} \cite{pytorch}, 
 and the Gaussian process $G_t$ was constructed using the library \emph{neural tangents 0.6.5} \cite{neuraltangents2020} for the kernels $\mathcal K$ and $k_\infty$, 
needed to construct $\mu_t$ and $\Sigma_t$ in \eqref{characterization}.
The operator $I_t(B)$ was programmed by solving the linear system of equations $BX = \mathbbm 1_n - e^{-B t}$ with the \emph{linalg} package of \emph{NumPy} \cite{numpy}.

\subsection{Experiment 2: $\mathcal W_2(f(x;\theta_t),G_t(x))$ decays with $n_1$ for any $x$}

In our second experiment (Figure \ref{fig:experiment}, right), we compute the quadratic Wasserstein distance between the trained shallow network and $G_t$ for a variety of widths, ranging from 2 to 256.
To do this we drew $10^4$ samples of $G_t$, given by the mean $\mu_t$ and covariance $\Sigma_t$ in \eqref{characterization}, which were computed with the \emph{neural tangents 0.6.5} library \cite{neuraltangents2020}.
Then, we trained indepently over a single training point, for 100 epochs and with a learning rate of $0.1$ (hence $t=10$), $10^4$ neural networks for each width and then calculated the empirical Wasserstein distance with the \emph{Python Optimal Transport 0.9.5} library \cite{pot}. 
%\textcolor{red}{The rightmost image corresponds to the same experiment for 500 epochs and with a learning rate of $0.01$, together with the power-law fit in red.
%For this experiments, the widths larger than 100 were excluded, in accordance to the discussion in \ref{paragraph:discussion}.}
As in Experiment 1, \emph{PyTorch} was used to construct the neural networks, and the activation chosen for $f$ and $G$ is once again the sigmoid.
%This plot suggests that our result holds too for ReLu activation, and thus Assumption \ref{ass:bounded} could be weakened.

\begin{figure}
    \centering
    \includegraphics[width=.32\textwidth]{seed3_700_1_20000}%
    \hspace{0.02in}% add some horizontal spacing
    \includegraphics[width=.32\textwidth]{seed5_1000_07_30000}%
    \hspace{0.02in}
    \includegraphics[width=.32\textwidth]{sim2_3}
    \caption{The Gaussian process approximates the neural networks during training (left and center images), and it converges in 2-Wasserstein space to $f_t$ (right image). 
    On the rightmost image, the blue points represent the empirical Wasserstein distance between $f$ and $G$ for increasing widths, and the red plot is the power-law fit between the blue points.}
    %\caption{The Gaussian process approximates the neural networks during training.}
    %Given two training points, the above graphs show a number of trained neural networks after 300 epochs, together with the mean (black) and standard deviation (gray) of the associated Gaussian process $G_t$.
    %On the left the choice for $(\sigma, n_1)$ is (ReLu, $100$), and on the right it is (sigmoid, $200$).
    %The choice of the activation function and the width of the hidden layer have an impact on convergence.} 
    \label{fig:experiment}
 \end{figure}

 %\begin{figure}
  %  \centering
   % \includegraphics[width=.4\textwidth]{sim2}
    %\hspace{0.05in}
    %\includegraphics[width=.4\textwidth]{sim2_2}
    %\caption{Convergence of $G_t(x)$ to $f(x;\theta_t)$ in 2-Wasserstein space. The rightmost experiment includes a power-law fit in red.}
        %Given two training points, the above graphs show a number of trained neural networks after 300 epochs, together with the mean (black) and standard deviation (gray) of the associated Gaussian process $G_t$.
    %On the left the choice for $(\sigma, n_1)$ is (ReLu, $100$), and on the right it is (sigmoid, $200$).
    %The choice of the activation function and the width of the hidden layer have an impact on convergence.} 
    %\label{fig:experiment2}
 %\end{figure}

\paragraph{Discussion on the choice of $n_1$ and the number of samples}
\label{paragraph:discussion}

The authors of \cite{Fournier2015} provide an estimation of the error between a probability measure $\mu$ and its empirical counterpart $\hat \mu_N = \frac{1}{N}\sum_{i=1}^N\delta_{X_i}$, 
given an i.i.d. sequence $(X_i)_{i=1}^N$, $X_i \sim \mu$.
More precisely, Theorem 1 in \cite{Fournier2015}, for $p = 2$ and $n_0=1$, provided that $\mu$ has finite third moment, reads:
\[
\E[\mathcal W_2^2(\mu,\hat\mu_N)] \le \frac{c_1}{\sqrt{N}},
\]
for some constant $c_1$.
This estimation applied to $f$ and $G$ can be combined with our main result to find a minimal ratio of samples per width needed for our numerical estimations to be noise-free.
There exist constants $c_2,c_3$ not depending on $n_1, N$ nor $t$ such that:
\[
    \mathcal W_2^2(f(x;\theta_t),G_t(x)) \approx \mathcal W_2^2(\widehat{f(x;\theta_t)}_N,\widehat{G_t(x)}_N) + \frac{c_2}{\sqrt{N}} \le \frac{c_3 \log n_1}{n_1}.
\]
Hence our computations make sense when $N \gg \left(\frac{n_1}{\log n_1}\right)^2$. 
For $N = 10^4$, this means an upper bound for the widths for which our experiments make sense is, approximately, $650$.
%In effect, in Figure \ref{fig:experiment}, rightmost image, the noise in the signal can be appreciated as the width $n_1$ grows.

\section{Discussion}
\label{sec:discussion}

In this paper we provided quantitative convergence rates in $2$-Wasserstein distance between trained shallow neural networks with standard Gaussian initialization and an appropriate Gaussian process, for any positive training time.
This was proven for Lipschitz, bounded activations with bounded derivative and for sufficiently large hidden layer width.
We now address some limitations of our work and possible future research directions:
\begin{enumerate}
    \item Our main result is not uniform in time.
        Although the dependence on time can be minimized at the price of including a sufficiently big multiplicative constant in the right-hand side of our inequality as discussed in Remark \ref{rem:time},
        a general result holding uniformly in $t>0$, in the limit when $t$ tends to infinity exponentially on $n_1$ is not available. 
        %\textcolor{red}{Furthermore, we have no evidence that the time dependence $t^8$ is sharp.}
        %\textcolor{teal}{Why is this important? Perché non è un "cambiamento di fase" a Feature learning?  cit. Huang Yau. Parlare dell'evento coda dove non si ha controllo della dinamica - Potrebbe esser legato a feature learning}
        
        %Moreover, we have no evidence that the $t^6$ dependence in our estimate is sharp. 

        %Understanding the behavior of the dynamics in this long-time regime is particularly intriguing, as it may involve rare events in the tails of the distribution where our current analysis provides no control. 
        This dependence on time could be related to the transition from the NTK regime to a feature-learning regime, as suggested by the work of \cite{huang_yau}. 
        Their analysis, however, does not address the tails of the distributions, which in our proof correspond to the set $S^C$ and are responsible for the $t^8$ scaling.
        Moreover, \cite{yang4} show that under standard and NTK parameterizations, wide networks cannot perform feature learning in the infinite-width limit. 
        %Their work reveals a fundamental dichotomy: a network either behaves like a kernel machine or evolves its representations during training, but not both.
        
        This suggests that our observed $t^8$ scaling might reflect the boundary of the NTK regime: 
        in the ``bad event'' $S^C$ or for sufficiently large times, the training dynamics may drift into feature-learning, where purely kernel-based control breaks down.
        Our main result remains consistent with works such as \cite{Bartlett_Montanari_Rakhlin_2021,bachchizat}, which hold with high probability, whereas our analysis explicitly incorporates the contribution of the event $S^C$.
        At present, it is unclear whether the $t^8$ scaling is sharp.
        
        We would like to address this problem in future work.

    \item The bound in our Theorem \ref{thm:main} depends on the test point $x$.
    This dependence is explicitly stated on the proof of the auxiliary results Proposition \ref{prop:main_linear} and Theorem \ref{thm:stime_polinomiali} in the Supplementary Material.
    Locally uniform bounds on the test point $x$ might follow from functional inequalities such as the ones found by \cite{Favaro2025} if extended to the NTK regime.
    \item We conjecture that our main result remains valid even without Assumption~\ref{ass:bounded}, as suggested by our numerical experiments with the ReLU activation. 
    In this work, we deliberately focused on a specialized setting with mild hypotheses to obtain a novel and technically precise result while maintaining a clear exposition. 
    Future research will aim to relax the regularity assumptions on the activation and extend our analysis to a more general setting.
    %\textcolor{teal}{Anche un problema di regolarità (LIpschitzianità del jacobiano)}
    %\textcolor{teal}{Risultato tecnico --> ci siamo ristretti a un setting semplificato per chiarezza.}
    
    
    \item When Assumption \ref{ass:gaussian_initialization} 
    and \ref{ass:positivity_kinf} 
    hold, $k_\infty(\X,\X)$ is strictly positive definite and $\Phi$ is Lipschitz,
     the rate of convergence at initialization found in \cite{trevisan2023wide} for the squared $2$-Wasserstein distance is of $n_1^{-2}$.
        This fact suggests that, in the proof of our main Theorem, a better estimation of $\mathcal W_2(f(x;\theta_t),f^\lin(x;\overline \theta_t))$ can be found;
        either by improving Proposition \ref{quasi_main} or
        by improving the estimation of the Lipschitz constant and the norm of the Jacobian $\nabla_\theta f(x;\theta_0)$ in the ``good event'' of the proof of Theorem \ref{thm:main},
        possibly by choosing a more restrictive ``good event'' that still makes the infinite sums in the proof of Theorem \ref{thm:main} converge.
        This last possibility calls for finer concentration inequalities for the parameters and the difference between the empirical NTK and its infinite-width limit.
    \item Our results could be extended to deep, fully connected neural networks, as done for initialization in \cite{trevisan2023wide}, \cite{BasteriTrevisan2023} and \cite{Favaro2025}.
        We hypothesize that 
        %(\textcolor{red}{Oppure ``we are currently working on''? Questo in per il teorema \ref{prop:main_linear} è già fatto nel file vecchio, manca la parte di montanari.}) 
        Proposition \ref{prop:main_linear} can be easily adapted to the deep setting exploiting recursive characterizations of $k_t$ and $k_\infty$ as the ones available in \cite{jacot}, \cite{nguyen} or \cite{Lee2020}.
        The other half of our proof, though, relies in Proposition \ref{quasi_main}, which would need a new proof for the deep setting.
    \item Another desirable step could be to study how well our result extends to other architectures, such as convolutional neural networks or the more modern attention-based architectures.
        This approach is present in \cite{yang2021tensor} but to the best of our knowledge no quantitative results in this direction are available.
    %\item We believe our result to hold in $p$-Wasserstein distance, for any finite $p$; and similar results could be obtained for other metrics between probability measures as the ones used in \cite{Favaro2025}.
\end{enumerate}


%========================================================================
%========================================================================
%========================================================================
\section{Acknowledgements}
\label{sec:acknowledgements}

E.M.G.\ and D.T. are members of GNAMPA group of the Istituto Nazionale di Alta Matematica (INdAM). 
All authors acknowledge the support of GNAMPA Project CUP E53C22001930001.
E.M.G also gratefully acknowledges the hospitality and support of the Institute of Mathematical Statistics and Actuarial Sciences at University of Bern during part of this work.
This research was supported by a Ph.D.\ fellowship funded by the Italian Ministry of University and Research (MUR) under the PNRR program “Transizioni digitali e ambientali (TDA)” (D.M.\ n.~118/2023, CUP~I51J23000400007), within the project “Limiti di scala di dinamiche stocastiche.”

D.T.\ acknowledges the MUR Excellence Department Project awarded to the Department of Mathematics, University of Pisa, CUP I57G22000700001, 
the HPC Italian National Centre for HPC, Big Data and Quantum Computing - Proposal code CN1 CN00000013, CUP I53C22000690001, 
the PRIN 2022 Italian grant 2022WHZ5XH - ``understanding the LEarning process of QUantum Neural networks (LeQun)'', CUP J53D23003890006, 
the project  G24-202 ``Variational methods for geometric and optimal matching problems'' funded by Università Italo Francese.
Research also partly funded by PNRR - M4C2 - Investimento 1.3, Partenariato Esteso PE00000013 - "FAIR - Future Artificial Intelligence Research" - Spoke 1 "Human-centered AI", 
funded by the European Commission under the NextGeneration EU programme. 
This research benefitted from the support of the FMJH Program Gaspard Monge for optimization and operations research and their interactions with data science.

The authors thank Lucia Celli and Luis Maia for pointing out an issue in the proof of Proposition 3.7.
%========================================================================
%========================================================================
%========================================================================

%\bibliographystyle{amsplain}
%\bibliography{biblio}
% \addbibresource{biblio.bib}
%\printbibliography

\bibliographystyle{plainnat}
\bibliography{biblio}


    %%% NEURIPS CHECKLIST
\begin{comment}
    
    \newpage
\section*{NeurIPS Paper Checklist}

\begin{enumerate}

\item {\bf Claims}
    \item[] Question: Do the main claims made in the abstract and introduction accurately reflect the paper's contributions and scope?
    \item[] Answer: \answerYes{}.
    \item[] Justification: Our main result Theorem \ref{thm:main} is a precise statement of the compressed version included in the introduction.
            Our abstract accurately reflects this result's content and scope.
    \item[] Guidelines:
    \begin{itemize}
        \item The answer NA means that the abstract and introduction do not include the claims made in the paper.
        \item The abstract and/or introduction should clearly state the claims made, including the contributions made in the paper and important assumptions and limitations. A No or NA answer to this question will not be perceived well by the reviewers. 
        \item The claims made should match theoretical and experimental results, and reflect how much the results can be expected to generalize to other settings. 
        \item It is fine to include aspirational goals as motivation as long as it is clear that these goals are not attained by the paper. 
    \end{itemize}

\item {\bf Limitations}
    \item[] Question: Does the paper discuss the limitations of the work performed by the authors?
    \item[] Answer: \answerYes{}.
    \item[] Justification: Limitations of our work are presented together with possible future work directions in Section \ref{sec:discussion}.
    \item[] Guidelines:
    \begin{itemize}
        \item The answer NA means that the paper has no limitation while the answer No means that the paper has limitations, but those are not discussed in the paper. 
        \item The authors are encouraged to create a separate "Limitations" section in their paper.
        \item The paper should point out any strong assumptions and how robust the results are to violations of these assumptions (e.g., independence assumptions, noiseless settings, model well-specification, asymptotic approximations only holding locally). The authors should reflect on how these assumptions might be violated in practice and what the implications would be.
        \item The authors should reflect on the scope of the claims made, e.g., if the approach was only tested on a few datasets or with a few runs. In general, empirical results often depend on implicit assumptions, which should be articulated.
        \item The authors should reflect on the factors that influence the performance of the approach. For example, a facial recognition algorithm may perform poorly when image resolution is low or images are taken in low lighting. Or a speech-to-text system might not be used reliably to provide closed captions for online lectures because it fails to handle technical jargon.
        \item The authors should discuss the computational efficiency of the proposed algorithms and how they scale with dataset size.
        \item If applicable, the authors should discuss possible limitations of their approach to address problems of privacy and fairness.
        \item While the authors might fear that complete honesty about limitations might be used by reviewers as grounds for rejection, a worse outcome might be that reviewers discover limitations that aren't acknowledged in the paper. The authors should use their best judgment and recognize that individual actions in favor of transparency play an important role in developing norms that preserve the integrity of the community. Reviewers will be specifically instructed to not penalize honesty concerning limitations.
    \end{itemize}

\item {\bf Theory assumptions and proofs}
    \item[] Question: For each theoretical result, does the paper provide the full set of assumptions and a complete (and correct) proof?
    \item[] Answer: \answerYes{}.
    \item[] Justification: Every new result stated as Theorem, Proposition or Lemma and is accompanied by its proof in the supplementary material.
            An additional theorem environment was created to present and cross-reference the assumptions in our results.
    \item[] Guidelines:
    \begin{itemize}
        \item The answer NA means that the paper does not include theoretical results. 
        \item All the theorems, formulae, and proofs in the paper should be numbered and cross-referenced.
        \item All assumptions should be clearly stated or referenced in the statement of any theorems.
        \item The proofs can either appear in the main paper or the supplemental material, but if they appear in the supplemental material, the authors are encouraged to provide a short proof sketch to provide intuition. 
        \item Inversely, any informal proof provided in the core of the paper should be complemented by formal proofs provided in appendix or supplemental material.
        \item Theorems and Lemmas that the proof relies upon should be properly referenced. 
    \end{itemize}

    \item {\bf Experimental result reproducibility}
    \item[] Question: Does the paper fully disclose all the information needed to reproduce the main experimental results of the paper to the extent that it affects the main claims and/or conclusions of the paper (regardless of whether the code and data are provided or not)?
    \item[] Answer: \answerYes{}.
    \item[] Justification: In Section \ref{sec:experiments} we provide a detailed exposition of how our numerical experiments were programmed and the Python libraries used on them.
    \item[] Guidelines:
    \begin{itemize}
        \item The answer NA means that the paper does not include experiments.
        \item If the paper includes experiments, a No answer to this question will not be perceived well by the reviewers: Making the paper reproducible is important, regardless of whether the code and data are provided or not.
        \item If the contribution is a dataset and/or model, the authors should describe the steps taken to make their results reproducible or verifiable. 
        \item Depending on the contribution, reproducibility can be accomplished in various ways. For example, if the contribution is a novel architecture, describing the architecture fully might suffice, or if the contribution is a specific model and empirical evaluation, it may be necessary to either make it possible for others to replicate the model with the same dataset, or provide access to the model. In general. releasing code and data is often one good way to accomplish this, but reproducibility can also be provided via detailed instructions for how to replicate the results, access to a hosted model (e.g., in the case of a large language model), releasing of a model checkpoint, or other means that are appropriate to the research performed.
        \item While NeurIPS does not require releasing code, the conference does require all submissions to provide some reasonable avenue for reproducibility, which may depend on the nature of the contribution. For example
        \begin{enumerate}
            \item If the contribution is primarily a new algorithm, the paper should make it clear how to reproduce that algorithm.
            \item If the contribution is primarily a new model architecture, the paper should describe the architecture clearly and fully.
            \item If the contribution is a new model (e.g., a large language model), then there should either be a way to access this model for reproducing the results or a way to reproduce the model (e.g., with an open-source dataset or instructions for how to construct the dataset).
            \item We recognize that reproducibility may be tricky in some cases, in which case authors are welcome to describe the particular way they provide for reproducibility. In the case of closed-source models, it may be that access to the model is limited in some way (e.g., to registered users), but it should be possible for other researchers to have some path to reproducing or verifying the results.
        \end{enumerate}
    \end{itemize}


\item {\bf Open access to data and code}
    \item[] Question: Does the paper provide open access to the data and code, with sufficient instructions to faithfully reproduce the main experimental results, as described in supplemental material?
    \item[] Answer: \answerYes{}.
    \item[] Justification: The code used to produce the experiments in Figure \ref{fig:experiment} is available at \href{https://anonymous.4open.science/r/quantitative_gaussian_trainedNN-F7A1/}{an anonymized repository}.
    \item[] Guidelines:
    \begin{itemize}
        \item The answer NA means that paper does not include experiments requiring code.
        \item Please see the NeurIPS code and data submission guidelines (\url{https://nips.cc/public/guides/CodeSubmissionPolicy}) for more details.
        \item While we encourage the release of code and data, we understand that this might not be possible, so “No” is an acceptable answer. Papers cannot be rejected simply for not including code, unless this is central to the contribution (e.g., for a new open-source benchmark).
        \item The instructions should contain the exact command and environment needed to run to reproduce the results. See the NeurIPS code and data submission guidelines (\url{https://nips.cc/public/guides/CodeSubmissionPolicy}) for more details.
        \item The authors should provide instructions on data access and preparation, including how to access the raw data, preprocessed data, intermediate data, and generated data, etc.
        \item The authors should provide scripts to reproduce all experimental results for the new proposed method and baselines. If only a subset of experiments are reproducible, they should state which ones are omitted from the script and why.
        \item At submission time, to preserve anonymity, the authors should release anonymized versions (if applicable).
        \item Providing as much information as possible in supplemental material (appended to the paper) is recommended, but including URLs to data and code is permitted.
    \end{itemize}


\item {\bf Experimental setting/details}
    \item[] Question: Does the paper specify all the training and test details (e.g., data splits, hyperparameters, how they were chosen, type of optimizer, etc.) necessary to understand the results?
    \item[] Answer: \answerYes{}.
    \item[] Justification: A detailed description of all the hyperparameters is included in Section \ref{sec:discussion}.
    \item[] Guidelines:
    \begin{itemize}
        \item The answer NA means that the paper does not include experiments.
        \item The experimental setting should be presented in the core of the paper to a level of detail that is necessary to appreciate the results and make sense of them.
        \item The full details can be provided either with the code, in appendix, or as supplemental material.
    \end{itemize}

\item {\bf Experiment statistical significance}
    \item[] Question: Does the paper report error bars suitably and correctly defined or other appropriate information about the statistical significance of the experiments?
    \item[] Answer: \answerYes{}.
    \item[] Justification: Our plot of the Gaussian process $G_t$ in \ref{fig:experiment} includes a 95\% confidence interval.
                     As for the plot involving the distance $\mathcal W_2(G_t,f_t)$, a subsection discussing the approximation error on the computation of the empirical Wasserstein distance has been included.
    \item[] Guidelines:
    \begin{itemize}
        \item The answer NA means that the paper does not include experiments.
        \item The authors should answer "Yes" if the results are accompanied by error bars, confidence intervals, or statistical significance tests, at least for the experiments that support the main claims of the paper.
        \item The factors of variability that the error bars are capturing should be clearly stated (for example, train/test split, initialization, random drawing of some parameter, or overall run with given experimental conditions).
        \item The method for calculating the error bars should be explained (closed form formula, call to a library function, bootstrap, etc.)
        \item The assumptions made should be given (e.g., Normally distributed errors).
        \item It should be clear whether the error bar is the standard deviation or the standard error of the mean.
        \item It is OK to report 1-sigma error bars, but one should state it. The authors should preferably report a 2-sigma error bar than state that they have a 96\% CI, if the hypothesis of Normality of errors is not verified.
        \item For asymmetric distributions, the authors should be careful not to show in tables or figures symmetric error bars that would yield results that are out of range (e.g. negative error rates).
        \item If error bars are reported in tables or plots, The authors should explain in the text how they were calculated and reference the corresponding figures or tables in the text.
    \end{itemize}

\item {\bf Experiments compute resources}
    \item[] Question: For each experiment, does the paper provide sufficient information on the computer resources (type of compute workers, memory, time of execution) needed to reproduce the experiments?
    \item[] Answer: \answerNo{}
    \item[] Justification: Our experiments were run on the author's personal computer and do not require any particular hardware specifications.
    \item[] Guidelines:
    \begin{itemize}
        \item The answer NA means that the paper does not include experiments.
        \item The paper should indicate the type of compute workers CPU or GPU, internal cluster, or cloud provider, including relevant memory and storage.
        \item The paper should provide the amount of compute required for each of the individual experimental runs as well as estimate the total compute. 
        \item The paper should disclose whether the full research project required more compute than the experiments reported in the paper (e.g., preliminary or failed experiments that didn't make it into the paper). 
    \end{itemize}
    
\item {\bf Code of ethics}
    \item[] Question: Does the research conducted in the paper conform, in every respect, with the NeurIPS Code of Ethics \url{https://neurips.cc/public/EthicsGuidelines}?
    \item[] Answer: \answerYes{}.
    \item[] Justification: Our research follows NeurIPS code of ethics.
    \item[] Guidelines:
    \begin{itemize}
        \item The answer NA means that the authors have not reviewed the NeurIPS Code of Ethics.
        \item If the authors answer No, they should explain the special circumstances that require a deviation from the Code of Ethics.
        \item The authors should make sure to preserve anonymity (e.g., if there is a special consideration due to laws or regulations in their jurisdiction).
    \end{itemize}


\item {\bf Broader impacts}
    \item[] Question: Does the paper discuss both potential positive societal impacts and negative societal impacts of the work performed?
    \item[] Answer: \answerNA{}.
    \item[] Justification: The innovations our paper introduces are of a theoretical nature.
    \item[] Guidelines:
    \begin{itemize}
        \item The answer NA means that there is no societal impact of the work performed.
        \item If the authors answer NA or No, they should explain why their work has no societal impact or why the paper does not address societal impact.
        \item Examples of negative societal impacts include potential malicious or unintended uses (e.g., disinformation, generating fake profiles, surveillance), fairness considerations (e.g., deployment of technologies that could make decisions that unfairly impact specific groups), privacy considerations, and security considerations.
        \item The conference expects that many papers will be foundational research and not tied to particular applications, let alone deployments. However, if there is a direct path to any negative applications, the authors should point it out. For example, it is legitimate to point out that an improvement in the quality of generative models could be used to generate deepfakes for disinformation. On the other hand, it is not needed to point out that a generic algorithm for optimizing neural networks could enable people to train models that generate Deepfakes faster.
        \item The authors should consider possible harms that could arise when the technology is being used as intended and functioning correctly, harms that could arise when the technology is being used as intended but gives incorrect results, and harms following from (intentional or unintentional) misuse of the technology.
        \item If there are negative societal impacts, the authors could also discuss possible mitigation strategies (e.g., gated release of models, providing defenses in addition to attacks, mechanisms for monitoring misuse, mechanisms to monitor how a system learns from feedback over time, improving the efficiency and accessibility of ML).
    \end{itemize}
    
\item {\bf Safeguards}
    \item[] Question: Does the paper describe safeguards that have been put in place for responsible release of data or models that have a high risk for misuse (e.g., pretrained language models, image generators, or scraped datasets)?
    \item[] Answer: \answerNA{}.
    \item[] Justification: No sensitive data is used in our article. The data used in our experiments in Section \ref{sec:experiments} is artificially generated by our code and consists of points in $\R$.
    \item[] Guidelines:
    \begin{itemize}
        \item The answer NA means that the paper poses no such risks.
        \item Released models that have a high risk for misuse or dual-use should be released with necessary safeguards to allow for controlled use of the model, for example by requiring that users adhere to usage guidelines or restrictions to access the model or implementing safety filters. 
        \item Datasets that have been scraped from the Internet could pose safety risks. The authors should describe how they avoided releasing unsafe images.
        \item We recognize that providing effective safeguards is challenging, and many papers do not require this, but we encourage authors to take this into account and make a best faith effort.
    \end{itemize}

\item {\bf Licenses for existing assets}
    \item[] Question: Are the creators or original owners of assets (e.g., code, data, models), used in the paper, properly credited and are the license and terms of use explicitly mentioned and properly respected?
    \item[] Answer: \answerYes{}.
    \item[] Justification: The python libraries used in our experiments are adequately referenced in Section \ref{sec:experiments}.
    \item[] Guidelines:
    \begin{itemize}
        \item The answer NA means that the paper does not use existing assets.
        \item The authors should cite the original paper that produced the code package or dataset.
        \item The authors should state which version of the asset is used and, if possible, include a URL.
        \item The name of the license (e.g., CC-BY 4.0) should be included for each asset.
        \item For scraped data from a particular source (e.g., website), the copyright and terms of service of that source should be provided.
        \item If assets are released, the license, copyright information, and terms of use in the package should be provided. For popular datasets, \url{paperswithcode.com/datasets} has curated licenses for some datasets. Their licensing guide can help determine the license of a dataset.
        \item For existing datasets that are re-packaged, both the original license and the license of the derived asset (if it has changed) should be provided.
        \item If this information is not available online, the authors are encouraged to reach out to the asset's creators.
    \end{itemize}

\item {\bf New assets}
    \item[] Question: Are new assets introduced in the paper well documented and is the documentation provided alongside the assets?
    \item[] Answer: \answerNA{} 
    \item[] Justification: No new assets are introduced in this paper.
    \item[] Guidelines:
    \begin{itemize}
        \item The answer NA means that the paper does not release new assets.
        \item Researchers should communicate the details of the dataset/code/model as part of their submissions via structured templates. This includes details about training, license, limitations, etc. 
        \item The paper should discuss whether and how consent was obtained from people whose asset is used.
        \item At submission time, remember to anonymize your assets (if applicable). You can either create an anonymized URL or include an anonymized zip file.
    \end{itemize}

\item {\bf Crowdsourcing and research with human subjects}
    \item[] Question: For crowdsourcing experiments and research with human subjects, does the paper include the full text of instructions given to participants and screenshots, if applicable, as well as details about compensation (if any)? 
    \item[] Answer: \answerNA{}.
    \item[] Justification: Our paper does not involve human subjects in any way.
    \item[] Guidelines:
    \begin{itemize}
        \item The answer NA means that the paper does not involve crowdsourcing nor research with human subjects.
        \item Including this information in the supplemental material is fine, but if the main contribution of the paper involves human subjects, then as much detail as possible should be included in the main paper. 
        \item According to the NeurIPS Code of Ethics, workers involved in data collection, curation, or other labor should be paid at least the minimum wage in the country of the data collector. 
    \end{itemize}

\item {\bf Institutional review board (IRB) approvals or equivalent for research with human subjects}
    \item[] Question: Does the paper describe potential risks incurred by study participants, whether such risks were disclosed to the subjects, and whether Institutional Review Board (IRB) approvals (or an equivalent approval/review based on the requirements of your country or institution) were obtained?
    \item[] Answer: \answerNA{}.
    \item[] Justification: Our paper does not involve human subjects in any way.
    \item[] Guidelines:
    \begin{itemize}
        \item The answer NA means that the paper does not involve crowdsourcing nor research with human subjects.
        \item Depending on the country in which research is conducted, IRB approval (or equivalent) may be required for any human subjects research. If you obtained IRB approval, you should clearly state this in the paper. 
        \item We recognize that the procedures for this may vary significantly between institutions and locations, and we expect authors to adhere to the NeurIPS Code of Ethics and the guidelines for their institution. 
        \item For initial submissions, do not include any information that would break anonymity (if applicable), such as the institution conducting the review.
    \end{itemize}

\item {\bf Declaration of LLM usage}
    \item[] Question: Does the paper describe the usage of LLMs if it is an important, original, or non-standard component of the core methods in this research? Note that if the LLM is used only for writing, editing, or formatting purposes and does not impact the core methodology, scientific rigorousness, or originality of the research, declaration is not required.
    %this research? 
    \item[] Answer: \answerNA{}.
    \item[] Justification: Our results do not regard LLMs.
    \item[] Guidelines:
    \begin{itemize}
        \item The answer NA means that the core method development in this research does not involve LLMs as any important, original, or non-standard components.
        \item Please refer to our LLM policy (\url{https://neurips.cc/Conferences/2025/LLM}) for what should or should not be described.
    \end{itemize}

\end{enumerate}

%----------------------------------------------------------
%----------------------------------------------------------
%----------------------------------------------------------

\end{comment}

\newpage

\appendix

\section{Gradient flow of the feature function}
\label{ap:eq1}

In this section we provide a closed analytical solution to \eqref{eq1} when $f = f^\lin$.
%$f$ is linear on its parameters.

Fix $x \in \R^{n_0}$ a test point.
For the sake of clearness, we will use the following notation for each $t\ge 0$: $\overline y_t=f^\lin(x;\overline \theta_t)$, $f^\lin_t=f_t(\X;\overline \theta_t)$, $\ktrain = k_0(\X,\X)$ and $\ktest = k_0(x,\X)$.
Note that $k_t=k_0$ for each $t$ when $f=f^\lin$.
% is linear on its parameters.

We begin by stating a lemma that shows that our definition of $I_t$ is consistent and commutes with the wide limit:
\begin{lem}\label{lem_It_properties}
    For any real symmetric matrix $B \in \R^{n \times n}$ we have $I_t(B)B=BI_t(B) = \mathbbm{1}_n - e^{-Bt}$; and for real symmetric matrix sequence $(B_n)_{n\in\N}$ with $B_n\to B$ we have $\lim_{n\to \infty}I_t(B_n)=I_t(B)$.
\end{lem}
The proof of this result follows from properties of the matrix exponential and is left to Supplementary Material \ref{ap:proofs}.

\begin{comment}
    The next result is a corollary of well-known Gronwall-like inequalities (see, for example, \cite{teschl} or \cite{Hartman}) for the dependence of the flow \eqref{eq1} with respect to its parameters.

    \begin{prop} \label{hartman-teschl}
    Denote $f_t(x_0,k), f_t(x_0',k')$ the solutions for \eqref{eq1} at time $t>0$ for different initial values $x_0,x_0' \in \R^{n_0}$ and choices of the matrix $k,k' \in \R^{n \times n}$. 
    Let $\mathcal L$ be the Lipschitz constant of $f_t$ with respect to the variables $k,x_0$.
    Then
    \[
    \vert f_t(x_0,k)-f_t(x_0',k')\vert \le \left(\norm{k-k'}{op} + \norm{x_0-x_0'}{}\right)e^{\mathcal Lt}.
    \]
\end{prop}

Supplementary Material \ref{ap:hartman} is devoted to the details of this computation.
\end{comment}

Consider the system of ODEs in \eqref{eq1} given by:
\begin{align}
    \frac{\partial}{\partial t} f^\lin_t &= - \ktrain(f^\lin _t-y),\\
    \frac{\partial}{\partial t} {\overline y_t} &= -\ktest(f^\lin_t-y).
\end{align}

Recall that, in general, the solution to the initial value problem $f'(t)=A(t)f(t),\; f(0)=f_0$ can be written as $f(t) = \exp(\int_0^tA(s)ds)f_0$,
where $f(t)\colon \R \to \R^m, \; A(t) \colon \R\to \R^{m \times m}$ are integrable functions, and $t > 0$. 
Therefore in our case, by letting $u_t = f^\lin_t-y$:
\begin{equation}
\label{eq1:soltrain_a}
    u_t = \exp\left(-\int_0^t\ktrain ds\right)u_0 =\exp(-\ktrain t)u_0.
\end{equation}

Moreover, by letting $v_t = \overline y_t-y$ and substituting the solution for $u_t$, we obtain the following expression for $v_t$:
\[
\frac{\partial}{\partial t} v_t = -\ktest\exp(-\ktrain t)(f_0-y), \quad v_0 = y_0,
\]
which, by integrating and by using Definition \ref{def:It}, becomes
\begin{align}
    v_t - v_0 &= \overline y_t - \overline y_0 \\
    &= - \ktest \int_0^t \exp\left(-\ktrain s\right)ds (f_0 -y)\\
    \label{eq1:soltest_a} &= -\ktest I_t(\ktrain)(f_0-y).
\end{align}

Note that formulae \eqref{eq1:soltrain_a} and \eqref{eq1:soltest_a} agree with the formulae found by the authors of \cite{Lee2020}.

When $\ktrain$ is not degenerate, taking the limit when $t$ tends to infinity, we get a prediction for the output of the linearized network at the end of the training:
\[
    f^\lin_\infty(x) = \lim_{t \to \infty}\overline y_t = \overline y_0-\ktest \ktrain^{-1}(f_0-y).
\]

%Note that, when $\ktrain$ is degenerate, there exists $i$, $1\le i\le n$ such that $I_t(\ktrain)_{ii} = t$.
%In this case, the limit in $t$ exists only when $t$ grows slower than \textcolor{red}{$??$}.

%We will call $f_\infty(x)$ the \emph{NTK solution} for $x \in \R^{n_0}$.

\subsection{Proof of the characterization of $G_t$}
\label{ap:sub:proofGt}

    Here we prove the formulae in \eqref{characterization}.

    Define $B_t = -\ktest I_t(\ktrain)$ and $C_t = \ktest I_t(\ktrain)$, so that $\overline y_t-\overline y_0 = B_tf_0 + C_ty$.
    Also recall that $k_t=k_0$ for all $t\ge 0$ since $f$ is linear on $\theta$.
    Note that $f_0$ and $\overline y_0$ are centered Gaussian processes and hence $\E[\overline y_0+B_tf_0] = 0$.
    Therefore, taking the expected value of the wide limit yields:
    \begin{align}
        \E[\lim_{n_1 \to \infty}\overline y_t] &= \E[\lim_{n_1 \to \infty}C_t y] = k_\infty(x,\X)I_t(k_\infty)y,
    \end{align}
    where $k_\infty = k_\infty(\X,\X)$.
    This limit is well defined thanks to Lemma \ref{lem_It_properties}.
    
    Now let $x'\in \R^{n_0}$ and put $y'_t=f_t(x')$ and $B'_t = -k_{x'\X}I_t(\ktrain)$ for each $t\ge 0$.
    Then,
    \begin{align}
        \textup{Cov}(\lim_{n_1 \to \infty}\overline y_t,\lim_{n_1 \to \infty}y_t')&= \E[\lim_{n_1 \to \infty}(\overline y_t - \E[\overline y_t])(y'_t - \E[y'_t]) ] \\
        &= \E[\lim_{n_1 \to \infty}(\overline y_0 + B_tf_0)(y'_0 + B'_tf_0) ] \\
        &= \E[\lim_{n_1 \to \infty}\overline y_0y'_0 ] + \E[\lim_{n_1  \to \infty} \overline y_0f_0 B'_t]  \\
        &\quad +\E[\lim_{n_1  \to \infty} y'_0f_0B_t  ]+ \E[\lim_{n_1 \to \infty}f_0^2B_tB'_t] \\
        &=\mathcal K(x,x') - \mathcal K(x,\X)I_t(k_\infty)k_\infty(\X,x')\\
        &\quad- k_\infty(x,\X)I_t(k_\infty)\mathcal K(\X,x')\\
        &\quad + k_\infty(x,\X)I_t(k_\infty)\mathcal K(\X,\X)I_t(k_\infty)k_\infty(\X,x').
    \end{align}
    Again, Lemma \ref{lem_It_properties} ensures the limit exists.

\begin{comment}
    \section{Lipschitz-continuity and smoothness of the gradient flow solutions}
\label{ap:hartman}

In this Supplementary Material we obtain a bound for the flow of Equation \eqref{eq1} at positive time $t$.
Recall the following Gronwall-type bound for dependence on the initial value (Theorem 2.8 in \cite{teschl}):

\begin{thm}[Dependence on initial conditions]\label{thm:lipschitz}
    Let $U \subset \R$ be an open set and let $f,g \in C(U,\R^n)$ be locally Lipschitz on the second argument, uniformly with respect to the first.
    Let $x(t), y(t)$ be the respective solutions of the initial value problems
    \[
    \begin{cases}
        \frac{\partial}{\partial t} x = f(t,x), \\
        x(t_0) = x_0,
    \end{cases}
    \qquad
    \begin{cases}
        \frac{\partial}{\partial t} y = g(t,y), \\
        y(t_0) = y_0.
    \end{cases}
    \]

    Then
    \[
    \norm{x(t)-y(t)}{} \le \norm{x_0 - y_0}{} e^{L \vert t - t_0 \vert} + \frac{M}{L}\left(e^{L \vert t - t_0 \vert}-1\right),
    \]
    where
    \[
    L = \sup_{(t,x) \ne (t.y) \in U} \frac{\norm{f(t,x) - f(t,y)}{}}{\vert x-y \vert}, \qquad M = \sup_{(t,x)\in U} \norm{f(t,x) - g(t,x)}{}. 
    \]

\end{thm}

To obtain an analogous bound for the dependence on the operator $k$, we follow the procedure described on \cite{Hartman}, Preliminaries of Chapter 5. 
The $n$-dimensional initial value problem depending on a continuous parameter $z \in \R^d$:
\begin{equation}
    \begin{cases}
        y' = g(t,y,z) \\
        y(t_0) = y_0,
    \end{cases}
\end{equation}

can be expressed as an initial value problem for the $(n+d)$-dimensional vector $(y,z)$
\begin{equation}
    \begin{cases}
        \begin{pmatrix}
            y' \\ z'
        \end{pmatrix}
        = \begin{pmatrix}
            g(t,y,z) \\ 0
        \end{pmatrix} \\
        \begin{pmatrix}
            y(t_0) \\ z(t_0)
        \end{pmatrix} = 
        \begin{pmatrix}
            y_0 \\ z_0
        \end{pmatrix},
    \end{cases}
\end{equation}
in which no extra parameters occur.

Hence Theorem \ref{thm:lipschitz} can be applied to study the Lipschitz continuity on the operator $k$ in our setting.
In the notation of Theorem \ref{thm:lipschitz}, $f$ and $g$ coincide (hence $M = 0$) with the right-hand side of the matrix equation \eqref{eq1}, which has bounded partial derivatives with respect to $t,y,\ktrain$ and $\ktest$ on any compact subset of $\R^{m+n}$.
Therefore they are Lipschitz on those variables.
If $L$ is the Lipschitz constant (with respect to $k$) of $f$; and $f_{k_1}, f_{k_2}$ denote, respectively, the solutions to $\eqref{eq1}$ with kernel $k = k_1$ and $k=k_2$, 
$k_1,k_2 \colon \R^{n_0}\times \R^{n_0} \to \R$ then Theorem \ref{thm:lipschitz} reads:

\begin{equation}
\label{eq:lipschitz}
    \norm{f_{k_1}(x) - f_{k_2}(x)}{} \le \norm{k_1-k_2}{op}e^{Lt}.
\end{equation}

Using this bound, by triangle inequality the following corollary follows:

\begin{cor}
    Denote $f_t(x_0,k), f_t(x_0',k')$ the solutions for \eqref{eq1} at time $t>0$ for different initial values $x_0,x_0' \in \R^{n_0}$ and choices of the kernel $k,k' \colon \R^{n_0}\times \R^{n_0} \to \R$. 
    Then
    \[
    \vert f_t(x_0,k)-f_t(x_0',k') \vert \le \left(\norm{k-k'}{op} + \norm{x_0-x_0'}{}\right)e^{Lt}.
    \]
\end{cor}

\end{comment}

\section{Auxiliary and related results}
\label{ap:auxiliary}

In this Supplementary Material we state intermediate results in the proof of our main theorem and recall some useful results.
Throughout this section we will use the following notation for each $t\ge 0$: $y_t=f(x;\theta_t)$, $f_t=f(\X;\theta_t)$, $k_t = k_t(\X,\X)$ and $k_\infty = k_\infty(\X,\X)$.
All the proofs are deferred to Supplementary Material \ref{ap:proofs}.

In the next lemma we collect some well-known properties of the $p$-Wasserstein distance:

\begin{lem}\label{thm:wass_properties}
    Let $p \in [1, \infty[$ and let $X,Y$ be random variables with values in $\R^n$ and $Z$ be a random variable with values in $\R^m$. 
    Let $\mathbb P_\xi$ denote the law of the random variable $\xi$ for each $\xi \in \{X,Y,Z\}$. Then
    \begin{enumerate}
        \item If $X, Y$ are defined on the same probability space, then $\mathcal W_p(X,Y) \le \E[\norm{X - Y}{}^p]^{\frac{1}{p}}$.
        \item If $Z$ is independent from $X$ and $Y$ then $\mathcal W_p(X+Z,Y+Z) \le \mathcal W_p(X,Y)$.
        \item Convexity of $\mathcal W_p^p$: $\mathcal W_p^p(X,Y) \le \int_{\R^m} \mathcal W_p^p(\mathbb P_{X \vert Z = z},\mathbb P_Y)d\mathbb P_Z(z)$. 
         \item Let $\lambda \in \R^m$ be a constant vector and consider the joint random variables $\tilde X = (X,Z), \tilde Y = (Y, \lambda)$. Then 
        \[\mathcal W_p^p(\tilde X, \tilde Y) \le \mathcal W_p^p(X, Y) + \mathcal W_p^p(Z, \lambda) = \mathcal W_p^p(X, Y) + \left(\int_{\R^m} \norm{z-\lambda}{}^p d\mathbb P(z)\right)^{\frac{1}{p}},\]
        %where $d$ denotes the euclidean distance in $\R^m$. 
        \item Let $V$ be a random variable with values in $\R^m$ and consider the joint random variables $\tilde X = (X,Z), \tilde Y = (Y, V)$. Then, for $p \ge 2$, 
        \[\mathcal W_p^p(\tilde X, \tilde Y) \le 2^{\frac{p}{2}-1}\left(\mathcal W_{2p}^{2p}(X, Y) + \mathcal W_{2p}^{2p}(Z, V)\right).\]

        Moreover, for $p=1$,
        \[\mathcal W_1(\tilde X, \tilde Y) \le \mathcal W_1(X, Y) + \mathcal W_1(Z, V).\]
    \end{enumerate}
\end{lem}

The following result provides explicit formulae for the components of the Jacobian of $f$ and the NTK:

\begin{lem}[Gradients $f$ and explicit formulae for $\tilde k$ and $k$]
    \label{lemma1}
    The following hold for each $x,x'\in \R^{n_0}$:
    \begin{align}
        \nabla_{\theta^{(0)}}f(x,\theta) &= \frac{1}{\sqrt{n_1n_0}}x^\top\Phi'(\frac{1}{\sqrt{n_0}}x\theta^{(0)})\theta^{(1)} \in \R^{n_0}, \\
        \nabla_{\theta^{(1)}}f(x,\theta) &= \frac{1}{\sqrt{n_1}}\Phi(\frac{1}{\sqrt{n_0}}x\theta^{(0)}) \in \R^{n_1}.
    \end{align}
    Moreover, $\tilde k(x,x')$ is a diagonal $n_1 \times n_1$ matrix with
    \begin{align}
        \tilde k_{ii}(x,x') = \frac{1}{n_0}\sum_{u=1}^{n_0}x_ux'_u,
    \end{align}
    and $k(x,x')$ is a real function given by:
    \begin{align}
        k(x,x') &=\frac{1}{n_1n_0}\sum_{u=1}^{n_0} x_ux'_u \sum_{v=1}^{n_1}\Phi'(h_v(x))\Phi'(h_v(x'))(\theta^{(1)}_v)^2 + \frac{1}{n_1}\sum_{v=1}^{n_1}\Phi(h_v(x))\Phi(h_v(x')).
    \end{align}
\end{lem}

\subsection{Results at initialization}

Throughout this subsection, we assume $t=0$ and omit the subindex $t$ unless needed.
Our results \ref{prop:main_linear} and \ref{thm:main} aim to generalize Theorem 4.1 in \cite{ trevisan2023wide} in different directions.
%which generalizes the preceding Theorem 3.1 in \cite{BasteriTrevisan2023}. 
For reference, we reproduce the main result from \cite{trevisan2023wide} here in a simplified version:

\begin{thm}[Trevisan]
    \label{thm:trevisanoriginal}
    %For a fixed $1 \le \ell \le L$, let $f_0(\X) = f^{(\ell)}_0(\X)$ and $G_0(\X) = G^{(\ell)}_0(\X)$.
    %Then, for each $p \in \N$ there exists a constant $c_{p,\ell}$ not depending on the network widths $n_1, \ldots, n_L$ such that
    Then, for each $p \in \N$ there exists a constant $c_p$ not depending on the network width $n_1$ such that:
    \begin{align}
        \mathcal W_p(f_0(\X), G_0(\X)) \le c_p \frac{1}{\sqrt{n_1}}.
    \end{align}
    %Moreover, if $\mathcal K$ is not degenerate, it holds
    %\begin{align}
    %    \mathcal W_p(f_0(\X), G_0(\X)) \le c_p \frac{1}{n_1},
    %\end{align}
    %for $1 \le \ell \le L$. 
    Furthermore, if $\p \colon \R \to \R$ is a Lipschitz function, then
    \begin{align}
        %\mathcal W_p\left(\frac{1}{n_i}\sum_{i=1}^{n_i}\p((f_0(\X))_i)^{\otimes 2}, \p(G_0(\X))^{\otimes 2}\right) \le c_p \frac{(\textup{Lip}h + h(0))^2}{\sqrt{n_1}}.
        \mathcal W_p\left(\p((f_0(\X)))^{\otimes 2}, \p(G_0(\X))^{\otimes 2}\right) \le c_p \frac{(\textup{Lip}\p + \p(0))^2}{\sqrt{n_1}}.
    \end{align}
\end{thm}

Theorem \ref{thm:trevisanoriginal} provides a quantitative bound for the Gaussian approximation of the neural network $f_t$.
This can be upgraded to a bound for the joint distribution of the empirical kernel and the output of the neural network.

From now to the end of this subsection assume $\Phi$ and $\Phi'$ are bounded and $x,x'\in \R^{n_0}$ are fixed.
Also, we adopt the notation introduced in Supplementary Material \ref{ap:eq1}.
We state a helpful estimation:

\begin{prop}[$L^p$ bound for the kernel difference]
    \label{prop:kernelp}
    Fix $x,x'$ in $\R^{n_0}$ and let $\tilde k_{11} = \tilde k_{11}(x,x')$, $k=k(x,x')$, $\mathcal K=\mathcal K(x,x')$ and $k_\infty = k_\infty(x,x')$.
    There exists a constant $C > 0$ independent of $n_1$ such that:
    \begin{align}
        %\E[\vert \tilde k_{11}-\mathcal K\vert^p] &= 0, \\
        \tilde k_{11}-\tilde{\mathcal K} &= 0, \\
        \E[\vert k-k_\infty\vert^p] &\le \frac{C}{n_1^\frac{p}{2}}.
    \end{align}
\end{prop}

\begin{rem}
    Note that since $k_\infty$ is deterministic, we have $\mathcal W_p^p(k,k_\infty) = \E[\norm{k -  k_\infty}{}^p]$.
    The constant $C$ in Proposition \ref{prop:kernelp} depends on the constants produced by applications of Theorem \ref{thm:trevisanoriginal}.
\end{rem}

The proof of Proposition \ref{prop:kernelp} goes by triangle inequality combined with Theorem \ref{thm:trevisanoriginal}, 
and by exploiting the independence between the entries of $\theta^{(1)}$ and those of $\theta^{(0)}$.
An auxiliary result to prove \ref{prop:kernelp} is the following:

\begin{prop}[$L^p$ bounds for the empirical kernel]
    \label{prop:kernelbound}
    Let $\tilde k_{ij} = \tilde k_{ij}(x,x')$, for each $1 \le i,j \le n_1$, and $k=k(x,x')$.
    Then, the following inequalities hold:
    \begin{align}
        \E[\vert k\vert] &= \norm{\Phi}{\infty}^2,\\
        \E[\vert k\vert^p] &\le 2^{p-1}(2p-1)!!\norm{\Phi'}{\infty}^{2p}\vert \tilde k_{11}\vert^p + 2^{p-1}\norm{\Phi}{\infty}^{2p}.
    \end{align}
    \end{prop}

Now we are ready to show the following result:

\begin{prop}[Joint distribution Basteri-Trevisan]
\label{prop:trevisan}
There exist a positive constant $C$ such that:
\[
\mathcal W_p \left(\left(k_0, f_0\right), \left(k_\infty, G_0\right)\right)  \le  \frac{C}{\sqrt{n_1}}.
\]
with $C$ not depending on the width $n_1$.
\end{prop}

The proof is by using Dudley's lemma (also referred to as the gluing lemma from optimal transport) as outlined in \cite{villani_transport}.
This lemma is used to decompose the Wasserstein distance into two terms. 
The summand regarding the network and its Gaussian approximation is bounded with Theorem \ref{thm:trevisanoriginal}, 
and the other summand is bounded by Proposition \ref{prop:kernelp}.

Lastly, the following lemma is used in the proof of Proposition \ref{prop:main_linear}.

\begin{lem}
The following inequalities hold:
    \label{lem:4moment}
    \begin{align}
        \E[\norm{f_0-y}{}]&\le \sqrt{n}\norm{\Phi}{\infty} + \norm{y}{},\\
        \E[\norm{f_0-y}{}^4]&\le 32n^2\norm{\Phi}{\infty}^4 + 8\norm{y}{}^4.
    \end{align}
\end{lem}

\subsection{Approximation by linearization at training time $t>0$}

In this subsection we state two results paramount to prove Proposition \ref{quasi_main}, along with all their auxiliary lemmas.
The following theorem resembles Theorem 5.4 in \cite{Bartlett_Montanari_Rakhlin_2021}, or Theorem 2.2 in \cite{bachchizat},
but we would like to remark that those result differ from ours since the first applies to neural networks in which the training is restricted to $\theta^{(0)}$ while keeping $\theta^{(1)}$ frozen,
and in both of them the loss function used for training is different than the one considered in our work.
The proof is similar to the one in \cite{Bartlett_Montanari_Rakhlin_2021} and is carried in full detail in Supplementary Material \ref{ap:proofs}.

In this result we prove quenched estimations for the dynamics of the parameters, 
the linearization error both in the parameters and in the network valued on test points, 
and the convergence of the network to the labels over the training set.

\begin{ass}
    \label{assumption_montanari}
    The smallest eigenvalue of $k_0$ is bounded from below by:
    \[
        4 L(\X) \norm{f_0-y}{} < \lammin(k_0),
    \]
    where $L(\X)$ the Lipschitz constant of $\nabla_\theta f_0$ seen as a function of $\theta$.
\end{ass}
\begin{thm}
    \label{thm:montanari}
    %Consider $L(\X)$ the Lipschitz constant of $\nabla_\theta f_0$ seen as a function of $\theta$.
    Let $\lammin = \lammin(k_0), \sigmin = \sigmin(k_0)$ and $\sigmax = \sigmax(k_0)$
    and let Assumption \ref{assumption_montanari} hold.
    %Assume
    %\begin{align}
    %    \label{assumption_montanari}
    %    4 \L(\X) \norm{f_0-y}{} < \lammin.
    %\end{align} 
    Then the following hold for $t>0$ and for any test point $x\in \R^{n_0}$:
    \begin{align}
        \norm{\theta_t-\theta_0}{} &\le \frac{2}{\sigmin}\norm{f_0-y}{}, \\
        \norm{\theta_t-\overline \theta_t}{} &\le \frac{(8+20\sigmax^2)L(\X)}{\sigmin^3} \norm{f_0-y}{}^2,\\
        \norm{f_t-y}{}^2 &\le \norm{f_0-y}{}^2 \exp\left(-\frac{\lammin}{2} t\right), \\
        \norm{f(x;\theta_t) - f^\lin(x;\overline \theta_t)}{} &\le \frac{4 L(x)}{\sigmin^2}\norm{y-f_0}{}^2 \\
        &\quad + \frac{(8+20\sigmax^2)L(\X)}{\sigmin^3} \norm{f_0-y}{}^2 \norm{\nabla_\theta f_0(x)}{}.
    \end{align}
\end{thm}

Proposition \ref{thm:montanari} relies on a strong assumption involving the Lipschitz constant of the Jacobian of the network, 
the norm of the network at initialization and the positive-definiteness of the limiting kernel.
The following rougher estimations depend on $t$, but they do not require Assumption \ref{assumption_montanari} and will be key to prove our main theorem.

\begin{thm}
    \label{thm:stime_polinomiali}
    Assume that $\Phi$ and $\Phi'$ are bounded.
    Let $L(\X)$ be the Lipschitz constant of $\nabla_\theta f_0$, seen as a function of $\theta$,
    and let $\s(\theta_0) = \norm{f_0-y}{}$.
    Then, for each $t > 0$ there exist positive constants $A_1, \ldots, A_5,B_1, \ldots, B_9,C_1, \ldots, C_8$ and $C_9$ not depending on $n_1, n_0$ nor $t$ such that, $\mathbb P$-almost everywhere:
    \begin{align}
        \norm{y_t- f_t}{}^2 &\le \frac{A_0}{n_0}\norm{\theta^{(0)}_0\theta^{(1)}_0}{}^2 +\frac{A_1t^2}{n_0n_1}\norm{\theta^{(0)}_0}{}^2\s(\theta_0)^2 + \frac{A_2\norm{\theta^{(1)}_0}{}^2t^4}{n_1^2n_0}\s(\theta_0)^4 \\
        &\quad + \frac{A_3t^6}{n_1^2n_0}\s(\theta_0)^6 + \frac{A_4t^2}{n_1n_0}\norm{\theta^{(1)}_0}{}^4\s(\theta_0)^2 + \frac{A_5t^4\norm{\theta^{(1)}_0}{}^2}{n_1^2n_0}\s(\theta_0)^4,
    \end{align}
    \begin{align}
        \norm{f^\lin-\overline y_t}{}^2 &\le \frac{B_0}{n_1n_0}\norm{\theta_0^{(0)}\theta_0^{(1)} }{}^2 + \frac{B_1}{n_1^2n_0^2}\norm{\theta_0^{(0)}\theta_0^{(1)} }{}^2\s(\theta_0)^4t^4\\
        &\quad +\frac{B_2}{n_1^2n_0^2}\norm{\theta_0^{(0)}}{}^2\norm{\theta_0^{(1)} }{}^4\s(\theta_0)^2t^2+\frac{B_3}{n_1n_0^2}\norm{\theta_0^{(0)}\theta_0^{(1)} }{}^2\s(\theta_0)^2 t^2\\
        &\quad +\frac{B_4}{n_1^2n_0}\norm{\theta_0^{(1)} }{}^2\s(\theta_0)^4t^4+\frac{B_5}{n_1^2n_0}\norm{\theta_0^{(1)} }{}^4\s(\theta_0)^2t^2 \\
        &\quad +\frac{B_6}{n_1n_0}\norm{\theta_0^{(1)} }{}^2\s(\theta_0)^2t^2+\frac{B_7}{n_1^2n_0}\norm{\theta_0^{(0)} }{}^2\s(\theta_0)^4t^4 \\
        &\quad +\frac{B_8}{n_1^2n_0}\norm{\theta_0^{(0)} }{}^2\norm{\theta_0^{(1)} }{}^2\s(\theta_0)^2t^2+\frac{B_9}{n_1n_0}\norm{\theta_0^{(0)} }{}^2\s(\theta_0)^2t^2,
    \end{align}
    \begin{align}
        \norm{f_t- f^\lin}{}^2 &\le \frac{L(\X)^2}{n_1^2}\left(\frac{C_1 \s(\theta_0)^4\norm{\theta^{(1)}_0}{}^2t^2}{n_1n_0^2}+\frac{C_2 \s(\theta_0)^6t^8}{n_1^2n_0^2}+\frac{C_3 \s(\theta_0)^4t^6}{n_1n_0}\right. \\
    &\left.\quad +\frac{C_4 \norm{\theta^{(1)}_0}{}^4t^4}{n_0^2}+\frac{C_5 \s(\theta_0)^2\norm{\theta^{(1)}_0}{}^2t^6}{n_1n_0^2}+\frac{C_6 \norm{\theta^{(1)}_0}{}^2t^6}{n_0}\right.\\
    &\left.\quad +\frac{C_7 \s(\theta_0)^2\norm{\theta^{(1)}_0}{}^2t^4}{n_0}+\frac{C_8 \s(\theta_0)^4t^6}{n_1n_0}+C_9\s(\theta_0)^2t^4\right),
\end{align}
where $\theta^{(0)}\theta^{(1)} \in \R^{n_0}$ denotes the usual product of the matrices $\theta^{(0)}$ and $\theta^{(1)}$.
\end{thm}
\begin{rem}
    The dependence on time of the right-hand side of the above formulae comes from Lemma \ref{lem:gronwallbeta}.
    Indeed, by definition of the operator $I_t$, the sharpest upper bound for the matrices $I_t(k_t)$ and $I_t(k_0)$ when no lower bound for $\lammin(k_t)$ and $\lammin(k_0)$ is available is $\mathbbm 1_n t$.
\end{rem}
The proof of the first two inequalities in this theorem is by exploiting an expression for the gradients of the network, contained in Lemma \ref{lemma1};
together with an integral result describing the behaviour of the parameters at time $t$ with respect to the parameters at initialization.
This is the content of Lemma \ref{lem:gronwallbeta}.
The third inequality uses an integral argument together with the semipositive-definiteness of $k_t$ to redirect the problem to studying $\norm{k_t-k_0}{}$.
All the constants in Theorem \ref{thm:stime_polinomiali} are multiples of the norms and Lipschitz constants of $\Phi$ and $\Phi'$, and of the norms of $x$ and $\X$.

Now we state Lemma \ref{lem:gronwallbeta} along with some concentration inequalities and related auxiliary results.

    \begin{lem}[Inequalities for $\theta^{(i)}_t$]
        \label{lem:gronwallbeta}
        Fix $u$ and $v$ with $1 \le u \le n_0$ and $1 \le v \le n_1$ and put $\X_u = ((x_i)_u)_{i=1}^n\in \R^n$.
        Let $\lammin$ and $\lammin^0$ be the smallest eigenvalues of $k_t$ and $k_0$, respectively.
        Then the following inequalities hold:
        \begin{align}
            (\theta^{(1)}_v)_t &\le (\theta_v^{(1)})_0 + \frac{\norm{\Phi}{\infty}\norm{f_0-y}{}}{\sqrt{n_1}}I_t(\lammin),\\
            (\theta^{(0)}_{uv})_t &\le (\theta_{uv}^{(0)})_0 + \frac{\norm{\Phi}{\infty}\norm{\Phi'}{\infty}\norm{f_0-y}{}^2\norm{\X_u}{}}{2n_1\sqrt{n_0}}I_t(\lammin)^2\\
        &\quad + \frac{\norm{\Phi'}{\infty}\norm{f_0-y}{}\norm{\X_u}{}}{\sqrt{n_1n_0}}(\theta^{(1)}_v)_0{}I_t(\lammin),\\
        (\overline \theta^{(1)}_v)_t &\le (\theta_v^{(1)})_0 + \frac{\norm{\Phi}{\infty}\norm{f_0-y}{}}{\sqrt{n_1}}I_t(\lammin^0),\\
        (\overline \theta^{(0)}_{uv})_t &\le (\theta_{uv}^{(0)})_0 + \frac{\norm{\Phi}{\infty}\norm{\Phi'}{\infty}\norm{f_0-y}{}^2\norm{\X_u}{}}{2n_1\sqrt{n_0}}I_t(\lammin^0)^2 \\
    &\quad + \frac{\norm{\Phi'}{\infty}\norm{f_0-y}{}\norm{\X_u}{}}{\sqrt{n_1n_0}}(\theta^{(1)}_v)_0{}I_t(\lammin^0).
\end{align}
    \end{lem} 
    \begin{rem}
        Recall from the definition of $I_t$ that for $t>0$, $I_t(a) > 0$, even if $a = 0$.
        Moreover, $I_t(0) = t$ by definition.
        Hence $I_t(\lammin^0) \le t$ and $I_t(\lammin^0)^2 \le t^2$.

        %Moreover, $I_t(0) = t$ by definition and $I_t^2(0) = I_t(t) \le 1$ when $t>0$. 
        %In effect:
        %\begin{align}
        %    I_t(t) &= \frac{1-e^{-t^2}}{t} \le \frac{1}{t}.
        %\end{align}
        %On the other hand,
        %\begin{align}
        %    I_t(t)& \le t \sup_{0\le s \le t}e^{-st} \le t.
        %\end{align}
        %Multiplying the two estimations, $(I_t(t))^2 \le 1$, which implies $I_t(t) \le 1$.
        %Hence in the right-hand side of the above formulae, the dependence in time is at most linear when $t>0$.
    \end{rem}

Recall the well known concentration inequality for $\chi^2$-distributed random variables (see, for example, \cite{laurentmassart}).
    For each $\gamma > 0$:
    \begin{equation}
        \label{eqmassart_laurent}
        \mathbb P (\norm{\theta_0^{(1)}}{}^2 \ge 2\gamma + 2 \sqrt{ \gamma n_1 }+ n_1) \le \exp(-\gamma). 
    \end{equation}

    The following is a concentration inequality for the sup-norm of $\theta_0^{(1)}$; 
    and as a consequence we get an estimation of the norm and Lipschitz constant of the Jacobian of $f$ at initialization.
\begin{lem}
    \label{lem:supnorm}
    For any $\gamma > 0$:
    \begin{align}
        \norm{\theta_0^{(1)}}{\infty} 
        %&\le \sqrt{\gamma \log n_1} + \sqrt{2\log(2n_1)},\\
        &\le\sqrt{r\gamma \log n_1},
    \end{align}
    with probability bigger or equal than $1 - \frac{1}{n_1^{\frac{r\gamma}{2}-1}}$.
    %Moreover, if $\gamma \ge 1$ and $n_1 \ge 2$, then $\norm{\theta_0^{(1)}}{\infty} \le \sqrt{5 \gamma \log n_1}$ with the same probability.
\end{lem}

This concentration inequality is proven using the fact that $\norm{\theta_0^{(1)}}{\infty}$ is the supremum of $n_1$ Gaussian variables in absolute value.

\begin{lem}[Norm and Lipschitz constant of the Jacobian at $t=0$]
    \label{lem:jacobian}
    Fix $r\ge 1$.
    Then for each $x \in \R^{n_0}$,
    \begin{align}
        \norm{\nabla_\theta f_0(x)}{} &\le \frac{\norm{x}{}\norm{\Phi'}{\infty}\sqrt{\gamma}}{\sqrt{n_0}} + \frac{\norm{\Phi}{\infty}}{\sqrt{n_1}},\\
        \textup{Lip}\nabla_\theta f_0(x) &\le \frac{\norm{x}{}(\norm{\Phi'}{\infty}+\textup{Lip}\Phi)}{\sqrt{n_1n_0}} + \frac{\norm{x}{}^2\textup{Lip}\Phi'}{\sqrt{n_1}n_0}\sqrt{r\gamma \log n_1}.
    \end{align}
    with probability greater or equal than $1 - \frac{1}{n_1^{\frac{r\gamma}{2}-1}}-\exp(-\gamma n_1)$, where $f_0(x) \colon \R^N \to \R$ is understood as a function of $\theta$.
\end{lem}

Now we state a concentration inequality controlling the norm of the difference between the NTK and its limit.
\begin{lem}
    \label{lem:lammin}
    Let $k=k_0(\X,\X)$ and $k_\infty = k_\infty(\X,\X)$ and let $\gamma \in \N$.
    Put $\lammin = \lammin(k)$ and $\lammin^\infty = \lammin(k_\infty)$.
    Then, for each $p\in \N$,
    \begin{align}
        \norm{k-k_\infty}{} \le \frac{\gamma \lammin^\infty}{2},
    \end{align}
    with probability greater or equal than $1 - \left(\frac{2}{\gamma \lammin^\infty}\right)^p\frac{C}{n_1^\frac{p}{2}}$, 
    where $C$ is a positive constant not depending on $n_1$.
\end{lem}
\begin{rem}
    The previous lemma provides useful bounds for the smallest and largest eigenvalues of the empirical kernel at initialization, with arbitrarily high probability when the width diverges.
    Recall that the smallest eigenvalue of a matrix is a 1-Lipschitz function of the operator norm, which is bounded by the Frobenius norm:
    \begin{align}
        \vert \lammin(A) - \lammin(B) \vert \le \norm{A-B}{op} \le \norm{A-B}{}.
    \end{align}
    In particular, the previous Lemma implies, for $\gamma = 1$:
    \begin{align}
        \lammin &\ge \lammin^\infty - \norm{k-k_\infty}{} \ge \frac{\lammin^\infty}{2}.
    \end{align}
    Conversely, an upper bound for the largest eigenvalue of a matrix using the operator norm is given by:
    \begin{align}
        \lammax(A) &\le \lammax(B) + \norm{A-B}{op} \le \lammax(B) + \norm{A-B}{}.
    \end{align}
    Again, taking $\gamma=1$ in the previous lemma yields:
    \begin{align}
        \lammax &\le \lammax^\infty + \frac{\lammin^\infty}{2}.
    \end{align}
    Both inequalities hold with probability greater or equal than $1 - \left(\frac{2}{\lammin^\infty}\right)^p\frac{C}{n_1^\frac{p}{2}}$.
\end{rem}
%\begin{rem}
%    The constant $C$ in this Lemma comes from Proposition \ref{prop:kernelp} by taking the maximum of all the constants $c_1,d_1$ and $C$ in the right hand side of that result.
%    Note that there, the constants $c_1,d_1$ do not depend on $n_1$, while $C$ tends to zero as $n_1$ diverges.
%    Explicitly, $C \lesssim \norm{\Phi'}{\infty}^2 \frac{cp^{1+\frac{1}{p}}}{\log p} \nu_p^\frac{4}{p}\frac{1}{n_1^\frac{1}{p}}$, where $c>0$ is an absolute constant coming from Rosenthal's inequality and $\nu_p$ is the $p$-th moment of a standard Gaussian.
%    Explicit expressions for all the constants involved are available on the proof of Proposition \ref{prop:kernelp}.
%\end{rem}
 
\begin{comment} %Lo ho spostato nel main body
    Now we are ready to state the main result of this subsection:
\begin{thm}
    \label{quasi_main}
    Make the following assumption on $n_0$ and $n_1$:
    \[
    \frac{4\norm{\X}{}(\sqrt 5\norm{\Phi}{\infty}+\norm{y}{})}{\sqrt{n_1n_0}}
        \left(\norm{\Phi'}{\infty} + \textup{Lip}\Phi + \frac{\norm{\X}{}\textup{Lip}\Phi'\sqrt{r \log n_1}}{\sqrt{n_0}}\right) < \lammin^\infty,
    \]
    for some $r \ge 5$.
    Then, there exist positive constants $\alpha_1$ and $\alpha_2$ not depending on $n_0,n_1$ nor $t$ such that:
    \begin{align}
        \mathcal W_2^2(y_t,\overline y_t) \le \left(\frac{\alpha_1\log n_1}{n_1n_0}+\frac{\alpha_2 n_0n_1^{-\frac{3}{4}}}{1-n_1^{-\frac{3}{4}}}(1+t^2+t^4+t^6)\right).
    \end{align}
\end{thm}
\begin{rem}
    The hypothesis in this theorem is quite mild: notice that the left hand side tends to zero as $n_1$ grows.
    In \cite{carvalho2024}, the authors showed that the analytic NTK is positive definite when the points in the training set are in general position in $\R^{n_0}$.
    This implies that our hypothesis is satisfied for sufficiently big $n_0$ or $n_1$.
\end{rem}

The proof is as follows: we split $\R^N$ into a "good" event $S \subset \R^N$, in which the assumptions of Proposition \ref{thm:montanari} hold along some other concentration properties of the parameters,
and a "bad" event $S^C$ in which they do not.
As $n_1$ diverges, $\mathbb P(S)$ converges to 1.
Splitting the Wasserstein distance into integrals over these two sets correspond to the two summands in the right-hand side of the result.
The first integral is limited by Proposition \ref{thm:montanari} combined with Lemmas \ref{lem:jacobian} and \ref{lem:supnorm}.
In $S^C$ the strategy is to show that the measure of $S^C$ decreases faster than how the upper bounds in Theorem \ref{thm:stime_polinomiali} grow.
We manage to do so by partitioning $S^C$ into countable discs parametrized by $\gamma \in \N$, and summing over $\gamma$ while exploiting concentration inequalities that hold on each disc.

\end{comment}

\section{Proof of Theorems \ref{thm:main} and \ref{quasi_main}}
\label{ap:proof_1}

Here we prove Theorems \ref{thm:main} and \ref{quasi_main}, which share some auxiliary lemmas.
Throughout this and the remaining appendices we will use the following notation for each $t\ge 0$: $y_t=f_t(x)$, $f_t=f_t(\X)$, $\overline y_t=f^\lin(x;\overline \theta_t)$, $f^\lin_t=f^\lin(\X;\overline \theta_t)$
$k_t = k_t(\X,\X)$ and $k_\infty = k_\infty(\X,\X)$.
The gradient $\nabla$ and the expectation $\E$ will always be taken with respect to the parameters $\theta$, unless otherwise indicated.

Let us begin by Proposition \ref{quasi_main}:

\begin{proof}[Proof of Proposition \ref{quasi_main}]
    Fix $p, r \in \N$.
    Consider the following subset of $\R^N$:
    \[S = \{\theta \mid \frac{\norm{\theta}{}^2}{n_1} \le 5, \norm{\theta}{\infty} \le \sqrt{r\log n_1}, \norm{k-k_\infty}{}\le \frac{\lammin^\infty}{2}\}.\]
    By the inequality \eqref{eqmassart_laurent} and Lemmas \ref{lem:supnorm} and \ref{lem:lammin}, the probability of $S$ is bounded from below by 
    $1-\exp(-n_1)-\frac{1}{\sqrt{n_1^{r-2}}}-\frac{\hat c}{(\lammin^\infty \sqrt{n_1})^p}$, for a positive constant $\hat c$ not depending on $n_1, n_0$ nor $t$.
    Then,
    \begin{align}
        \mathcal W_2^2(y_t,\overline y_t) &\le \inf_{\theta} \E_\theta[\norm{y_t-\overline y_t}{}^2] \\
        \label{eqaux_wassfflin}& = \int_S \norm{y_t-\overline y_t}{}^2 d\theta + \int_{S^C}\norm{y_t-\overline y_t}{}^2 d\theta \\
        &\le \sup_{\theta \in S} \norm{y_t-\overline y_t}{}^2 + \left(\exp(-n_1)+\frac{1}{\sqrt{n_1^{r-2}}}+\frac{\hat c}{(\lammin^\infty \sqrt{n_1})^p}\right)\sup_{\theta \in S^C}\norm{y_t-\overline y_t}{}^2.
    \end{align}
    
    Let us begin by estimating the supremum over $S$.
    Note that by Lemma \ref{lem:lammin}, $\lammin\ge \frac{\lammin^\infty}{2} > 0$ in $S$.
    The hypothesis for Proposition \ref{thm:montanari} are satisfied when $n_1,n_0$ are large enough.
    In particular, thanks to Lemma \ref{lem:jacobian}, Assumption \ref{ass:r}
    \[
        \frac{4\norm{\X}{}(\sqrt 5\norm{\Phi}{\infty}+\norm{y}{})}{\sqrt{n_1n_0}}
        \left(\norm{\Phi'}{\infty} + \textup{Lip}\Phi + \frac{\norm{\X}{}\textup{Lip}\Phi'\sqrt{r \log n_1}}{\sqrt{n_0}}\right) < \lammin^\infty
    \]
    implies Assumption \ref{assumption_montanari} for any $\theta$ in $S$.
    Indeed, by Lemmas \ref{lemma1} and \ref{lem:jacobian}:
    \begin{align}
        4 L(\X) &\norm{f_0-y}{} \\
        &\le \frac{4\norm{\X}{}}{\sqrt{n_1n_0}}
        \left(\norm{\Phi'}{\infty} + \textup{Lip}\Phi + \frac{\norm{\X}{}\textup{Lip}\Phi'\sqrt{r \log n_1}}{\sqrt{n_0}}\right)(\sqrt 5\norm{\Phi}{\infty}+\norm{y}{}).
    \end{align}
    Moreover, Lemma \ref{lem:lammin} implies $\lammin \ge \frac{\lammin^\infty}{2}$ in $S$.
    These two inequalities together show that Assumption \ref{ass:r}, which holds for sufficiently big $n_1$, is a sufficient condition for Proposition \ref{thm:montanari} to hold.

    On the other hand, by Lemmas \ref{lem:lammin} and \ref{lem:jacobian}, the following inequalities are satisfied in $S$, for $Z \in \{x,\X\}$:
    \begin{align}
        L(Z) &\le \frac{\norm{Z}{}}{\sqrt{n_1n_0}}\left(\norm{\Phi'}{\infty} + \textup{Lip}\Phi + \frac{\norm{\X}{}\textup{Lip}\Phi'\sqrt{r \log n_1}}{\sqrt{n_0}}\right),\\
        \norm{\nabla_\theta f_0}{} &\le \frac{\norm{\X}{}\norm{\Phi'}{\infty}}{\sqrt{n_0}} + \frac{\norm{\Phi}{\infty}}{\sqrt{n_1}},\\
        \lammax &\le \lammax^\infty + \frac{\lammin^\infty}{2}.
    \end{align}
    By substituting the estimations of $L(x),L(\X),\nabla f(x;\theta_0), \lammax$ and $\lammin$ above in Proposition \ref{thm:montanari}, for each $\theta \in S$:
    \begin{align}
        \norm{y_t - \overline y_t}{}&\le\frac{8(\sqrt{5}\norm{\Phi}{\infty} + \norm{y}{})^2}{\lammin^\infty \sqrt{n_1n_0}}\left(\norm{x}{} + \frac{2 \norm{\X}{}}{\sqrt{\lammin^\infty}}(\sqrt 2 + 15 \lammax^\infty) \right. \\
        &\left.\quad \cdot\left(\norm{\Phi'}{\infty} + \textup{Lip}\Phi + \frac{\norm{\X}{}\textup{Lip}\Phi'\sqrt{r \log n_1}}{\sqrt{n_0}}\right)\left(\frac{\norm{\X}{}\norm{\Phi'}{\infty}}{\sqrt{n_0}} + \frac{\norm{\Phi}{\infty}}{\sqrt{n_1}}\right)\right) \\
        &\le c \sqrt{\frac{r\log n_1}{n_1n_0(\lammin^\infty)^3}},
    \end{align}
    where the constant $c$ is independent of $r,t,n_0$ and $n_1$ and can be determined explicitly:
    \[c = 384\norm{\X}{}(\sqrt{5}\norm{\Phi}{\infty} + \norm{y}{})^2\max\{\norm{x}{},\sqrt 2, 15\lammax^\infty,\norm{\Phi'}{\infty}, \textup{Lip}\Phi, \norm{\X}{}\textup{Lip}\Phi',\norm{\X}{}\norm{\Phi'}{\infty},\norm{\Phi}{\infty}\}.
    \]

    Hence,
    \begin{align}
        \int_S \norm{y_t-\overline y_t}{}^2 d\theta &\le \mathbb P(S)\sup_{\theta\in S} \norm{y_t-\overline y_t}{}^2 \\
        &\le \sup_{\theta\in S} \norm{y_t-\overline y_t}{}^2 \\
        &\le \frac{c^2r\log n_1}{(\lammin^\infty)^3n_1n_0}.
    \end{align}
    %with $c = 256\sqrt{2}(2+5 \lammax^\infty + \frac{5\lammin^\infty}{2}) (\sqrt 5 \norm{\Phi}{\infty} + \norm{y}{})^4 (\norm{x}{}^2 + \norm{\X}{}^2)\max\{2 \textup{Lip} \Phi', \norm{\Phi'}{\infty}\}\frac{(\textup{Lip}\Phi)^2\norm{\X}{}\norm{\Phi}{\infty}}{n_0}$.
    %Note that the upper bound $\lammax \le \lammax^\infty + \frac{\lammin^\infty}{2}$ holds for $\theta\in S$, thanks to Lemma \ref{lem:lammin}.
    Put $\alpha_1 = c^2$.

    Now it only remains to estimate the second summand in \eqref{eqaux_wassfflin}.
    For each $\gamma \in \N$ let $\hat \gamma = 2\gamma + \sqrt{2 \gamma} + 1$ and define the subsets:
    \begin{align}
        \Omega_\gamma = \{\theta \mid \frac{\norm{\theta^{(1)}}{}^2}{n_1} >\hat \gamma, \norm{\theta^{(1)}}{\infty} > \sqrt{r\gamma \log n_1}\}.
        %\Omega_\gamma = \{\theta \mid \norm{k-k_\infty}{}> \frac{\gamma \lammin^\infty}{2}, \norm{\theta^{(1)}}{\infty} > \sqrt{\gamma \log n_1} + \sqrt{2 \log(2n_1)}\}.
    \end{align}
    Also, define the subset
    \begin{align}
        \Omega_* = \{\theta \mid\norm{k-k_\infty}{}> \frac{\lammin^\infty}{2}\},
    \end{align}
    and let $\Omega = \Omega_* \setminus \bigcup_{\gamma \in \N}\Omega_\gamma$.

    Intuitively, $\Omega_*$ and $\bigcup_{\gamma \in \N}\Omega_\gamma$ are the events in which the lower bound for the smallest eigenvalue of the empirical kernel
    and the upper bound of the Frobenius and sup-norm of the parameters at initialization, respectively, do not hold.
    Notice that $\R^N = S \sqcup \Omega \sqcup \bigcup_{\gamma \in \N}\Omega_\gamma$.
    We will use this partition of $\R^N$ to finish the proof.

    By Lemma \ref{lem:supnorm} and by \eqref{eqmassart_laurent} we have 
    $\mathbb P (\Omega_\gamma) \le \exp(-\gamma n_1) + \frac{1}{n_1^{\frac{r\gamma}{2}-1}}$ and $\mathbb P(\Omega) \le \mathbb P(\Omega_*) \le \frac{\hat c}{(\lammin^\infty\sqrt{n_1})^p}$.
    Moreover, the family $(\Omega_\gamma)_{\gamma \in \N}$ is a descending filtration of $S^C \setminus \Omega$.
    Let $D_\gamma = \Omega_\gamma \setminus \Omega_{\gamma + 1}$, for each $\gamma \in \N$.
    This allows us to write:
    \begin{align}
        \int_{S^C} \norm{y_t-\overline y_t}{}^2 d\theta &\le \int_{\Omega}\norm{y_t-\overline y_t}{}^2 d\theta + \sum_{\gamma \in \N}\int_{D_\gamma} \norm{y_t-\overline y_t}{}^2 d\theta \\
        &\le \left(\mathbb P(\Omega) + \sum_{\gamma \in \N}\mathbb P(D_\gamma)\right) \sup_{\theta \in D_\gamma}\norm{y_t-\overline y_t}{}^2\\
        &\le 3\mathbb P(\Omega)\sup_{\theta \in \Omega}\left(\norm{y_t-f_t}{}^2 + \norm{f_t-f_t^\lin}{}^2+\norm{f^\lin_t-\overline y_t}{}^2\right)  \\
        &\quad + 3\sum_{\gamma \in \N}\mathbb P(D_\gamma)\sup_{\theta \in D_\gamma}\left(\norm{y_t-f_t}{}^2 + \norm{f_t-f_t^\lin}{}^2+\norm{f^\lin_t-\overline y_t}{}^2\right) \\
        \label{eqaux_sum_eventocattivo}&\le 3\sum_{\gamma \in \N}\exp(-\gamma n_1) \sup_{\theta \in D_\gamma}\norm{y_t-f_t}{}^2 + 3\sum_{\gamma \in \N}\exp(-\gamma n_1) \sup_{\theta \in D_\gamma}\norm{f_t-f^\lin_t}{}^2 \\
        &\quad  + 3\sum_{\gamma \in \N}\exp(-\gamma n_1) \sup_{\theta \in D_\gamma}\norm{f^\lin_t-\overline y_t}{}^2 + 3\sum_{\gamma \in \N}\frac{1}{n_1^{\frac{r\gamma}{2}-1}} \sup_{\theta \in D_\gamma}\norm{y_t-f_t}{}^2 \\
        &\quad + 3\sum_{\gamma \in \N}\frac{1}{n_1^{\frac{r\gamma}{2}-1}} \sup_{\theta \in D_\gamma}\norm{f_t-f^\lin_t}{}^2 + 3\sum_{\gamma \in \N}\frac{1}{n_1^{\frac{r\gamma}{2}-1}} \sup_{\theta \in D_\gamma}\norm{f^\lin_t-\overline y_t}{}^2 \\
        &\quad +3\frac{\hat c}{(\lammin^\infty \sqrt{n_1})^p}\sup_{\theta \in \Omega}\norm{y_t-f_t}{}^2 + 3\frac{\hat c}{(\lammin^\infty \sqrt{n_1})^p}\sup_{\theta \in \Omega}\norm{f_t-f^\lin_t}{}^2 \\
        &\quad +3\frac{\hat c}{(\lammin^\infty \sqrt{n_1})^p}\sup_{\theta \in \Omega}\norm{f^\lin_t-\overline y_t}{}^2 
    \end{align}

    The 6 series in the previous expression are convergent.
    We compute an upper bound for each of the 6 series in \eqref{eqaux_sum_eventocattivo} with the aid of Theorem \ref{thm:stime_polinomiali}, Lemma \ref{lem:supnorm} and the inequality \eqref{eqmassart_laurent}.
    Let $A = \max\{A_i\}, B = \max\{B_i\}$ and $C = \max\{C_i\}$ the maximums among the constants in Theorem \ref{thm:stime_polinomiali}.
    \begin{enumerate}
        \item Let us estimate the first series.
        Observe that we can bound $\widehat{\gamma+1} \le 7\gamma$.
        Moreover, since $A(\sqrt{\gamma+1} +\norm{y}{}) \le 2A\sqrt{\gamma}$ for $\gamma$ large enough, up to adding to the constant $A$ a multiple of $\norm{y}{}$ we can bound:
        \begin{align}
            \sum_{\gamma \in \N}&\exp(-\gamma n_1) \sup_{\theta \in D_\gamma}\norm{y_t-f_t}{}^2 \\
            &\le A\sum_{\gamma \in \N}\exp(-\gamma n_1)\left((\widehat{\gamma+1})^2n_1^2n_0 +\widehat{\gamma+1}(\sqrt{\gamma+1}+\norm{y}{})^2 t^2 \right.\\
            &\left.\quad + \frac{\widehat{\gamma+1}(\sqrt{\gamma+1}+\norm{y}{})^4t^4}{n_1n_0} + \frac{(\sqrt{\gamma+1}+\norm{y}{})^6t^6}{n_1^2n_0} \right.\\
            &\left.\quad + \frac{(\widehat{\gamma+1})^2(\sqrt{\gamma+1}+\norm{y}{})^2n_1t^2}{n_0} + \frac{\widehat{\gamma+1}(\sqrt{\gamma+1}+\norm{y}{})^4t^4}{n_1n_0}\right)\\
            &\le 7A\sum_{\gamma \in \N}\exp(-\gamma n_1)\left(7\gamma^3 n_1^2n_0 +4\gamma^2 t^2 \right.\\
            &\left.\quad + \frac{16\gamma^3t^4}{n_1n_0} + \frac{64\gamma^3t^6}{7n_1^2n_0} + \frac{28\gamma^3n_1t^2}{n_0} + \frac{16\gamma^3t^4}{n_1n_0}\right).
        \end{align}
        Since the negative exponential function decraeases faster than any polynomial, for each $p\in \N$ there exists a positive constant $N_p$ such that $x^p\le \exp(-\frac{x}{2})$ for each $x\ge N_p$.
        Therefore, up to a multiplicative constant not depending on $n_1,n_0,t$ nor $\gamma$:
        \begin{align}
            \sum_{\gamma \in \N}\exp(-\gamma n_1) \sup_{\theta \in D_\gamma}\norm{y_t-f_t}{}^2 &\le 7\cdot 64 \cdot 6An_0(1+t^6)\sum_{\gamma \in \N}\exp(-\frac{\gamma n_1}{2}) \\
            &\le \frac{7\cdot 64 \cdot 6An_0(1+t^6) e^{-\frac{n_1}{2}}}{1-e^{-\frac{n_1}{2}}}.
        \end{align}

        \item Now we estimate the second series.
        Using the bounds from the first series combined with Theorem \ref{thm:stime_polinomiali}:
        \begin{align}
            \sum_{\gamma \in \N}&\exp(-\gamma n_1) \sup_{\theta \in D_\gamma}\norm{f_t^\lin - \overline y_t}{}^2 \\
            &\le 7Bn_1\sum_{\gamma \in \N}\exp(-\gamma n_1) \left(7\gamma^2 n_1+ \frac{112\gamma^4}{n_1}+\frac{196\gamma^4n_1t^2}{n_0}+\frac{28\gamma^3n_1t^2}{n_0}\right.\\
            &\left.\quad +\frac{16\gamma^3t^4}{n_1n_0}+\frac{28\gamma^3 t^2}{n_0} +\frac{4\gamma^2 t^2}{n_0}+16\gamma^3t^4n_0 +28\gamma^3t^2+4\gamma^2t^2\right) \\
            &\le 7\cdot 196\cdot 9Bn_1(1+t^4)\sum_{\gamma \in \N}\exp(-\frac{\gamma n_1}{2}) \\
            &=\frac{7\cdot 196\cdot 9Bn_1(1+t^4)\exp(-\frac{n_1}{2})}{1-\exp(-\frac{n_1}{2})}.
        \end{align}

        \item Now we estimate the third series in the same fashion as we did with the preceding series.
        For an upper bound of $L(\X)$, recall Lemma \ref{lem:jacobian}, which reads 
        \[ L(\X) \le \frac{d}{\sqrt{n_1n_0}}\left(1 + \frac{\sqrt{r(\gamma+1) \log n_1}}{\sqrt{n_0}}\right),
        \] 
        for $\theta \in D_\gamma$ and $d$ a constant not depending on $n_1, n_0, t$ nor $\gamma$.
        For $n_1$ large enough, we can suppose $1 + \sqrt{\frac{r(\gamma+1) \log n_1}{n_0}} \le 2\sqrt{\frac{r\gamma \log n_1}{n_0}}$.
        Then,
        \begin{align}
            \sum_{\gamma \in \N}&\exp(-\gamma n_1) \sup_{\theta \in D_\gamma}\norm{f_t-f^\lin_t}{}^2 \\
            &\le \frac{4Cdr \log n_1}{n_1^3n_0^2}\sum_{\gamma \in \N}\gamma \exp(-\gamma n_1)\left(\frac{112 \gamma^3 t^2}{n_0^2}+\frac{64 \gamma^3t^8}{n_1^2n_0^2}+\frac{16 \gamma^2t^6}{n_1n_0}\right. \\
            &\left.\quad +\frac{49\gamma^2n_1^2t^4}{n_0^2}+\frac{28\gamma^2t^6}{n_0^2}+\frac{7 \gamma n_1t^6}{n_0}+\frac{28\gamma^2n_1t^4}{n_0}+\frac{16 \gamma^2 t^6}{n_1n_0}+4\gamma t^4\right)\\
            &\le \frac{4\cdot 112 \cdot 9Cdr \log n_1(1+t^8)}{n_1^3n_0^2}\sum_{\gamma \in \N}\exp(-\frac{\gamma n_1}{2})\\
            &=\frac{4\cdot 112 \cdot 9Cdr (1+t^8) e^{-\frac{n_1}{2}}}{n_1^2n_0^2(1- e^{-\frac{n_1}{2}})}.
        \end{align}

        \item Now we estimate the fourth series. 
        Following the reasoning from the first series:
        \begin{align}
            \sum_{\gamma \in \N}\frac{1}{n_1^{\frac{r\gamma}{2}-1}} \sup_{\theta \in D_\gamma}\norm{y_t-f_t}{}^2 &\le 7A\sum_{\gamma \in \N}\frac{1}{n_1^{\frac{r\gamma}{2}-1}}\left(7\gamma^3 n_1^2n_0 +4\gamma^2 t^2 + \frac{16\gamma^3t^4}{n_1n_0} \right.\\
            &\left.\quad + \frac{64\gamma^3t^6}{7n_1^2n_0} + \frac{28\gamma^3n_1t^2}{n_0} + \frac{16\gamma^3t^4}{n_1n_0}\right).
        \end{align}
        Recall that we can choose $r$ large enough so that all the summands have $n_1$ on the denominator.
        In particular, it is enough to choose $r\ge 5$ in this case.
        Moreover, by reasoning like in the proof of the first series,
        up to a multiplicative constant the terms of the form $\gamma^pn_1^{-\frac{r\gamma}{2}}$ are bounded from above by $n_1^{-\frac{r\gamma}{4}}$.
        Therefore, up to a multiplicative constant not depending on $n_1$,$n_0$,$t$ nor $\gamma$:
        \begin{align}
            \sum_{\gamma \in \N}\exp(-\gamma n_1) \sup_{\theta \in D_\gamma}\norm{y_t-f_t}{}^2 &\le 7\cdot 64 \cdot 6An_0(1+t^6)\sum_{\gamma \in \N}\frac{1}{n_1^\frac{r\gamma}{4}}\\
            &=\frac{7\cdot 64 \cdot 6An_0(1+t^6)n_1^{-\frac{r}{4}}}{1-n_1^{-\frac{r}{4}}}.
        \end{align}

        \item Now we estimate the fifth series.
        Using the bounds from the first series combined with Theorem \ref{thm:stime_polinomiali}, up to a multiplicative constant:
        \begin{align}
            \sum_{\gamma \in \N}&\frac{1}{n_1^{\frac{r\gamma}{2}-1}} \sup_{\theta \in D_\gamma}\norm{f_t^\lin - \overline y_t}{}^2 \\
            &\le 7Bn_1\sum_{\gamma \in \N}\frac{1}{n_1^\frac{r\gamma}{2}} \left(7\gamma^2 n_1+ \frac{112\gamma^4t^4}{n_1}+\frac{196\gamma^4n_1t^2}{n_0}+\frac{28\gamma^3n_1t^2}{n_0}\right.\\
            &\left.\quad +\frac{16\gamma^3t^4}{n_1n_0}+\frac{28\gamma^3 t^2}{n_0} +\frac{4\gamma^2 t^2}{n_0}+16\gamma^3t^4n_0 +28\gamma^3t^2+4\gamma^2t^2\right).
        \end{align}
    
    As we did for the fourth series, we can choose $r$ large enough so that the previous series is convergent.
    $r\ge 5$ is sufficient.
    Up to a multiplicative constant:
    \begin{align}
        \sum_{\gamma \in \N}\frac{1}{n_1^{\frac{r\gamma}{2}-1}} \sup_{\theta \in D_\gamma}\norm{y_t-f_t}{}^2 &\le 196\cdot7\cdot 9 Bn_0(1+t^4)\sum_{\gamma \in \N}\frac{1}{n_1^{\frac{r\gamma}{4}}}\\
        &\le \frac{196\cdot7\cdot 9 Bn_0(1+t^4) n_1^{-\frac{r}{4}}}{1-n_1^{-\frac{r}{4}}}.
    \end{align}

    \item Now we estimate the sixth and last series.
    Following the reasoning in the third and fourth series, up to a multiplicative constant:
    \begin{align}
        \sum_{\gamma \in \N}&\frac{1}{n_1^{\frac{r\gamma}{2}-1}} \sup_{\theta \in D_\gamma}\norm{f_t-f^\lin_t}{}^2 \\
        &\le \frac{4Cdr \log n_1}{n_1^2n_0^2}\sum_{\gamma \in \N}\frac{\gamma}{n_1^{\frac{r\gamma}{2}}} \left(\frac{112 \gamma^3 t^2}{n_0^2}+\frac{64 \gamma^3t^8}{n_1^2n_0^2}+\frac{16 \gamma^2t^7}{n_1n_0}\right. \\
            &\left.\quad +\frac{49\gamma^2n_1^2t^4}{n_0^2}+\frac{28\gamma^2t^6}{n_0^2}+\frac{7 \gamma n_1t^6}{n_0}+\frac{28\gamma^2n_1t^4}{n_0}+\frac{16 \gamma^2 t^6}{n_1n_0}+4\gamma t^4\right).
        \end{align}
    In this case, any $r \ge 3$ makes the series converge.
    \begin{align}
        \sum_{\gamma \in \N}&\frac{1}{n_1^{\frac{r\gamma}{2}-1}} \sup_{\theta \in D_\gamma}\norm{f_t-f^\lin_t}{}^2 \\
        &\le \frac{4\cdot 112 \cdot 9 Cdr(1+t^8)}{n_1n_0^2}\sum_{\gamma \in \N}\frac{\gamma}{n_1^{\frac{r\gamma}{4}}} \\
        &=\frac{4\cdot 112 \cdot 9 Cdr(1+t^8)n_1^{-\frac{r}{4}}}{n_1n_0^2(1-n_1^{-\frac{r}{4}})}.
    \end{align}
\end{enumerate}
    
    Now we finish the proof by estimating the last 3 terms in \eqref{eqaux_sum_eventocattivo}.
    Recall that, by definition of $\Omega$, the bounds for the Frobenius and the sup-norms of the parameters at initialization used in the estimation of $\int_S \norm{y_t-\overline y_t}{}d\theta$ hold also for $\theta \in \Omega$;
    and recall that $\mathbb P(\Omega) \le \frac{\hat c}{(\lammin^\infty \sqrt{n_1})^p}$.
    Let $R_1 = \frac{\hat c}{(\lammin^\infty)^p}$.
    Then it suffices to choose $p$ large enough to counterattack the biggest exponent of $n_1$ in each of the bounds given by Theorem \ref{thm:stime_polinomiali}.
    In particular, as seen while choosing $r$ when computing the six series, it is enough to take $p = r\ge 5$.
    Then, up to multiplicative constants,
    \begin{align}
        \frac{R_1}{n_1^\frac{p}{2}}\sup_{\theta \in \Omega}\norm{y_t-f_t}{}^2 &\le AR_1n_0(1+t^6) \frac{1}{\sqrt{n_1}},\\
        \frac{R_1}{n_1^\frac{p}{2}}\sup_{\theta \in \Omega}\norm{f_t}{}^2 &\le R_1Bn_0(1+t^4) \frac{1}{\sqrt{n_1}},\\
        \frac{R_1}{n_1^\frac{p}{2}}\sup_{\theta \in \Omega}\norm{y_t-f_t}{}^2 &\le CdrR_1(1+t^8) \frac{\log n_1}{n_1\sqrt{n_1}n_0^2}\\
        &\le CdrR_1(1+t^8) \frac{1}{\sqrt{n_1}n_0^2}.
    \end{align}

    Grouping the estimations for the nine summands in \eqref{eqaux_sum_eventocattivo} and taking 
    %$\alpha_2 = (\lammin^\infty)^rR_1\max\{A,B,Cd\}$,
    $\alpha_2 = R_1\max\{A,B,Cd\}$,
     up to a multiplicative constant,
    \begin{align}
        \int_{S^C} \norm{y_t-\overline y_t}{}^2d\theta &\le \frac{\alpha_2 r n_0}{(\lammin^\infty)^rn_1^{\frac{r}{4}}}(1+t^8).
    \end{align}
    %where in the last step the terms $t^2$ and $t^4$ were absorbed by the constant $\alpha_2$.
This concludes the proof.
\end{proof}

Then, our main result, Theorem \ref{thm:main}, is a direct application of Propositions \ref{quasi_main} and \ref{prop:main_linear}:

\begin{proof}[Proof of Theorem \ref{thm:main}]
    By triangle inequality and the elementary inequality $(a+b)^2\le 2a^2+2b^2$ for $a,b\ge 0$ decompose:
    \begin{align}
        \mathcal W_2^2(y_t,G_t) &\le 2\mathcal W_2^2(y_t,\overline y_t) + 2\mathcal W_2^2(\overline y_t,G_t).
    \end{align}
    Then the thesis follows estimating the first summand with Proposition \ref{quasi_main} and the second one with Proposition \ref{prop:main_linear}.
    For large enough $n_1$ we can take the constants $a_1$ and $a_2$ to be a multiple of $\alpha_1$ and $\alpha_2$ in Proposition \ref{quasi_main}; 
    since the right-hand side in the statement in that result decreases as $\frac{\log n_1}{n_1} + \frac{1}{n_1^\frac{r}{4}}$, 
    which is strictly slower than the right-hand side of the statement in Proposition \ref{prop:main_linear} for any $n_1 \ge 2$.
\end{proof}

\section{Proof of Proposition \ref{prop:main_linear}}
\label{ap:proof_2}

\begin{proof}[Proof of Proposition \ref{prop:main_linear}]

  %CORRECTED PROOF
%\color{teal}
Fix $x \in \R^{n_0}$ a test point.
During the proof we will use the notations $f_t = f_t^\lin(\X)$, $y_t = f_t^\lin(x)$, $G_t^x = G_t(x)$ and $G_t^\X = G_t(\X)$.
Also, we will denote by $\lammin$ and $\lammin^\infty$ the minimum eigenvalues of $\ktrain$ and $k_\infty$, respectively.
First we show the result for the training set.
With the aid of Equation \eqref{eq1}, we derive a closed ODE for $\vert f_t-G_t\vert^2$.
By Cauchy-Schwarz's and Young's inequalities:
    \begin{align}
        \frac{1}{2}\frac{\partial}{\partial t} \norm{f_t-G_t^\X}{}^2 &= \langle\ktrain (y-f_t)-k_\infty(y-G_t^\X), (f_t-G_t^\X)\rangle\\
        &= (\ktrain - k_\infty) (y-f_t) (f_t-G_t^\X)- k_\infty \norm{f_t-G_t^\X}{}^2 \\
        &\le \norm{\ktrain - k_\infty}{}\norm{y-f_t}{}\norm{f_t-G_t^\X}{}- \lammin^\infty \norm{f_t-G_t^\X}{}^2 \\
        &\le \frac{1}{2\eps}\norm{\ktrain - k_\infty}{}^2\norm{y-f_t}{}^2 + \frac{\eps}{2} \norm{f_t-G_t^\X}{}^2 - \lammin^\infty \norm{f_t-G_t^\X}{}^2.
    \end{align}

Note that, by gradient flow equation and since $\ktrain$ is semipositive definite, the following rough bound holds for each $t \ge 0$:
\[\norm{f_t-y}{}\le e^{-\lammin t}\norm{f_0-y}{} \le \norm{f_0-y}{}\]
Choosing $\eps = \lammin^\infty$ and putting $b=\frac{1}{\lammin^\infty}\norm{\ktrain - k_\infty}{}^2\norm{y-f_0}{}^2$ together yield:
\begin{equation}
    \label{eqaux_gronyoung1}
    \frac{\partial}{\partial t} \norm{f_t-G_t^\X}{}^2 \le b - \lammin^\infty \norm{f_t-G_t^\X}{}^2.
\end{equation}

Grönwall's inequality applied to \eqref{eqaux_gronyoung1} implies:
\begin{align}
    \norm{f_t-G_t}{}^2 &\le e^{-\lammin^\infty t}\left(\norm{f_0-G_0}{}^2 + b\int_0^te^{\lammin^\infty s}ds\right).
    %&= e^{-\lammin^\infty t}\left(\frac{\norm{\ktrain - k_\infty}{}^2\norm{f_0-y}{}^2 t}{\lammin^\infty} + \norm{f_0-G_0}{}^2 \right).
\end{align}

Recall the definition of $2$-Wasserstein distance.
By taking the expected value of the previous equation and the infimum on all the couplings between $f_t$ and $G_t$ we can bound by Hölder's inequality after exchanging the integrals with Fubini-Tonelli theorem:
\begin{align}
    \mathcal W^2_2(f_t,G_t) &\le \E[\norm{f_t-G_t}{}^2] \\ 
    &\le e^{-\lammin^\infty t} \left(\frac{e^{\lammin^\infty t}-1}{{\lammin^\infty}^2}\E[\norm{\ktrain - k_\infty}{}^2\norm{f_0-y}{}^2]  + \E[\norm{f_0-G_0}{}^2]\right)\\
    &\label{eq32_1} \le \frac{1-e^{-\lammin^\infty t}}{{\lammin^\infty}^2}\E[\norm{\ktrain - k_\infty}{}^4]^\frac{1}{2}\E[\norm{f_0-y}{}^4]^\frac{1}{2}  + e^{-\lammin^\infty t}\E[\norm{f_0-G_0}{}^2].
\end{align}

Now we can take the infimum on the couplings between $f_0$ and $G_0$ on both sides of the inequality to apply Theorem \ref{thm:trevisanoriginal}, Proposition \ref{prop:kernelp} and Lemma \ref{lem:4moment}
     to estimate the right-hand side of \eqref{eq32_1}.
    There are positive constants $c_1$ and $c_2$ not depending on $n_1$ nor $t$ and such that:
    \begin{align}
        \mathcal W^2_2(f_t,G_t) &\le \left(c_1\frac{1-e^{-\lammin^\infty t}}{{\lammin^\infty}^2 n_1}\sqrt{32n^2\norm{\Phi}{\infty}^4 + 8\norm{y}{}^4}  + \frac{c_2}{n_1}\right)
        \le \frac{C}{n_1}
        %&\le \frac{Ce^{-\lammin^\infty t}}{n_1}(t+1),
    \end{align}
    with $C = \max\{\frac{2c_1}{{\lammin^\infty}^2}\sqrt{8n^2\norm{\Phi}{\infty}^4 + 2\norm{y}{}^4},c_2\}$.
This proves the claim for $\X$.

Now we show the result for an arbitrary test point $x \in \R^{n_0}$.
Let $k_\infty^x = k_\infty(x,\X)$.
We can bound, by Equation \eqref{eq1},
    \[
        \frac{\partial}{\partial t} (y_t-G_t^x)^2 = 2(\ktest (y-f_t)-k_\infty^x(y-G_t^\X))(y_t-G_t^x).
    \]

By the formula for the derivative of the product and Cauchy-Schwarz's inequality, the following estimate holds almost everywhere:
    \begin{align}
        \frac{\partial}{\partial t} (y_t-G_t^x) &\le \ktest (y-f_t)-k_\infty^x(y-G_t^\X)
        = \norm{\ktest-k_\infty^x}{} \norm{y-f_t}{}-k_\infty^x(f_t-G_t^\X).
    \end{align}
    Put $\overline \lambda = \min_{1\le i\le n}k_\infty(x,x_i)$.
    The last summand can be further estimated with the first result in this theorem.
    There exists a positive constant $C$ not depending on $n_1$ such that:
    \begin{equation}
        \label{eqaux_test1}\frac{\partial}{\partial t} (y_t-G_t^x) \le \norm{\ktest-k_\infty^x}{} \norm{y-f_t}{}+\sqrt{\frac{C}{n_1}}.
    \end{equation}

    Integrating we obtain:
    \begin{align}
        \label{eqaux_test2}
        (y_t-G_t^x) &\le (y_0-G_0^x) + \norm{\ktest-k_\infty^x}{} \norm{y-f_0}{} t +\sqrt{\frac{C}{n_1}}t.
    \end{align}
    we can finish by applying the elementary inequality $(a+b+c)^2 \le 3a^2+3b^2+3c^2$ for $a,b,c\ge 1$ and Hölder's inequality.
    After that, Theorem \ref{thm:trevisanoriginal}, Proposition \ref{prop:kernelp} and Lemma \ref{lem:4moment} can be applied in the same fashion as in the proof of the training case, 
    yielding positive constants $d_1$ and $d_2$ such that:
    %\textcolor{red}{SENZA ELEVARE AL QUADRATO SI AVREBBE PURE UN RISULTATO PER LA $\mathcal W_1$. INCLUDERE?}
    \begin{align}
        \mathcal W_2^2(y_t,G_t^x) &= \E[\vert y_t-G_t^x\vert^2] \\
        &\le 3\E[\vert y_0-G_0^x\vert^2] + 3\E[\norm{\ktest-k_\infty^x}{}^4]^\frac{1}{2} \E[\norm{y-f_0}{}^4]^\frac{1}{2}t^2+\frac{3Ct^2}{n_1}\\
        &\le \frac{3d_1}{n_1} + \frac{6d_2}{n_1}\sqrt{8n^2\norm{\Phi}{\infty}^4 + 2\norm{y}{}^4}t^2+\frac{3Ct^2}{n_1}\\
        &\le C'(1+t^2)\frac{1}{n_1},
    \end{align}
    with $C' = \max\{3d_1,6d_2\sqrt{8n^2\norm{\Phi}{\infty}^4 + 2\norm{y}{}^4},3C\}$.
    Note that $C$ and $C'$ do not depend neither on $t$ nor $n_1$.
    This finishes the proof.
%\color{black}

%\input{proof_flin_G_old}
\end{proof}


\section{Proofs of the auxiliary results}
\label{ap:proofs}

We present here all the  remaining proofs.
For clearness, we will use the following notation: $X_v = h(x)_v$ and $X'_v = h(x')_v$ for each $1 \le v \le n_1$, for any $x,x'\in \R^{n_0}$.


\begin{proof}[Proof of Lemma \ref{lem_It_properties}]

    It is trivial when $B$ is nonsingular.
    Let $\lambda_1, \ldots, \lambda_n$ be the (possible repeated) ordered eigenvalues of $B$ and suppose $\lambda_j = 0$. Then, by using the eigenvalue decomposition of $B$ and elementary properties of the matrix exponential:
    \begin{align}
        I_t(B)B &= U \begin{pmatrix}
            \frac{1 - e^{-\lambda_1t}}{\lambda_1} & & & & \\
            & \ddots & & & \\
            & & t & & \\
            & & & \ddots & \\
            & & & & \frac{1 - e^{-\lambda_nt}}{\lambda_n} 
        \end{pmatrix} U^\top  U \begin{pmatrix}
            \lambda_1 & & & & \\
            & \ddots & & & \\
            & & 0 & & \\
            & & & \ddots & \\
            & & & & \lambda_n 
        \end{pmatrix} U^\top  \\
        &= U \begin{pmatrix}
            1 - e^{-\lambda_1t} & & & & \\
            & \ddots & & & \\
            & & 0 & & \\
            & & & \ddots & \\
            & & & & 1 - e^{-\lambda_nt}
        \end{pmatrix} U^\top  \\
        &= UU^\top - U \begin{pmatrix}
            e^{-\lambda_1t} & & & & \\
            & \ddots & & & \\
            & & 0 & & \\
            & & & \ddots & \\
            & & & & e^{-\lambda_nt}
        \end{pmatrix} U^\top  \\
        &= \mathbbm{1}_n - e^{-Bt}.
    \end{align}
    The converse equality is proven in an analogous way.

    As for the limit property, it is enough to show it for real numbers.
    Let $(a_n)_{n\in\N}$ be a real sequence converging to $a\in \R$.
    If $a\ne 0$, of if $a = a_n = 0$ for each $n$ the result is trivial.
    Thus, assume $a=0$ and, up to taking a subsequence, that $a_n\ne 0$.
    Then
    \[
        \lim_{n\to \infty}I_t(a_n) =\lim_{n\to \infty}\frac{1-e^{-a_nt}}{a_n} = t = I_t(0).
    \]
\end{proof}

\begin{proof}[Proof of Lemma \ref{thm:wass_properties}]
    For the proof of the first three points we refer to any monograph on the Wasserstein distance such as \cite{villani_transport}.
    In order to show (4), let $\pi^{XY}$ be an optimal transport plan between $X$ and $Y$, let $\pi^{XZ}$ be the law of $\tilde X$ and let $\mu^X = \mathbb P_X$ be the marginal law of $X$.
    Applying the gluing lemma of optimal transport produces a probability measure $\pi$ on $\R^n \times \R^m \times \R^n$ given by
    \[
    \pi(x,z,y) = \frac{\pi^{XY}(x,y)}{\mu^X(x)}\mu^X(x)\frac{\pi^{XZ}(x,z)}{\mu^X(x)} = \frac{\pi^{XY}(x,y)\pi^{XZ}(x,z)}{\mu^X(x)}.
    \]

    Note that integrating with respect to $y$ yields $\pi^{XZ}$ and integrating with respect to $z$ yields $\pi^{XY}$. 
    Note also that there is a unique coupling between $Y$ and $\lambda$, which is given by $\mathbb P_Y \times \delta_\lambda$, hence,
    \begin{align}
        \mathcal W_p^p \left(\tilde X, \tilde Y\right) &\le \int_{\R^{n+m+n}} \int_{\R^m} \norm{\tilde X(x,z) - \tilde Y(y,w)}{}^p \delta_\lambda(dw)d\pi(dx,dz,dy) \\
        &= \int_{\R^{n+m+n}} \norm{\tilde X(x,z) - \tilde Y(y,\lambda)}{}^p d\pi(dx,dz,dy) \\
        &\le \int_{\R^{n+n}} \norm{X - Y}{}^p d\pi^{XY}(dx,dy) + \int_{\R^{m}} \norm{Z-\lambda}{}^p d\mu^Z(dz) \\
        &= \mathcal W_p^p\left(X,Y\right) + \mathcal W_p^p\left(Z,\lambda\right),
    \end{align} 
    where $\mu^Z$ is the marginal probability on $Z$.

    On the other hand, to show (5) fix $\mu \in \mathcal P(\R^n)$ and $\nu \in \mathcal P(\R^m)$ two probability measures, and denote $\mu \times \nu$ the product measure on $\R^{n+m}$ with the tensor product $\sigma$-algebra.
    Then
    \begin{align}
        \mathcal W^p_p(\tilde X, \tilde Y) &\le \E_{\mu \times \nu}[\norm{\tilde X - \tilde Y}{}^p] \\
        &\le \E_{\mu \times \nu}[(\norm{X-Y}{}^2 + \norm{Z-V}{}^2)^\frac{p}{2}]\\
        &\le 2^{\frac{p}{2}-1}\left(\E_\mu[\norm{X - Y}{}^{2p}] + \E_\nu[\norm{Z - V}{}^{2p}]\right),
    \end{align}
    where in the last step we applied the elementary inequality $(a+b)^p \le 2^{p-1}(a^p + b^p)$, for $a,b \ge 0$ and $p \ge 1$.
    For the $1$-Wasserstein distance instead, we apply the elementary inequality $\sqrt{a+b} \le \sqrt{a} + \sqrt{b}$:
    \begin{align}
        \mathcal W_1(\tilde X, \tilde Y) &\le \E_{\mu \times \nu}[\norm{\tilde X - \tilde Y}{}] \\
        &\le \E_{\mu \times \nu}[\left(\norm{X-Y}{}^2 + \norm{Z-V}{}^2\right)^\frac{1}{2}]\\
        &\le \E_\mu[\norm{X - Y}{}] + \E_\nu[\norm{Z - V}{}].
    \end{align}
    
Lastly, taking the infimum over $(\mu, \nu) \in \mathcal P(\R^n) \times \mathcal P(\R^m)$ finishes the proof.
\end{proof}

\begin{proof}[Proof of Lemma \ref{lemma1}]
    The computation of the gradients is by chain rule and definition of $f$.
    The claims about the kernels follow from taking the dot product on the gradients we just calculated, for each $1\le i,j \le n_1$:
    \begin{align}
        \tilde k_{ij}(x,x') &= \left(\nabla_{\theta^{(0)}}X_i\right)\left(\nabla_{\theta^{(0)}}X'_j\right) \\
        &=\frac{1}{n_0}\sum_{\substack{u=1, \ldots, n_0 \\v=1, \ldots, n_1}} \partv{\theta^{(0)}_{uv}}(x\theta^{(0)}_{\_ i})\partv{\theta^{(0)}_{uv}}(x'\theta^{(0)}_{\_ j}) \\
        &=\frac{1}{n_0}\sum_{\substack{u=1, \ldots, n_0 \\v=1, \ldots, n_1}} x_u\delta_{iv}x'_u\delta_{jv} \\
        &=\frac{1}{n_0}\sum_{u=1}^{n_0}x_ux'_u\delta_{ij}.
    \end{align}
    Lastly,
    \begin{align}
        k(x,x') &= \left(\nabla_{\theta^{(0)}}f^{(2)}(x)\right)\left(\nabla_{\theta^{(0)}}f^{(2)}(x')\right) + \left(\nabla_{\theta^{(1)}}f^{(2)}(x)\right)\left(\nabla_{\theta^{(1)}}f^{(2)}(x')\right) \\
        &=\frac{1}{n_1}\sum_{\substack{u=1, \ldots, n_0 \\v=1, \ldots, n_1}} \partv{\theta^{(0)}_{uv}}(\Phi(\frac{1}{\sqrt{n_0}}x\theta^{(0)})\theta^{(1)})\partv{\theta^{(0)}_{uv}}(\Phi(\frac{1}{\sqrt{n_0}}x'\theta^{(0)})\theta^{(1)}) \\
        &\quad + \frac{1}{n_1}\sum_{z=1}^{n_1} \partv{\theta^{(1)}_z}(\Phi(\frac{1}{\sqrt{n_0}}x\theta^{(0)})\theta^{(1)})\partv{\theta^{(1)}_z}(\Phi(\frac{1}{\sqrt{n_0}}x'\theta^{(0)})\theta^{(1)}) \\
        &=\frac{1}{n_1n_0}\sum_{\substack{u=1, \ldots, n_0 \\v=1, \ldots, n_1}} x_ux'_u\Phi'(X_v)\Phi'(X_v)(\theta^{(1)}_v)^2 + \frac{1}{n_1}\sum_{z=1}^{n_1}\Phi(X_z)\Phi(X'_z).
    \end{align}
\end{proof}
    
\subsection{Proof of results at initialization}

In this subsection we prove the useful Proposition \ref{prop:kernelp}, which generalises the second part of Theorem \ref{thm:trevisanoriginal}.
This step is paramount for the proof of the rest of our results.
From now to the end of this subsection assume $\Phi$ and $\Phi'$ are bounded and $x,x'\in \R^{n_0}$ are fixed.

\begin{comment}
    In the next Lemma we collect some well-known properties of the $p$-Wasserstein distance:

\begin{lem}\label{thm:wass_properties}
    Let $p \in [1, \infty[$ and let $X,Y$ be random variables with values in $\R^n$ and $Z$ be a random variable with values in $\R^m$. 
    Let $\mathbb P_\xi$ denote the law of the random variable $\xi$ per $\xi \in \{X,Y,Z\}$. Then
    \begin{enumerate}
        \item If $X, Y$ are defined on the same probability space, then $\mathcal W_p(X,Y) \le \E[\norm{X - Y}{}^p]^{\frac{1}{p}}$.
        \item If $Z$ is independent from $X$ and $Y$ then $\mathcal W_p(X+Z,Y+Z) \le \mathcal W_p(X,Y)$.
        \item Convexity of $\mathcal W_p^p$: $\mathcal W_p^p(X,Y) = \int_{\R^m} \mathcal W_p^p(\mathbb P_{X \vert Z = z},\mathbb P_Y)d\mathbb P_Z(z)$. 
         \item Let $\lambda \in \R^m$ be a constant vector and consider the joint random variables $\tilde X = (X,Z), \tilde Y = (Y, \lambda)$. Then 
        \[\mathcal W_p^p(\tilde X, \tilde Y) \le \mathcal W_p^p(X, Y) + \mathcal W_p^p(Z, \lambda) = \mathcal W_p^p(X, Y) + \left(\int_{\R^m} d(z,\lambda)^p d\mathbb P(z)\right)^{\frac{1}{p}},\]
        where $d$ denotes the euclidean distance in $\R^m$. 
        \item Let $V$ be a random variable with values in $\R^m$ and consider the joint random variables $\tilde X = (X,Z), \tilde Y = (Y, V)$. Then, for $p \ge 2$, 
        \[\mathcal W_p^p(\tilde X, \tilde Y) \le 2^{\frac{p}{2}-1}\left(\mathcal W_{2p}^{2p}(X, Y) + \mathcal W_{2p}^{2p}(Z, V)\right).\]

        Moreover, for $p=1$,
        \[\mathcal W_1(\tilde X, \tilde Y) \le \mathcal W_1(X, Y) + \mathcal W_1(Z, V).\]
    \end{enumerate}
\end{lem}
\end{comment}

\begin{proof}[Proof of Proposition \ref{prop:kernelbound}]
    Note that, from Lemma \ref{lemma1}, $k$ can be written as:
    \[
    k = \frac{1}{n_1}\tilde k_{11}\sum_{v=1}^{n_1}\Phi'(X_v)\Phi'(X'_v)(\theta^{(1)}_v)^2 + \frac{1}{n_1}\sum_{v=1}^{n_1}\Phi(X_v)\Phi(X'_v),
    \]
    and recall that $\tilde k$ is deterministic, for any fixed inputs $x,x'$.
    Hence, by boundedness of $\Phi$ and $\Phi'$, and independence of $\theta^{(1)}$ and $\theta^{(0)}$,
    \begin{align}
        \E[\vert k\vert] &\le (\norm{\Phi'}{\infty}^2\E[\vert \tilde k_{11}\vert] + \norm{\Phi}{\infty}^2) = \norm{\Phi}{\infty}^2,\\
        \E[\vert k\vert^p] &= \E[\vert \frac{1}{n_1}\tilde k_{11}\sum_{v=1}^{n_1}\Phi'(X_v)\Phi'(X'_v)(\theta^{(1)}_v)^2 + \frac{1}{n_1}\sum_{v=1}^{n_1}\Phi(X_v)\Phi(X'_v)\vert^p]\\
        &\le 2^{p-1}\norm{\Phi'}{\infty}^{2p}\E[\vert \tilde k_{11}\vert^p]\E[(\theta^{(1)}_1)^{2p}] + 2^{p-1}\norm{\Phi}{\infty}^{2p}\\
        &\le 2^{p-1}(2p-1)!!\norm{\Phi'}{\infty}^{2p}\vert \tilde k_{11}\vert^p + 2^{p-1}\norm{\Phi}{\infty}^{2p},
    \end{align}
    where in the last inequality we used that the $2p$-th moment of a standard Gaussian variable is equal to $(2p-1)!!$.
    
\end{proof}

\begin{proof}[Proof of Proposition \ref{prop:kernelp}]
    We will use the notation instroduced in Proposition \ref{prop:kernelbound}.
    The first claim is trivial since $\tilde k$ coincides with $\mathcal K$.
    To show the second claim, we split:
    \begin{align}
        \E[\vert k_{11}-k_\infty\vert^p] &= \E[\vert \frac{1}{n_1}\tilde k_{11}\sum_{v=1}^{n_1}\Phi'(X_v)\Phi'(X'_v)(\theta^{(1)}_v)^2 + \frac{1}{n_1}\sum_{v=1}^{n_1}\Phi(X_v)\Phi(X'_v) \\
        &\quad - \mathcal K\E_G[\Phi'(G(x))\Phi'(G(x'))] - \E_G[\Phi(G(x))\Phi(G(x'))]\vert^p] \\
        &\label{AUXLPKER1}\le 2^{p-1}\E[\vert \frac{1}{n_1}\tilde k_{11}\sum_{v=1}^{n_1}\Phi'(X_v)\Phi'(X'_v)(\theta^{(1)}_v)^2 - \mathcal K\E_G[\Phi'(G(x))\Phi'(G(x'))]\vert^p] \\
        &\quad  + 2^{p-1}\E[\frac{1}{n_1}\sum_{v=1}^{n_1}\Phi(X_v)\Phi(X'_v)- \E_G[\Phi(G(x))\Phi(G(x'))]\vert^p].
    \end{align}

    To bound the first summand in \eqref{AUXLPKER1} we split again by adding and substracting an auxiliary term:
    \begin{align}
        \E[\vert &\frac{1}{n_1}\tilde k_{11}\sum_{v=1}^{n_1}\Phi'(X_v)\Phi'(X'_v)(\theta^{(1)}_v)^2 - \mathcal K\E_G[\Phi'(G(x))\Phi'(G(x'))]\vert^p] \\
        &\le 2^{p-1}\E[\vert \frac{1}{n_1}\tilde k_{11}\sum_{v=1}^{n_1}\Phi'(X_v)\Phi'(X'_v)(\theta^{(1)}_v)^2 - \frac{1}{n_1}\mathcal K\sum_{v=1}^{n_1}\Phi'(X_v)\Phi'(X'_v)(\theta^{(1)}_v)^2\vert^p] \\
        &\quad + 2^{p-1}\E[\vert \frac{1}{n_1}\mathcal K\sum_{v=1}^{n_1}\Phi'(X_v)\Phi'(X'_v)(\theta^{(1)}_v)^2 - \mathcal K\E_G[\Phi'(G(x))\Phi'(G(x'))]\vert^p] \\
        &\label{AUXLPKER2}= 2^{p-1}\E[\vert \frac{1}{n_1}(\tilde k_{11}-\mathcal K)\sum_{v=1}^{n_1}\Phi'(X_v)\Phi'(X'_v)(\theta^{(1)}_v)^2\vert^p] \\
        &\quad + 2^{p-1}(\mathcal K)^p\E[\vert \frac{1}{n_1}\sum_{v=1}^{n_1}\Phi'(X_v)\Phi'(X'_v)(\theta^{(1)}_v)^2 - \E_G[\Phi'(G(x))\Phi'(G(x'))]\vert^p].
    \end{align}

    The first summand in \eqref{AUXLPKER2} vanishes since $\tilde k$ equals $\mathcal K$.
    We estimate the second summand in \eqref{AUXLPKER2}, once again by adding and substracting an auxiliary term:
    \begin{align}
        \E[\vert &\frac{1}{n_1}\sum_{v=1}^{n_1}\Phi'(X_v)\Phi'(X'_v)(\theta^{(1)}_v)^2 - \E_G[\Phi'(G(x))\Phi'(G(x'))]\vert^p]\\
        &\label{AUXLPKER3}\le 2^{p-1}\E[\vert \frac{1}{n_1}\sum_{v=1}^{n_1}\Phi'(X_v)\Phi'(X'_v)((\theta^{(1)}_v)^2-1)\vert^p] \\
        &\quad + 2^{p-1}\E[\vert\frac{1}{n_1}\sum_{v=1}^{n_1}\Phi'(X_v)\Phi'(X'_v) - \E_G[\Phi'(G(x))\Phi'(G(x'))]\vert^p].
    \end{align}

    The first summand in \eqref{AUXLPKER3} vanishes, by boundedness of $\Phi'$ and independence of the parameters $\theta^{(1)}_v$:
    \begin{align}
        \E[\vert \frac{1}{n_1}\sum_{v=1}^{n_1}\Phi'(X_v)\Phi'(X'_v)((\theta^{(1)}_v)^2-1)\vert^p] &\le \frac{1}{n_1^p}\norm{\Phi'}{\infty}^{2p}\E[\vert \sum_{v=1}^{n_1}(\theta^{(1)}_v)^2-1 \vert^p]\\
        &= \frac{1}{n_1^p}\norm{\Phi'}{\infty}^{2p}\sum_{\alpha_1,\ldots, \alpha_p = 1}^{n_1}\prod_{j=1}^p\E[(\theta^{(1)}_{\alpha_j})^2-1]\\
        &= 0.
    \end{align}

    As for the second summand in \eqref{AUXLPKER3}, by Theorem \ref{thm:trevisanoriginal} there exists a constant $c_1$ not depending on $n_1$ such that:
    \begin{align}
        \mathcal W^p_p(&\frac{1}{n_1}\sum_{v=1}^{n_1}\Phi'(X_v)\Phi'(X'_v),\E_G[\Phi'(G(x))\Phi'(G(x'))])\\
        &=\E[\vert\frac{1}{n_1}\sum_{v=1}^{n_1}\Phi'(X_v)\Phi'(X'_v) - \E_G[\Phi'(G(x))\Phi'(G(x'))]\vert^p]\\
        &\le c_1 \left(\frac{\textup{Lip}\Phi' + \Phi'(0)}{\sqrt{n_1}}\right)^p.
    \end{align}

    It remains only to bound the second summand in \eqref{AUXLPKER1}.
    This is done by using again Theorem \ref{thm:trevisanoriginal}.
    There exists a constant $c_2$ not depending on $n_1$ such that:
    \begin{align}
        \mathcal W^p_p(&\frac{1}{n_1}\sum_{v=1}^{n_1}\Phi(X_v)\Phi(X'_v),\E_G[\Phi(G(x))\Phi(G(x'))])\\
        &=\E[\frac{1}{n_1}\sum_{v=1}^{n_1}\Phi(X_v)\Phi(X'_v)- \E_G[\Phi(G(x))\Phi(G(x'))]\vert^p]\\
        &\le c_2  \left(\frac{\textup{Lip}\Phi + \Phi(0)}{\sqrt{n_1}}\right)^p.
    \end{align}

    Putting together all the preceding estimations we obtain: 
    \begin{align}
        \E[\vert k_{11}-k_\infty\vert^p] &\le C \frac{1}{n_1^\frac{p}{2}},
    \end{align}
    with $C = 2^{p-1}\max\{2^{2p-2}c_1(\textup{Lip}\Phi' + \Phi'(0)), c_2(\textup{Lip}\Phi + \Phi(0))\}$.
\end{proof}

These results suffice to prove Proposition \ref{prop:trevisan}.

\begin{proof}[Proof of Proposition \ref{prop:trevisan}]

    Consider the joint random variables $\tilde X = (\ktrain,\hat f(\X))$ and $\tilde Y = (k_\infty(\X,\X),G(\X))$.
Then Lemma \ref{thm:wass_properties}.4 together with Proposition \ref{prop:kernelp} and Theorem \ref{thm:trevisanoriginal} yield
\begin{align}
    \mathcal W_p(\tilde X, \tilde Y) &\le \mathcal W_p(\ktrain,k_\infty(\X,\X)) + \mathcal W_p(f(\X),G(\X)) \le \frac{C+D}{\sqrt{n_1}},
\end{align}
where $C$ is the constant in Proposition \ref{prop:kernelp} and $D$ the one in Theorem \ref{thm:trevisanoriginal}.
Both constants do not depend on $n_1$.
\end{proof}

Lastly, we prove Lemma \ref{lem:4moment}.

\begin{proof}[Proof of Lemma \ref{lem:4moment}]
    Let $\theta^{(0)}_{ij}$ denote the $ij$-th entry of $\theta^{(0)}_0 \in \R^{n_0\times n_1}$, and let $\theta^{(1)}_j$ denote the $j$-th component of $\theta^{(1)}_0 \in \R^{n_1}$.
    By Jensen's inequality and independence of the parameters and $x_1, \ldots x_n$:
    \begin{align}
        \E[\norm{f_0}{}] &\le \sqrt{\E[\norm{\frac{1}{\sqrt{n_1}}\Phi(\X\theta^{(0)}_0)\theta^{(1)}_0}{}^2]} \\
        &\le \sqrt{n}\sqrt{\E[\vert \Phi(x_1\theta^{(0)}_{\_ 1})\vert^2]\E[\vert \theta^{(1)}_1\vert^2]} \\
        &\le \sqrt{n}\norm{\Phi}{\infty}.
    \end{align}

    As for the fourth moment,
    \begin{align}
        \E[\norm{f_0}{}^4] &= \E\left[\left(\sum_{i=1}^n\frac{1}{n_1}\left(\sum_{j=1}^{n_1}\Phi(x_i\theta^{(0)}_{\_ j})\theta^{(1)}_j\right)^2\right)^2\right] \\
        &\le \frac{n^2}{n_1^2}\E\left[\left(\sum_{j=1}^{n_1}\Phi(x_1\theta^{(0)}_{\_ j})\theta^{(1)}_j\right)^4\right] \\
        &\le \frac{\norm{\Phi}{\infty}^4n^2}{n_1^2}(3n_1 + n_1^2) \\
        &\le 4 n^2 \norm{\Phi}{\infty}^4.
    \end{align}
    Triangular inequality finishes the proof.
\end{proof}

\subsection{Approximation of the network by linearization}

In this subsection we prove the results involved in the proof of Proposition \ref{quasi_main}.

With a slight abuse of notation, we will denote by $\norm{x-\X}{}$ the positive quantity $\sup_{1 \le i \le n}\norm{x-x_i}{}$.
Also, given any matrix $A = (a_{ij})_{\substack{1 \le i \le n \\1 \le j \le m}}$ we will denote by $\partv{A}f$ the matrix $\nabla_A f = (\partv{A_{ij}})_{\substack{1 \le i \le n \\1 \le j \le m}}$.
We will consider the \emph{linearized gradient flow}, given by
\[
    \frac{\partial}{\partial t} {\overline \theta_t} = -\nabla_\theta f_0 (f^\lin(\X;\overline \theta_t)-y).
\]

For this subsection introduce the following notations: $f^\lin_t =f^\lin(\X;\overline \theta_t)$ and $\overline y_t = f^\lin(x;\overline \theta_t)$.

    \begin{proof}[Proof of Lemma \ref{lem:gronwallbeta}]
        Let $\lammin$ be the smallest eigenvalue of $k_t$.
        By gradient flow equations for the parameters $\theta_t$ and $\overline \theta_t$: 
        \begin{align}
            \label{eqaux_gftrucco}
            \norm{f_t-y}{} &\le e^{-\lammin t} \norm{f_0 - y}{}, \\
            \norm{f_t^\lin-y}{} &\le e^{-\lammin^0 t} \norm{f_0 - y}{}.
        \end{align}
        
        On the other hand,
        Lemma \ref{lemma1} combined with Cauchy-Schwarz's inequality and the gradient flow equations produces the following system of differential inequalities:
    \begin{align}
        \label{eqaux_gfineq}
        \frac{\partial}{\partial t} {(\theta_v^{(1)})_t} &\le \frac{1}{\sqrt{n_1}}\norm{\Phi}{\infty}\norm{f_0-y}{} e^{-\lammin t},\\
        \frac{\partial}{\partial t} {(\theta_{uv}^{(0)})_t} &\le \frac{1}{\sqrt{n_1n_0}}\norm{\Phi'}{\infty}\norm{f_0-y}{}\norm{\X_u}{}(\theta^{(1)}_v)_t e^{-\lammin t}.
    \end{align}
    The previous is a triangular system of differential inequalities of the form
    \[\begin{cases}
        \frac{\partial}{\partial t} {(\theta_v^{(1)})_t} &\le B_1  e^{-\lammin t} \\
        \frac{\partial}{\partial t} {(\theta_{uv}^{(0)})_t} &\le B_0 (\theta^{(1)}_v)_t  e^{-\lammin t},
    \end{cases}\] 
    with $B_1=\frac{1}{\sqrt{n_1}}\norm{\Phi}{\infty}\norm{f_0-y}{}$ and $B_0=\frac{1}{\sqrt{n_1n_0}}\norm{\Phi'}{\infty}\norm{f_0-y}{}\norm{\X_u}{}$.
    
    By integration on $[0,t]$ and substitution we get:,
    \begin{align}
        \label{eqaux_gronwalls}
        (\theta_v^{(1)})_t &\le (\theta_v^{(1)})_0 + B_1I_t(\lammin)\\
        &\le(\theta_v^{(1)})_0 + \frac{\norm{\Phi}{\infty}\norm{f_0-y}{}}{\sqrt{n_1}}I_t(\lammin),\\
        (\theta_{uv}^{(0)})_t &\le (\theta_{uv}^{(0)})_0 + B_0B_1\int_0^tI_s(\lammin)ds + B_0\norm{\theta^{(1)}_0}{}I_t(\lammin)\\
        &\le (\theta_{uv}^{(0)})_0 + \frac{\norm{\Phi}{\infty}\norm{\Phi'}{\infty}\norm{f_0-y}{}^2\norm{\X_u}{}}{2n_1\sqrt{n_0}}I_t(\lammin)^2 \\
        &\quad + \frac{\norm{\Phi'}{\infty}\norm{f_0-y}{}\norm{\X_u}{}}{\sqrt{n_1n_0}}I_t(\lammin)(\theta^{(1)}_v)_0.
    \end{align}
    Note that in the last inequality, we used $\int_0^t I_s(b)ds \le \frac{I_t(b)^2}{2}$, for any $b\ge 0$.
    
    Thanks to \eqref{eqaux_gftrucco}, the linearised parameters $\overline \theta_t$ also satisfy the preceding inequalities, and hence the thesis holds.
    \end{proof}

\begin{proof}[Proof of Lemma \ref{lem:supnorm}]
    Let $\Phi$ denote the CDF of a standard Gaussian variable.
    For each $a>0$, since the entries of $\theta_0^{(1)}$ are $n_1$ i.i.d. standard Gaussian variables,
    \begin{align}
        \mathbb P(\norm{\theta_0^{(1)}}{\infty} \le a) &= (1-2(1-\Phi(a)))^{n_1}.
    \end{align}
    Bernouilli's inequality and standard estimations for Gaussian tails yield
    \begin{align}
        \mathbb P(\norm{\theta_0^{(1)}}{\infty}\le a) &\ge 1 - 2n_1(1-\Phi(a)) \\
        &\ge 1 - n_1\exp\left(-\frac{a^2}{2}\right).
    \end{align}

    Let $r\ge 1$ and put $a = \sqrt{r\gamma \log n_1}$. 
    Then:
    \begin{align}
        \mathbb P(\norm{\theta_0^{(1)}}{\infty}\le a) &\ge 1 - n_1\exp\left(-\frac{r\gamma \log n_1}{2}\right) \\
        &= 1 - n_1\exp\left(\log n_1^{-\frac{r\gamma}{2}}\right)\\
        &= 1 - \frac{1}{n_1^{\frac{r\gamma}{2}-1}}.
    \end{align}

\end{proof}

\begin{proof}[Proof of Lemma \ref{lem:jacobian}]
    We will write $f$ for short of $f_0(x)(\theta)$.
    By Lemma \ref{lemma1} and Cauchy-Schwarz's inequality,
    \begin{align}
        \norm{\nabla_\theta f}{}^2 &= \norm{\partv{\theta^{(0)}}f}{}^2 + \norm{\partv{\theta^{(1)}}f}{}^2 \\
        &\le \frac{1}{n_0n_1}\norm{x}{}^2\norm{\Phi'}{\infty}^2\norm{\theta^{(1)}}{}^2 + \frac{1}{n_1}\norm{\Phi}{\infty}^2.
    \end{align}
    
    Then the first claim follows by \eqref{eqmassart_laurent} and the elementary inequality $\sqrt{a+b} \le \sqrt a + \sqrt b$, for $a,b\ge 0$.

    Now we prove the second inequality.
    Let $\theta, \tilde \theta \in \R^N$, then,
    \begin{align}
        \norm{\nabla_\theta f(\theta)-\nabla_\theta f(\tilde \theta)}{}^2 &= \norm{\partv{\theta^{(0)}} f(\theta)-\partv{\theta^{(0)}} f(\tilde \theta)}{}^2 + \norm{\partv{\theta^{(1)}} f(\theta)-\partv{\theta^{(1)}} f(\tilde \theta)}{}^2.
    \end{align}
    
    Let us estimate the first summand in the previous expression.
    By Lemma \ref{lemma1}, for each $1 \le u \le n_0, 1 \le v \le n_1$,
    \begin{align}
        \vert \partv{\theta^{(0)}_{uv}} &f(\theta )-\partv{\theta^{(0)}_{uv}} f(\tilde \theta) \vert \\
        &\le \frac{1}{\sqrt{n_1n_0}}x_u(\Phi'(\frac{1}{\sqrt{n_0}}\sum_{j=1}^{n_0}x_j\theta^{(0)}_{jv})\theta^{(1)}_v-\Phi'(\frac{1}{\sqrt{n_0}}\sum_{j=1}^{n_0}x_j\tilde \theta^{(0)}_{jv})\tilde\theta^{(1)}_v)\\
        &\le \frac{x_u}{\sqrt{n_1n_0}}\Phi'(\frac{1}{\sqrt{n_0}}\sum_{j=1}^{n_0}x_j\theta^{(0)}_{jv})(\theta^{(1)}_v - \tilde \theta^{(1)}_v)\\
        &\quad + \frac{x_u\tilde \theta^{(1)}_v}{\sqrt{n_1n_0}}x_u(\Phi'(\frac{1}{\sqrt{n_0}}\sum_{j=1}^{n_0}x_j\theta^{(0)}_{jv})-\Phi'(\frac{1}{\sqrt{n_0}}\sum_{j=1}^{n_0}x_j\tilde \theta^{(0)}_{jv}))\\
        &\le \frac{x_u\norm{\Phi'}{\infty}(\theta^{(1)}_v-\tilde \theta^{(1)}_v)}{\sqrt{n_1n_0}} + \frac{x_u\textup{Lip}\Phi'\tilde \theta^{(1)}_v}{\sqrt{n_1}n_0}\sum_{j=1}^{n_0}x_j(\theta^{(0)}_{jv}-\tilde\theta^{(0)}_{jv}).
    \end{align}
    Hence,
    \begin{align}
        \norm{\partv{\theta^{(0)}} &f(\theta )-\partv{\theta^{(0)}} f(\tilde \theta)}{}^2 \\
        &= \sum_{\substack{u=1, \ldots, n_0 \\v=1, \ldots, n_1}}\vert \partv{\theta^{(0)}_{uv}} f(\theta )-\partv{\theta^{(0)}_{uv}} f(\tilde \theta) \vert^2 \\
        &\le \frac{\norm{x}{}^2\norm{\Phi'}{\infty}^2\norm{\theta^{(1)}-\tilde \theta^{(1)}}{}^2}{n_1n_0} + \frac{\norm{x}{}^4(\textup{Lip}\Phi')^2\norm{\tilde \theta^{(1)}}{\infty}^2\norm{\theta^{(0)}-\tilde \theta^{(0)}}{}^2}{n_1n_0^2}\\
        &\le \frac{\norm{x}{}^2\norm{\Phi'}{\infty}^2\norm{\theta^{(1)}-\tilde \theta^{(1)}}{}^2}{n_1n_0} + \frac{\norm{x}{}^4(\textup{Lip}\Phi')^2\norm{\theta^{(0)}-\tilde \theta^{(0)}}{}^2}{n_1n_0^2}r\gamma\log n_1,
    \end{align}
    with probability greater or equal than $1-\frac{1}{n^{\frac{r\gamma}{2}-1}}$, where in the last step we used Lemma \ref{lem:supnorm}.

    Similarly, the second summand can be estimated as follows.
    First compute the partial derivatives by using Lemma \ref{lemma1}, for each $1 \le v \le n_1$:
    \begin{align}
        \vert \partv{\theta^{(1)}_v} &f(\theta )-\partv{\theta^{(1)}_v} f(\tilde \theta) \vert \\
        &\le \frac{1}{\sqrt{n_1}}(\Phi(\frac{1}{\sqrt{n_0}}\sum_{j=1}^{n_0}x_j\theta^{(0)}_{jv})-\Phi(\frac{1}{\sqrt{n_0}}\sum_{j=1}^{n_0}x_j\tilde \theta^{(0)}_{jv}))\\
        &\le \frac{\textup{Lip}\Phi}{\sqrt{n_1n_0}}\sum_{j=1}^{n_0}x_j(\theta^{(0)}_{jv}-\tilde \theta^{(0)}_{jv}).
    \end{align}
    Therefore,
    \begin{align}
        \norm{\partv{\theta^{(1)}} &f(\theta )-\partv{\theta^{(1)}} f(\tilde \theta)}{}^2 \\
        &= \sum_{v=1, \ldots, n_1}\vert \partv{\theta^{(1)}_{v}} f(\theta )-\partv{\theta^{(1)}_{v}} f(\tilde \theta) \vert^2 \\
        &\le \frac{(\textup{Lip}\Phi)^2 \norm{x}{}^2}{n_1n_0}\norm{\theta^{(0)}-\tilde\theta^{(0)}}{}^2.
    \end{align}
    The preceding estimations, together with $\frac{\norm{\theta^{(i)}-\tilde\theta^{(i)}}{}^2}{\norm{\theta - \tilde \theta}{}^2} \le 1$ for $i = 0,1$, yield:
    \begin{align}
        &\frac{\norm{\nabla_\theta f(\theta)-\nabla_\theta f(\tilde \theta)}{}^2}{\norm{\theta - \tilde \theta}{}^2} \\
        &\le \frac{\norm{x}{}^2\norm{\Phi'}{\infty}^2}{n_1n_0} + \frac{\norm{x}{}^4(\textup{Lip}\Phi')^2}{n_1n_0^2}r\gamma \log n_1 + \frac{(\textup{Lip}\Phi)^2 \norm{x}{}^2}{n_1n_0}.
    \end{align}
    Taking the square root in the last inequality yields the thisis.
\end{proof}
 
\begin{proof}[Proof of Lemma \ref{lem:lammin}]
    Fix $\gamma \in \N$.
    The probability of $Z= \norm{k-k_\infty}{} > \frac{\gamma \lammin^\infty}{2}$ can be estimated with Markov's inequality and Proposition \ref{prop:kernelp}.
    There exists a constant $C>0$ not depending on $n_1$ such that:
    \begin{align}
        \mathbb P(Z > \frac{\gamma\lammin^\infty}{2}) &= \mathbb P\left(Z^p > \left(\frac{\gamma\lammin^\infty}{2}\right)^p\right)\\
        &\le \left(\frac{2}{\gamma\lammin^\infty}\right)^p\E[\norm{Z}{}^p]\\
        &=\left(\frac{2}{\gamma \lammin^\infty}\right)^p\mathcal W^p_p(k,k_\infty) \\
        &\le \left(\frac{2}{\gamma \lammin^\infty}\right)^p\frac{C}{n_1^\frac{p}{2}}.
    \end{align}
    Note that Proposition \ref{prop:kernelp} holds for every natural $p$.
    This concludes the proof.
\end{proof}

Now are ready to prove Proposition \ref{thm:montanari}:

\begin{proof}[Proof of Proposition \ref{thm:montanari}]

For the sake of clearness we introduce the following abbreviations for the remainder of the proof.
    Let $y_t = f(x;\theta_t), \overline y_t = f^\lin(x;\overline \theta_t), f_t = f(\X;\theta_t)$ and $f^\lin_t = f^\lin(\X;\overline \theta_t)$.
Also, let $k_t = k(\X,\X;\theta_t)$, $\nabla = \nabla_\theta$ and let $L(\X)$ denote the Lipschitz constant of $\nabla f$, seen as a function of $\theta$.

Consider the empirical risk for the quadratic loss $\mathcal R_{\mathcal D}(\theta_t) = \frac{1}{2} \sum_{i=1}^{n}(f^{(L)}(x_i;\theta_t)-y)^2$.

From gradient flow equations we have:
\begin{align}
    \frac{\partial}{\partial t} f_t &= -k_t(f_t-y),\\
    \label{eqaux_montanariB7}\partv{t}\norm{f_t-y}{}^2 &= -2\langle f_t-y, k_t(f_t-y)\rangle.
\end{align}

Let $t_* = \inf\{t \mid \norm{\theta_t-\theta_0}{} > \frac{\sigmin}{2L(\X)}\}$ 
%\textcolor{teal}{$t_* = \inf\{t \mid \norm{\theta_t-\theta_0}{} > \frac{\sigmin}{2L_n} + \int \vert S(t)\vert dt\}$}. 
Then for each $t \le t_*$, by $1$-Lipschitzianity of the smallest eigenvalue with respect to the operator norm, and by definition of $t_*$, we obtain an upper bound for $\lammin(k_t)$:
\[
    \vert \lammin(k_t) - \lammin \vert \le \norm{k_t-k_0}{op} \le\norm{k_t-k_0}{} \le L(\X)\norm{\theta_t-\theta_0}{}\le \frac{\sigmin}{2},
\]
which implies:
\[
\lammin(k_t) \ge \lammin - \frac{\sigmin}{2} \ge \frac{\lammin}{4}.
\]
%following the reasoning in Theorem 2.4 in \cite{bachchizat}, $\lammin(k_t) \ge \frac{\lammin}{4}$.
This estimation combined with Grönwall's inequality applied to \eqref{eqaux_montanariB7} yield:
\begin{align}
    \label{equax_montanariB8}
    \norm{f_t-y}{}^2 &\le \norm{f_0-y}{}^2 \exp\left(-\frac{\lammin}{2} t\right).
\end{align}

 From \eqref{eqaux_montanariB7} and Cauchy-Schwarz we deduce:
\begin{align}
    \partv{t}\norm{f_t-y}{} &= -\frac{\norm{\nabla f_t (f_t-y)}{}^2}{\norm{f_t-y}{}} \\
    &\le -\frac{\sigmin}{2}\norm{\nabla f_t(f_t-y)}{}.
\end{align}

Hence,
\begin{align}
    \label{eqaux_montanari_preB3}
    \partv{t}\left(\norm{f_t-y}{}+ \frac{\sigmin}{2}\norm{\theta_t-\theta_0}{}\right) &\le\partv{t}\norm{f_t-y}{}+\frac{\sigmin}{2}\norm{\frac{\partial}{\partial t} \theta_t}{}
    %&\le -\frac{\lambda}{n} \frac{\langle y_t-y,  Df_n(\theta_t)\theta_t\rangle}{\norm{y_t-y}{}} \textcolor{red}{= \frac{\lambda}{n\norm{y_t-y}{}}S(t)},
    \le 0.
\end{align}
for all $t \le t_*$.

Thus, for all $t \le t_*$:
\begin{align}
    \label{eqaux_montanariB3}
    \norm{\theta_t-\theta_0}{} \le \frac{2}{\sigmin}\norm{f_0-y}{}.
\end{align}
    
    Let us show that this property holds for all $t>0$.
    By contradiction assume $t_* < \infty$.
    \eqref{eqaux_montanariB3} with Assumption \ref{assumption_montanari} implies
    \begin{align}
    	\norm{\theta_{t_*}-\theta_0}{} &< \frac{2}{\sigmin}\frac{\sigmin^2}{4L(\X)}\\
    	&= \frac{\sigmin}{2L(\X)}.
    \end{align}
    In particular the last inequality holds for $t_*$, which contradicts its definition.
    Hence $t_* = \infty$.
     
Let us now prove the rest of the inequalities in the theorem.

    The gradient flow equation for the linearised network reads:
    \begin{align}
        \label{eqaux_montanariflinGF}
     	\frac{\partial}{\partial t} {f^\lin_t}=-k_0(f^\lin_t - y).
     \end{align}
     Define the difference $r_t=f_t-f^\lin_t$.
     Then 
     \begin{align}
		\frac{\partial}{\partial t} {r_t}&=-k_t(f_t-y)+k_0(f^\lin_t-y)\\
		&= -k_tr_t - (k_t-k_0)(f^\lin_t - y).
	\end{align}     
         Then, by Cauchy-Schwarz and \eqref{equax_montanariB8} combined with \eqref{eqaux_montanariflinGF},
         \begin{align}
         	\frac{1}{2}\frac{\partial}{\partial t}\norm{r_t}{}^2 &= -\langle r_t, k_tr_t\rangle -\langle r_t, (k_t-k_0)(f^\lin_t-y)\rangle\\
         	&\le -\lammin(k_t)\norm{r_t}{}^2 +\norm{r_t}{}\norm{k_t-k_0}{}\norm{f^\lin_t-y}{}\\
         	&\le  -\frac{\lammin}{4}\norm{r_t}{}^2 +\norm{r_t}{}\norm{k_t-k_0}{}\norm{f_0-y}{}\exp\left(-\frac{\lammin t}{4}\right). 
         \end{align}
         Hence,
         \begin{align}
         	\frac{\partial}{\partial t}\norm{r_t}{} &\le  -\frac{\lammin}{4}\norm{r_t}{} +\norm{k_t-k_0}{}\norm{f_0-y}{}\exp\left(-\frac{\lammin t}{4}\right).
         \end{align}
         
         Now let us bound separately the different factors in the previous equation.
         The norm of the difference between the kernels can be estimated as:
         \begin{align}
         	\norm{k_t-k_0}{} &\le \norm{\nabla f_t \nabla f_t^\top - \nabla f_0 \nabla f_0^\top}{}\\
            &\le 2\norm{\nabla f_0}{}\norm{\nabla f_t-\nabla f_0}{} + \norm{\nabla f_t-\nabla f_0}{}^2 \\
         	&\le 2\sigmax L(\X) \norm{\theta_t-\theta_0}{} + L(\X)^2 \norm{\theta_t-\theta_0}{}^2 \\
         	&\le \label{eqaux_montanariancora}2\sigmax L(\X) \norm{\theta_t-\theta_0}{} + L(\X) \norm{\theta_t-\theta_0}{}\frac{\sigmin}{2} \\
         	&\le \frac{5}{2}\sigmax L(\X) \norm{\theta_t-\theta_0}{},
         \end{align}
         where in \eqref{eqaux_montanariancora} we applied the definition of $t_*$.
         
         Moreover, by Grönwall and Cauchy-Schwarz inequalities we have
         \begin{align}
         	\norm{r_t}{} &\le \exp\left(-\frac{\lammin t}{4}\right)\norm{f_0-y}{}\int_0^t\norm{k_s-k_0}{}ds \\
            &\le\exp\left(-\frac{\lammin t}{4}\right)\norm{f_0-y}{}\sup_{s \ge 0}\norm{k_s-k_0}{} \\
            &\le\exp\left(-\frac{\lammin t}{4}\right)\frac{5}{2}\sigmax L(\X) \norm{f_0-y}{}\sup_{s \ge 0}\norm{\theta_s-\theta_0}{} \\
            &\le\exp\left(-\frac{\lammin t}{4}\right)\frac{5}{2}\sigmax L(\X) \norm{f_0-y}{}\sup_{s \ge 0}\frac{2}{\sigmin}\norm{f_0-y}{} \\
            &\label{eqaux_montanarib12}\le\exp\left(-\frac{\lammin t}{4}\right)\frac{5\sigmax}{\sigmin} L(\X) \norm{f_0-y}{}^2 \\
         \end{align}

    Moreover, by taking the difference of the gradient flow equations for $\theta_t$ and $\overline \theta_t$ we obtain:

\begin{align}
    \partv t \norm{\theta_t-\overline \theta_t}{} &\le \norm{\nabla f_t - \nabla f_0}{}\norm{f_t-y}{} + \norm{\nabla f_0}{}\norm{f_t-f_t^\lin}{} \\
    &\le L(\X) \norm{\theta_t-\theta_0}{}\norm{f_t-y}{} + \sigmax \norm{f_t-f_t^\lin}{}\\
    &\label{eqaux_montanaripostb12}\le \frac{2L(\X)}{\sigmin} \norm{f_0-y}{}^2 \exp\left(-\frac{\lammin t}{4}\right)\\
    &\quad + \frac{5\sigmax^2}{\sigmin} L(\X) \norm{f_0-y}{}^2 \exp\left(-\frac{\lammin t}{4}\right)\\
    &\le \frac{ (2+5\sigmax^2)L(\X)}{\sigmin} \norm{f_0-y}{}^2 \exp\left(-\frac{\lammin t}{4}\right).
\end{align}
where in \eqref{eqaux_montanaripostb12} we used \eqref{eqaux_montanariB3}, \eqref{equax_montanariB8} and \ref{eqaux_montanarib12}.

Integrating the previous inequality we obtain:
\begin{align}
    \norm{\theta_t-\overline \theta_t}{} &\le \frac{ (2+5\sigmax^2)L(\X)}{\sigmin} \norm{f_0-y}{}^2 \int_0^t \exp\left(-\frac{\lammin s}{4}\right)ds \\
    &\le \frac{ 4 (2+5\sigmax^2)L(\X)}{\sigmin^3} \norm{f_0-y}{}^2 \left(1- \exp\left(-\frac{\lammin s}{4}\right)\right)\\
    \label{eqaux_montanari_betabetalin}  &\le \frac{(8+20\sigmax^2)L(\X)}{\sigmin^3} \norm{f_0-y}{}^2.
\end{align}

    Now we are ready to prove the last inequality in the thesis. 
    Decompose by triangle inequality:
    \begin{align}
    \label{eqaux_montanari_final}
        \norm{y_t-\overline y_t}{} &\le \norm{y_t-f^\lin(x;\theta_t)}{}+\norm{f^\lin(x;\theta_t)-\overline y_t}{}.
    \end{align}

    First, let us focus on the first summand of \eqref{eqaux_montanari_final}.
    Denote by $L(x)$ the Lipschitz constant of $\nabla y_0$ seen as a function of $\theta$.
    Then, by Lemma \ref{lem:jacobian},
    \begin{align}
        \norm{y_t-f^\lin(x;\theta_t)}{} &= \norm{\int_0^t(\nabla f(x;\theta_s)-\nabla f(x;\theta_0))\frac{\partial}{\partial t} \theta_sds}{}\\
        &\le L(x) \sup_{t \ge 0}\norm{\theta_t-\theta_0}{}\int_0^t \norm{\frac{\partial}{\partial t} \theta_s}{}ds\\
        &\le L(x) \sup_{t \ge 0}\norm{\theta_t-\theta_0}{} \cdot \frac{2}{\sigmin} \norm{y-f_0}{}\\
        &\le L(x) \frac{4 \norm{y-f_0}{}^2}{\lammin},
    \end{align}
    where in the third inequality we used \eqref{eqaux_montanari_preB3} and \eqref{eqaux_montanariB3} on the last one.

    As for the second summand of \eqref{eqaux_montanari_final}, by \eqref{eqaux_montanari_betabetalin} and Lemma \ref{lem:jacobian}:
    \begin{align}
        \norm{f^\lin(x;\theta_t)-\overline y_t}{} &= \norm{\nabla f(x;\theta_0)(\theta_t-\overline \theta_t)}{}\\
        &\le \frac{(8+20\sigmax^2)L(\X)}{\sigmin^3} \norm{f_0-y}{}^2 \norm{\nabla f(x;\theta_0)}{}.
    \end{align}
 
Combining the two preceding estimations, we obtain the thesis.

\end{proof}

Lastly, we prove Theorem \ref{thm:stime_polinomiali}.
\begin{proof}[Proof of Theorem \ref{thm:stime_polinomiali}]
    We prove the three inequalities separately.
    Let $\lammin$ denote the smallest eigenvalue of $k_t$.
    \begin{itemize}
        \item By Lemma \ref{lemma1} and Cauchy-Schwarz's inequality,
        \begin{align}
            \label{eqaux_3triangle}
            \norm{y_t- f_t}{} &\le \frac{1}{\sqrt{n_1n_0}} \textup{Lip}\Phi \norm{x-\X}{} \norm{\theta_t^{(0)} \theta_t^{(1)} }{}.
        \end{align}
        Recall that $I_t(\lammin) \le t$.
        % and $I_t(I_t((\lammin))) \le 1$.
        Then the norm $\norm{\theta_t^{(0)} \theta_t^{(1)} }{}^2$ can be estimated with the aid of Lemma \ref{lem:gronwallbeta}:
        \begin{align}
            \norm{\theta_t^{(0)} \theta_t^{(1)}}{}^2 &= \sum_{u=1}^{n_0}\left(\sum_{v=1}^{n_1}(\theta^{(0)}_{uv})_t(\theta^{(1)}_v)_t\right)^2 \\
            &\le \sum_{u=1}^{n_0}\left(\sum_{v=1}^{n_1}(\theta_v^{(1)})_0(\theta_{uv}^{(0)})_0 + \frac{a_1(\theta_{uv}^{(0)})_0}{\sqrt{n_1}}\s(\theta_0)t + \frac{a_0(\theta_v^{(1)})_0}{n_1\sqrt{n_0}}\s(\theta_0)^2t^2 \right. \\
            &\quad \left. + \frac{a_0a_1}{n_1^\frac{3}{2}\sqrt{n_0}}\s(\theta_0)^3t^3 + \frac{a'_0(\theta_v^{(1)})^2_0}{\sqrt{n_1n_0}}\s(\theta_0)t + \frac{a'_0a_1(\theta^{(1)}_v)_0}{n_1\sqrt{n_0}}\s(\theta_0)^2t^2\right)^2\\
            &\le \sum_{u,v} n_1(\theta_v^{(1)})_0^2(\theta_{uv}^{(0)})_0^2 +a_1^2(\theta_{uv}^{(0)})^2_0\s(\theta_0)^2t^2 + \frac{a_0^2(\theta_v^{(1)})^2_0}{n_1n_0}\s(\theta_0)^4t^4 \\
            &\quad + \frac{a_0^2a_1^2}{n_1^2n_0}\s(\theta_0)^6t^6 + \frac{{a'}^2_0(\theta_v^{(1)})^4_0}{n_0}\s(\theta_0)^2t^2 + \frac{{a'}^2_0a_1^2(\theta^{(1)}_v)^2_0}{n_1n_0}\s(\theta_0)^4t^4 \\
            &\le n_1\norm{\theta^{(0)}_0\theta^{(1)}_0}{}^2 +a_1^2\norm{\theta^{(0)}_0}{}^2\s(\theta_0)^2t^2 + \frac{a_0^2\norm{\theta^{(1)}_0}{}^2}{n_1}\s(\theta_0)^4t^4 \\
            &\quad + \frac{a_0^2a_1^2}{n_1}\s(\theta_0)^6t^6 + {a'}^2_0\norm{\theta^{(1)}_0}{}^4\s(\theta_0)^2t^2 + \frac{{a'}^2_0a_1^2\norm{\theta^{(1)}_0}{}^2}{n_1}\s(\theta_0)^4t^4.
        \end{align}
        with $a_0 = \frac{1}{2}\norm{\Phi}{\infty}\norm{\Phi'}{\infty}\norm{\X_u}{}$, $a'_0 = \norm{\Phi'}{\infty}\norm{\X_u}{}$ and $a_1 = \norm{\Phi}{\infty}$.
        %with $\tilde a_1 = a_1^2$, 
        %$\tilde a_2=a_0^2$, 
        %$\tilde a_3=a_0^2a_1^2$,
        %$\tilde a_4={a'}^2_0$ and $\tilde a_5 = {a'}^2_0a_1^2$ positive constants not depending on $n_1$ nor $t$.

        Hence,
        \begin{align}
            \label{eqaux_1stsum}
            \norm{y_t- f_t}{}^2 &\le \frac{(\textup{Lip}\Phi)^2\norm{x-\X}{}}{n_0n_1}\norm{\theta^{(0)}_t\theta^{(1)}_t}{}^2\\
            &\le \frac{A_0}{n_0}\norm{\theta^{(0)}_0\theta^{(1)}_0}{}^2 +\frac{A_1t^2}{n_0n_1}\norm{\theta^{(0)}_0}{}^2\s(\theta_0)^2 + \frac{A_2\norm{\theta^{(1)}_0}{}^2t^4}{n_1^2n_0}\s(\theta_0)^4 \\
            &\quad + \frac{A_3t^6}{n_1^2n_0}\s(\theta_0)^6 + \frac{A_4t^2}{n_1n_0}\norm{\theta^{(1)}_0}{}^4\s(\theta_0)^2 + \frac{A_5t^4\norm{\theta^{(1)}_0}{}^2}{n_1^2n_0}\s(\theta_0)^4.
        \end{align}
        with $A_0 = (\textup{Lip}\Phi)^2\norm{x-\X}{}^2$,
        $A_1 = a_1^2$, $A_2= a_0^2$, $A_3 = a_1^2a_0^2$, $A_4 = {a'}_0^2$ and $A_5 = {a'}_0^2a_1^2$.

\item We follow a similar strategy to prove the second inequality in the Theorem. 
Put $\overline w_t = \overline \theta_t - \theta_0$.
By the triangle inequality and Cauchy-Schwarz we decompose:

\begin{align}
    \label{eqaux_flin}
        \norm{f^\lin-\overline y_t}{}^2 \le 2\norm{f_0 - y_0}{}^2 + 2\norm{\nabla_\theta f_0 - \nabla_\theta y_0}{}^2\norm{\overline w_t}{}^2.
    \end{align}
    The first summand in \eqref{eqaux_flin} is bounded exactly as the first summand in \eqref{eqaux_3triangle} by setting $t=0$:
    \begin{align}
        \norm{f_0 - y_0}{}^2 &\le  \frac{(\textup{Lip}\Phi)^2 \norm{x-\X}{}^2}{n_1n_0} \norm{\theta_0^{(0)}\theta_0^{(1)} }{}^2 \le \frac{A_0}{n_1n_0}\norm{\theta_0^{(0)}\theta_0^{(1)} }{}^2.
    \end{align}

    As for the second summand in \eqref{eqaux_flin}, we decompose by Lemma \ref{lemma1}.
    Factoring out $\max_i \{(\theta^{(1)}_0)_i\} = \norm{\theta^{(1)}_0}{\infty}$ permits us to write:
    \begin{align}
        \norm{\nabla_\theta f_0 - \nabla_\theta y_0}{}^2 &\le \norm{\partv{\theta^{(0)}}(f_0 - y_0)}{}^2+\norm{\partv{\theta^{(1)}}(f_0 - y_0)}{}^2 \\
        &\le \frac{1}{n_1n_0} \norm{(\X^\top\Phi'(\frac{1}{\sqrt{n_0}}\X\theta_0^{(0)})-x^\top\Phi'(\frac{1}{\sqrt{n_0}}x\theta_0^{(0)}))\theta_0^{(1)}}{}^2\\
        &\quad + \frac{1}{n_1}\norm{\Phi(\frac{1}{\sqrt{n_0}}\X\theta_0^{(0)})-\Phi(\frac{1}{\sqrt{n_0}}x\theta_0^{(0)})}{}^2 \\
        &\le \frac{2}{n_1n_0} \norm{(\X^\top\Phi'(\frac{1}{\sqrt{n_0}}\X\theta_0^{(0)})-\X^\top\Phi'(\frac{1}{\sqrt{n_0}}x\theta_0^{(0)}))\theta_0^{(1)}}{}^2\\
        &\quad +\frac{2}{n_1n_0} \norm{(\X^\top\Phi'(\frac{1}{\sqrt{n_0}}x\theta_0^{(0)})-x^\top\Phi'(\frac{1}{\sqrt{n_0}}x\theta_0^{(0)}))\theta_0^{(1)}}{}^2\\
        &\quad + \frac{A_0}{n_1n_0}\norm{\theta_0^{(0)}}{}^2 \\
        &\le \frac{2}{n_1n_0^2}(\textup{Lip}\Phi')^2\norm{\X}{}^2 \norm{x-\X}{}^2\norm{\theta_0^{(0)}\theta_0^{(1)}}{}^2\\
        &\quad +\frac{2}{n_1n_0} \norm{\Phi'}{\infty}^2\norm{x-\X}{}^2\norm{\theta_0^{(1)}}{}^2 + \frac{A_0}{n_1n_0}\norm{\theta_0^{(0)}}{}^2.
    \end{align}

    Moreover we can bound the norm of $\overline w_t$ with Lemma \ref{lem:gronwallbeta}:
    \begin{align}
        \label{eqaux_flinaux2}
        \norm{\overline w_t}{}^2 &\le \norm{\overline \theta_t^{(0)}-\theta_0^{(0)}}{}^2 + \norm{\overline \theta_t^{(1)}-\theta_0^{(1)}}{}^2 \\
        &\le \sum_{u,v}\frac{2a_0^2}{n_1^2n_0}\s(\theta_0)^4t^4 + \frac{2{a'}^2_0(\theta^{(1)}_v)^2_0}{n_1n_0}\s(\theta_0)^2t^2+ \sum_{v} \frac{a_1^2}{n_1}\s(\theta_0)^2t^2\\
        &\le \frac{2a_0^2}{n_1}\s(\theta_0)^4t^4 + \frac{2{a'}^2_0\norm{\theta^{(1)}}{}^2}{n_1}\s(\theta_0)^2t^2+ a_1^2\s(\theta_0)^2t^2.
    \end{align}

    Hence, \eqref{eqaux_flin} can be written as:
    \begin{align}
        \label{eqaux_3rdsum} \norm{f^\lin-\overline y_t}{}^2 &\le \frac{B_0}{n_1n_0}\norm{\theta_0^{(0)}\theta_0^{(1)} }{}^2 + \frac{B_1}{n_1^2n_0^2}\norm{\theta_0^{(0)}\theta_0^{(1)} }{}^2\s(\theta_0)^4t^4\\
        &\quad +\frac{B_2}{n_1^2n_0^2}\norm{\theta_0^{(0)}}{}^2\norm{\theta_0^{(1)} }{}^4\s(\theta_0)^2t^2+\frac{B_3}{n_1n_0^2}\norm{\theta_0^{(0)}\theta_0^{(1)} }{}^2\s(\theta_0)^2 t^2\\
        &\quad +\frac{B_4}{n_1^2n_0}\norm{\theta_0^{(1)} }{}^2\s(\theta_0)^4t^4+\frac{B_5}{n_1^2n_0}\norm{\theta_0^{(1)} }{}^4\s(\theta_0)^2t^2 \\
        &\quad +\frac{B_6}{n_1n_0}\norm{\theta_0^{(1)} }{}^2\s(\theta_0)^2t^2+\frac{B_7}{n_1^2n_0}\norm{\theta_0^{(0)} }{}^2\s(\theta_0)^4t^4 \\
        &\quad +\frac{B_8}{n_1^2n_0}\norm{\theta_0^{(0)} }{}^2\norm{\theta_0^{(1)} }{}^2\s(\theta_0)^2t^2+\frac{B_9}{n_1n_0}\norm{\theta_0^{(0)} }{}^2\s(\theta_0)^2t^2,
    \end{align}
    where the constants in the last inequality are, explicitly, $B_0=2A_0$,
    $B_1=8(\textup{Lip}\Phi'\norm{\X}{}\norm{x-\X}{}a_0)^2$,
    $B_2=8(\textup{Lip}\Phi'\norm{\X}{}\norm{x-\X}{}a'_0)^2$,
    $B_3=4(\textup{Lip}\Phi'\norm{\X}{}\norm{x-\X}{}a_1)^2$,
    $B_4=8(\norm{\Phi'}{\infty}\norm{x-\X}{}a_0)^2$,
    $B_5=8(\norm{\Phi'}{\infty}\norm{x-\X}{}a'_0)^2$,
    $B_6=4(\norm{\Phi'}{\infty}\norm{x-\X}{}a_1)^2$,
    $B_7=4A_0a_0^2$,
    $B_8=4A_0{a'}_0^2$, and $B_9 = 2A_0a_1^2$.

    \item It remains to estimate the last inequality.
        Consider $\Delta(t) = \norm{f_t - f_t^\lin}{}$
        Then by gradient flow equations and Cauchy-Schwarz,
        \begin{align}
            \partv{t}(\Delta(t)^2) & = \langle k_t(f_t-y) -k_0(f_t^\lin-y), f_t - f_t^\lin \rangle \\
            &= \sum_{i=1}^n (k_t(x_i,\X)(f_t-y) -k_0(x_i,\X)(f_t^\lin-y))(f_t(x_i) - f_t^\lin(x_i))\\
            &= \sum_{i=1}^n (k_t(x_i,\X)-k_0(x_i,\X))(f_t-y)(f_t(x_i) - f_t^\lin(x_i)) \\
            &\quad -k_0(x_i,\X)(f_t - f_t^\lin)(f_t(x_i) - f_t^\lin(x_i))\\
            &=\norm{(k_t-k_0)(f_t-y)(f_t-f_t^\lin)^\top}{1} - \norm{k_0(f_t-f_t^\lin)(f_t-f_t^\lin)^\top}{1}
        \end{align}
        By equivalence of the $1$-norm and the euclidean norm for $v\in \R^d$ we have $\norm{v}{} \le \norm{v}{1} \le \sqrt{d}\norm{v}{}$.
        Then, by Cauchy-Schwarz's inequality,
        \begin{align}
            \frac{\partial}{\partial t} \Delta(t) &= n\norm{k_t-k_0}{}\norm{f_t-y}{} -\lammin^0\Delta(t) \\
            &\label{eqaux_delta} \le n\norm{k_t-k_0}{}\s(\theta_0)e^{-\lammin t} -\lammin^0\Delta(t)
        \end{align}

        Let us bound the norm of $k_t-k_0$:
        \begin{align}
            \norm{k_t-k_0}{}&= \norm{\nabla_\theta f_t(\X)\nabla_\theta f_t(\X)^\top-\nabla_\theta f_0(\X)\nabla_\theta f_0(\X)^\top }{}\\
            &\le \norm{\nabla_\theta f_t(\X)+\nabla_\theta f_0(\X)}{}L(\X)\norm{\theta_t-\theta_0}{}\\
        \end{align}

        From Lemmas \ref{lemma1} and \ref{lem:gronwallbeta}, we have:
        \begin{align}
            \norm{\nabla_\theta f_t(\X)}{}^2 &=\norm{\nabla_{\theta^{(0)}}f_t(\X)}{}^2+\norm{\nabla_{\theta^{(1)}} f_t(\X)}{}^2\\
            &\le \frac{\norm{\X}{}^2\norm{\Phi'}{\infty}^2\norm{\theta^{(1)}_t}{}^2}{n_1n_0} + \frac{\norm{\Phi}{\infty}^2}{n_1}\\
            &\label{eqaux_delta1}\le \frac{\norm{\X}{}^2\norm{\Phi'}{\infty}^2}{n_1n_0}\left(\norm{\theta^{(1)}_0}{} + \frac{\norm{\Phi}{\infty}\s(\theta^0)}{\sqrt{n_1}}t\right)^2 + \frac{\norm{\Phi}{\infty}^2}{n_1}.
        \end{align}

        Analogously,
        \begin{align}
            \norm{\nabla_\theta f_0(\X)}{}^2 &=\norm{\nabla_{\theta^{(0)}}f_0(\X)}{}^2+\norm{\nabla_{\theta^{(1)}} f_0(\X)}{}^2\\
            &\label{eqaux_delta2}\le \frac{\norm{\X}{}^2\norm{\Phi'}{\infty}^2\norm{\theta^{(1)}_0}{}^2}{n_1n_0} + \frac{\norm{\Phi}{\infty}^2}{n_1}.
        \end{align}

        Moreover, again by Lemma \ref{lem:gronwallbeta},
        \begin{align}
            \norm{\theta_t-\theta_0}{}^2 &= \norm{\theta_t^{(0)}-\theta_0^{(0)}}{}^2 + \norm{\theta_t^{(1)}-\theta_0^{(1)}}{}^2 \\
            &\label{eqaux_delta3}\le \frac{2a_0^2}{n_1^2n_0}\s(\theta_0)^4t^4 + \frac{2{a'_0}^2}{n_1n_0}\norm{\theta^{(1)}_0}{}^2t^2 +\frac{a_1^2\s(\theta^0)^2}{n_1}t^2 \\
        \end{align}

        Inequalities \eqref{eqaux_delta1},\eqref{eqaux_delta2} and \eqref{eqaux_delta3} allow us to estimate:
        \begin{align}
            \label{eqaux_monster}
            \norm{k_t-k_0}{}&\le \frac{L(\X)}{n_1}\left(\frac{c_1 \s(\theta_0)^2\norm{\theta^{(1)}_0}{}}{\sqrt{n_1}n_0}+\frac{c_2 \s(\theta_0)^3t^3}{n_1n_0}+\frac{c_3 \s(\theta_0)^2t^2}{\sqrt{n_1n_0}} \right.\\
            &\left.\quad +\frac{c_4 \norm{\theta^{(1)}_0}{}^2t}{n_0}+\frac{c_5 \s(\theta_0)\norm{\theta^{(1)}_0}{}t^2}{\sqrt{n_1}n_0}+\frac{c_6 \norm{\theta^{(1)}_0}{}t}{\sqrt{n_0}}\right.\\
            &\left.\quad +\frac{c_7 \s(\theta_0)\norm{\theta^{(1)}_0}{}t}{\sqrt{n_0}}+\frac{c_8 \s(\theta_0)^2t^2}{\sqrt{n_1n_0}}+c_9\s(\theta_0)t\right),
        \end{align}
        with $c_1 = 2\sqrt{2}a_0\norm{\X}{}\norm{\Phi'}{\infty}$,
$c_2 = \sqrt{2}a_0\norm{\Phi}{\infty}$,
$c_3 = 2\sqrt{2}a_0\norm{\Phi}{\infty}$,
$c_4 = 2\sqrt{2}a'_0\norm{\X}{}\norm{\Phi'}{\infty}$,
$c_5 = \sqrt{2}a'_0\norm{\Phi}{\infty}$,
$c_6 = 2\sqrt{2}a'_0\norm{\Phi}{\infty}$,
$c_7 = 2a_1\norm{\X}{}\norm{\Phi'}{\infty}$,
$c_8 = a_1\norm{\Phi}{\infty}$ and $c_9 = 2a_1\norm{\Phi}{\infty}$.
Let $C(n_1,n_0,t,\theta_0)$ be the right hand side of \eqref{eqaux_monster}.
Then, the reight-hand side of \eqref{eqaux_delta} can be $\mathbb P$-almost surely bounded from above with:
\begin{align}
    \frac{\partial}{\partial t} \Delta(t) &\le n\norm{k_t-k_0}{}\s(\theta_0)e^{-\lammin t} -\lammin^0\Delta(t) \\
    &\le nC(n_1,n_0,t,\theta_0).
\end{align}
In the previous inequality we used that the event $\lammin=0$ has null measure.
Integrating, and using that $\Delta(0)=0$:
\begin{align}
    \Delta(t) &\le \frac{nL(\X)}{n_1}\left(\frac{c_1 \s(\theta_0)^2\norm{\theta^{(1)}_0}{}t}{\sqrt{n_1}n_0}+\frac{c_2 \s(\theta_0)^3t^4}{2n_1n_0}+\frac{c_3 \s(\theta_0)^2t^3}{\sqrt{n_1n_0}}\right. \\
    &\left.\quad +\frac{c_4 \norm{\theta^{(1)}_0}{}^2t^2}{2n_0}+\frac{c_5 \s(\theta_0)\norm{\theta^{(1)}_0}{}t^3}{3\sqrt{n_1}n_0}+\frac{c_6 \norm{\theta^{(1)}_0}{}t^3}{2\sqrt{n_0}}\right.\\
    &\left.\quad +\frac{c_7 \s(\theta_0)\norm{\theta^{(1)}_0}{}t^2}{2\sqrt{n_0}}+\frac{c_8 \s(\theta_0)^2t^3}{3\sqrt{n_1n_0}}+\frac{c_9\s(\theta_0)t^2}{2}\right)
\end{align}
Taking the square, applying the elementary inequality $(\sum^n_{i=1}a_i)^2\le n\sum^n_{i=1}a_i^2$, for $a_i\ge 0$,
and adjusting the constants yields the desired result.
    \end{itemize}
\end{proof}

\end{document}

